\clearpage\section*{索引}\phantomsection\label{app:index}
\addcontentsline{toc}{section}{索引}
按原书英文词条顺序保留主项与子项，附中文译名.下列页码由本译稿中相关概念、定理、公式或习题的位置生成；同一译稿页的重复定位合并.原词条及原书页码的完整映射保存在核验资料中.\footnote{译注：原索引的“$t$ 检验／一致最强／无偏”（原印刷331页）与“置信区间／一致最准确／正态”（原印刷372页）分类比所引内容的实际陈述更宽. 此处分别定位到\autoref{exercise:8.37}、\autoref{exercise:9.47}的相关讨论，保留原分类，不据索引额外断言其结论. “每次比较错误率”原指印刷458页，明确定义实际在下一页的\hyperref[item:11.17.b]{习题11.17(b)}；“回归／残差”的明确公式同样在原索引所指页的下一页. 页码已按实际语义位置生成.}
\begingroup\small\raggedright\interlinepenalty=10000\setlength{\parskip}{0pt}\setlength{\columnsep}{1.5em}
\begin{multicols}{2}
\sourcepageanchor{556}
\par\addvspace{2pt}
\par\noindent\hangindent=1em\hspace*{0em}\phantomsection\label{idx:1}接受／拒绝算法（Accept-Reject algorithm），\hyperref[ex:5.6.7]{\pageref*{ex:5.6.7}}, \hyperref[thm:5.6.8]{\pageref*{thm:5.6.8}}\par
\par\addvspace{2pt}
\par\noindent\hangindent=1em\hspace*{0em}\phantomsection\label{idx:2}验收抽样（Acceptance sampling），\hyperref[idxocc:2:333]{\pageref*{idxocc:2:333}}, \hyperref[ex:8.2.9]{\pageref*{ex:8.2.9}}\par
\par\noindent\hangindent=2em\hspace*{1em}\phantomsection\label{idx:3}超几何（hypergeometric），\hyperref[ex:3.2.1]{\pageref*{ex:3.2.1}}\par
\par\addvspace{2pt}
\par\noindent\hangindent=1em\hspace*{0em}\phantomsection\label{idx:4}行动空间（Action space），\hyperref[idxocc:4:355]{\pageref*{idxocc:4:355}}\par
\par\addvspace{2pt}
\needspace{3\baselineskip}
\par\noindent\hangindent=1em\hspace*{0em}\phantomsection\label{idx:5}算法（Algorithm）\par
\par\noindent\hangindent=2em\hspace*{1em}\phantomsection\label{idx:6}接受／拒绝（accept-reject），\hyperref[ex:5.6.7]{\pageref*{ex:5.6.7}}, \hyperref[thm:5.6.8]{\pageref*{thm:5.6.8}}\par
\par\noindent\hangindent=2em\hspace*{1em}\phantomsection\label{idx:7}EM，\hyperref[sec:7.2.4]{\pageref*{sec:7.2.4}}, \hyperref[sec:7.5.4]{\pageref*{sec:7.5.4}}\par
\par\noindent\hangindent=2.8em\hspace*{1.8em}\phantomsection\label{idx:8}完整数据（complete-data），\hyperref[idxocc:8:295]{\pageref*{idxocc:8:295}}\par
\par\noindent\hangindent=2.8em\hspace*{1.8em}\phantomsection\label{idx:9}完整数据似然（complete-data likelihood），\hyperref[eq:7.2.11]{\pageref*{eq:7.2.11}}\par
\par\noindent\hangindent=2.8em\hspace*{1.8em}\phantomsection\label{idx:10}收敛（convergence），\hyperref[thm:7.2.20]{\pageref*{thm:7.2.20}}, \hyperref[thm:7.5.2]{\pageref*{thm:7.5.2}}\par
\par\noindent\hangindent=2.8em\hspace*{1.8em}\phantomsection\label{idx:11}不完整数据（incomplete-data），\hyperref[idxocc:11:295]{\pageref*{idxocc:11:295}}\par
\par\noindent\hangindent=2.8em\hspace*{1.8em}\phantomsection\label{idx:12}不完整数据似然（incomplete-data likelihood），\hyperref[eq:7.2.13]{\pageref*{eq:7.2.13}}\par
\par\noindent\hangindent=2.8em\hspace*{1.8em}\phantomsection\label{idx:13}缺失数据（missing data），\hyperref[sec:7.2.4]{\pageref*{sec:7.2.4}}\par
\par\noindent\hangindent=2em\hspace*{1em}\phantomsection\label{idx:14}重要性抽样（importance sampling），\hyperref[exercise:5.64]{\pageref*{exercise:5.64}}\par
\par\noindent\hangindent=2em\hspace*{1em}\phantomsection\label{idx:15}Metropolis，\hyperref[idxocc:15:238]{\pageref*{idxocc:15:238}}\par
\par\addvspace{2pt}
\needspace{3\baselineskip}
\par\noindent\hangindent=1em\hspace*{0em}\phantomsection\label{idx:16}算法（Algorithms）\par
\par\noindent\hangindent=2em\hspace*{1em}\phantomsection\label{idx:17}马尔可夫链蒙特卡洛（Markov chain Monte Carlo），\hyperref[sec:5.8.5]{\pageref*{sec:5.8.5}}\par
\par\addvspace{2pt}
\par\noindent\hangindent=1em\hspace*{0em}\phantomsection\label{idx:18}辅助统计量（Ancillary statistic），\hyperref[def:6.2.16]{\pageref*{def:6.2.16}}, \hyperref[sec:6.6.2]{\pageref*{sec:6.6.2}}\par
\par\noindent\hangindent=2em\hspace*{1em}\phantomsection\label{idx:19}指数（exponential），\hyperref[ex:6.2.26]{\pageref*{ex:6.2.26}}\par
\par\noindent\hangindent=2em\hspace*{1em}\phantomsection\label{idx:20}位置族（location family），\hyperref[ex:6.2.18]{\pageref*{ex:6.2.18}}, \hyperref[exercise:6.11]{\pageref*{exercise:6.11}}\par
\par\noindent\hangindent=2em\hspace*{1em}\phantomsection\label{idx:21}位置尺度族（location–scale family），\hyperref[exercise:6.40]{\pageref*{exercise:6.40}}\par
\par\noindent\hangindent=2em\hspace*{1em}\phantomsection\label{idx:22}样本量（sample size），\hyperref[exercise:6.12]{\pageref*{exercise:6.12}}\par
\par\noindent\hangindent=2em\hspace*{1em}\phantomsection\label{idx:23}尺度族（scale family），\hyperref[ex:6.2.19]{\pageref*{ex:6.2.19}}\par
\par\noindent\hangindent=2em\hspace*{1em}\phantomsection\label{idx:24}均匀（uniform），\hyperref[ex:6.2.17]{\pageref*{ex:6.2.17}}\par
\par\noindent\hangindent=2em\hspace*{1em}\phantomsection\label{idx:25}用途（usefulness），\hyperref[ex:6.2.20]{\pageref*{ex:6.2.20}}\par
\par\addvspace{2pt}
\needspace{3\baselineskip}
\par\noindent\hangindent=1em\hspace*{0em}\phantomsection\label{idx:26}方差分析（ANOVA）\par
\par\noindent\hangindent=2em\hspace*{1em}\phantomsection\label{idx:27}假设（assumptions），\hyperref[sec:11.2.1]{\pageref*{sec:11.2.1}}\par
\par\noindent\hangindent=2em\hspace*{1em}\phantomsection\label{idx:28}Bonferroni不等式（Bonferroni Inequality），\hyperref[ex:11.2.9]{\pageref*{ex:11.2.9}}\par
\par\noindent\hangindent=2em\hspace*{1em}\phantomsection\label{idx:29}单元均值模型（cell means model），\hyperref[idxocc:29:454]{\pageref*{idxocc:29:454}}\par
\par\noindent\hangindent=2em\hspace*{1em}\phantomsection\label{idx:30}置信区间（confidence interval），\hyperref[eq:11.2.8]{\pageref*{eq:11.2.8}}\par
\par\noindent\hangindent=2em\hspace*{1em}\phantomsection\label{idx:31}对比平方和（contrast sum of squares），\hyperref[eq:11.2.18]{\pageref*{eq:11.2.18}}\par
\par\noindent\hangindent=2em\hspace*{1em}\phantomsection\label{idx:32}对比（contrasts），\hyperref[def:11.2.4]{\pageref*{def:11.2.4}}, \hyperref[ex:11.2.6]{\pageref*{ex:11.2.6}}, \hyperref[idxocc:32:459]{\pageref*{idxocc:32:459}}, \hyperref[sec:11.2.5]{\pageref*{sec:11.2.5}}\par
\par\noindent\hangindent=2em\hspace*{1em}\phantomsection\label{idx:33}数据窥探（data snooping），\hyperref[idxocc:33:464]{\pageref*{idxocc:33:464}}, \hyperref[exercise:11.18]{\pageref*{exercise:11.18}}\par
\par\noindent\hangindent=2em\hspace*{1em}\phantomsection\label{idx:34}方差相等（equality of variance），\hyperref[item:11.2.assumption.iii]{\pageref*{item:11.2.assumption.iii}}\par
\par\noindent\hangindent=2em\hspace*{1em}\phantomsection\label{idx:35}均方的期望（expected mean squares），\hyperref[exercise:11.22]{\pageref*{exercise:11.22}}\par
\par\noindent\hangindent=2em\hspace*{1em}\phantomsection\label{idx:36}F检验（F test），\hyperref[eq:11.2.14]{\pageref*{eq:11.2.14}}\par
\par\noindent\hangindent=2.8em\hspace*{1.8em}\phantomsection\label{idx:37}与t检验的关系（relation to t test），\hyperref[exercise:11.12]{\pageref*{exercise:11.12}}\par
\par\noindent\hangindent=2em\hspace*{1em}\phantomsection\label{idx:38}可识别性（identifiability），\hyperref[def:11.2.2]{\pageref*{def:11.2.2}}\par
\par\noindent\hangindent=2em\hspace*{1em}\phantomsection\label{idx:39}因子的水平（levels of a factor），\hyperref[idxocc:39:453]{\pageref*{idxocc:39:453}}\par
\par\noindent\hangindent=2em\hspace*{1em}\phantomsection\label{idx:40}线性组合（linear combinations），\hyperref[sec:11.2.3]{\pageref*{sec:11.2.3}}\par
\par\noindent\hangindent=2em\hspace*{1em}\phantomsection\label{idx:41}似然比检验（LRT），\hyperref[exercise:11.13]{\pageref*{exercise:11.13}}\par
\par\noindent\hangindent=2em\hspace*{1em}\phantomsection\label{idx:42}零假设（null hypothesis），\hyperref[eq:11.2.3]{\pageref*{eq:11.2.3}}\par
\par\noindent\hangindent=2em\hspace*{1em}\phantomsection\label{idx:43}单因素（oneway），\hyperref[idxocc:43:452]{\pageref*{idxocc:43:452}}\par
\par\noindent\hangindent=2em\hspace*{1em}\phantomsection\label{idx:44}正交对比（orthogonal contrast），\hyperref[idxocc:44:465]{\pageref*{idxocc:44:465}}\par
\par\noindent\hangindent=2em\hspace*{1em}\phantomsection\label{idx:45}其他模型（other models），\hyperref[sec:11.5.4]{\pageref*{sec:11.5.4}}\par
\par\noindent\hangindent=2em\hspace*{1em}\phantomsection\label{idx:46}过参数化模型（overparameterized model），\hyperref[eq:11.2.2]{\pageref*{eq:11.2.2}}\par
\par\noindent\hangindent=2em\hspace*{1em}\phantomsection\label{idx:47}Scheffé方法（Scheffé’s method），\hyperref[thm:11.2.10]{\pageref*{thm:11.2.10}}, \hyperref[exercise:11.14]{\pageref*{exercise:11.14}}\par
\par\noindent\hangindent=2em\hspace*{1em}\phantomsection\label{idx:48}Stein估计（Stein estimation），\hyperref[sec:11.5.6]{\pageref*{sec:11.5.6}}\par
\par\noindent\hangindent=2em\hspace*{1em}\phantomsection\label{idx:49}充分统计量（sufficient statistic），\hyperref[item:11.6.a]{\pageref*{item:11.6.a}}\par
\par\noindent\hangindent=2em\hspace*{1em}\phantomsection\label{idx:50}平方和（sum of squares），\hyperref[thm:11.2.11]{\pageref*{thm:11.2.11}}\par
\par\noindent\hangindent=2em\hspace*{1em}\phantomsection\label{idx:51}t检验（t test），\hyperref[eq:11.2.7]{\pageref*{eq:11.2.7}}\par
\par\noindent\hangindent=2.8em\hspace*{1.8em}\phantomsection\label{idx:52}与F检验的关系（relation to F test），\hyperref[exercise:11.12]{\pageref*{exercise:11.12}}\par
\par\noindent\hangindent=2em\hspace*{1em}\phantomsection\label{idx:53}表（table），\hyperref[ex:11.2.12]{\pageref*{ex:11.2.12}}\par
\par\noindent\hangindent=2em\hspace*{1em}\phantomsection\label{idx:54}数据变换（transformation of data），\hyperref[idxocc:54:455]{\pageref*{idxocc:54:455}}, \hyperref[exercise:11.1]{\pageref*{exercise:11.1}}\par
\par\noindent\hangindent=2em\hspace*{1em}\phantomsection\label{idx:55}处理均值（treatment means），\hyperref[idxocc:55:453]{\pageref*{idxocc:55:453}}\par
\par\noindent\hangindent=2em\hspace*{1em}\phantomsection\label{idx:56}并交检验（union–intersection test），\hyperref[thm:11.2.5]{\pageref*{thm:11.2.5}}, \hyperref[sec:11.2.4]{\pageref*{sec:11.2.4}}\par
\par\noindent\hangindent=2em\hspace*{1em}\phantomsection\label{idx:57}方差估计量（variance estimator），\hyperref[eq:11.2.5]{\pageref*{eq:11.2.5}}\par
\par\addvspace{2pt}
\needspace{3\baselineskip}
\par\noindent\hangindent=1em\hspace*{0em}\phantomsection\label{idx:58}近似（Approximation）\par
\par\noindent\hangindent=2em\hspace*{1em}\phantomsection\label{idx:59}以正态近似二项（binomial by normal），\hyperref[ex:3.3.2]{\pageref*{ex:3.3.2}}\par
\par\noindent\hangindent=2em\hspace*{1em}\phantomsection\label{idx:60}以泊松近似二项（binomial by Poisson），\hyperref[ex:2.3.13]{\pageref*{ex:2.3.13}}, \hyperref[eq:3.2.8]{\pageref*{eq:3.2.8}}\par
\par\noindent\hangindent=2em\hspace*{1em}\phantomsection\label{idx:61}以二项近似超几何（hypergeometric by binomial），\hyperref[idxocc:61:137]{\pageref*{idxocc:61:137}}\par
\par\noindent\hangindent=2em\hspace*{1em}\phantomsection\label{idx:62}以泊松近似超几何（hypergeometric by Poisson），\hyperref[itm:exercise:3.11:b]{\pageref*{itm:exercise:3.11:b}}\par
\par\noindent\hangindent=2em\hspace*{1em}\phantomsection\label{idx:63}以正态近似负二项（negative binomial by normal），\hyperref[ex:5.5.16]{\pageref*{ex:5.5.16}}\par
\par\addvspace{2pt}
\par\noindent\hangindent=1em\hspace*{0em}\phantomsection\label{idx:64}算术平均数（Arithmetic mean），\hyperref[ex:4.7.8]{\pageref*{ex:4.7.8}}\par
\par\addvspace{2pt}
\par\noindent\hangindent=1em\hspace*{0em}\phantomsection\label{idx:65}渐近效率（Asymptotic efficiency），\hyperref[def:10.1.11]{\pageref*{def:10.1.11}}\par
\par\addvspace{2pt}
\par\noindent\hangindent=1em\hspace*{0em}\phantomsection\label{idx:66}渐近相对效率（Asymptotic Relative Efficiency），\hyperref[def:10.1.16]{\pageref*{def:10.1.16}}\par
\par\noindent\hangindent=2em\hspace*{1em}\phantomsection\label{idx:67}Huber估计量（Huber estimator），\hyperref[ex:10.2.7]{\pageref*{ex:10.2.7}}\par
\par\noindent\hangindent=2em\hspace*{1em}\phantomsection\label{idx:68}中位数相对于均值（median to mean），\hyperref[ex:10.2.4]{\pageref*{ex:10.2.4}}\par
\par\addvspace{2pt}
\needspace{3\baselineskip}
\par\noindent\hangindent=1em\hspace*{0em}\phantomsection\label{idx:69}渐近方差（Asymptotic variance）\par
\par\noindent\hangindent=2em\hspace*{1em}\phantomsection\label{idx:70}与影响函数的关系（relation to influence function），\hyperref[idxocc:70:450]{\pageref*{idxocc:70:450}}\par
\par\addvspace{2pt}
\par\noindent\hangindent=1em\hspace*{0em}\phantomsection\label{idx:71}渐近方差（另见\hyperref[idx:1062]{方差}）（Asymptotic variance (see also Variance)），\hyperref[idxocc:71:410]{\pageref*{idxocc:71:410}}\par
\par\addvspace{2pt}
\par\noindent\hangindent=1em\hspace*{0em}\phantomsection\label{idx:72}自回归模型（Autoregressive model），\hyperref[itm:exercise:10.19:c]{\pageref*{itm:exercise:10.19:c}}\par
\par\addvspace{2pt}
\par\noindent\hangindent=1em\hspace*{0em}\phantomsection\label{idx:73}连续性公理（Axiom of Continuity），\hyperref[exercise:1.12]{\pageref*{exercise:1.12}}\par
\par\addvspace{2pt}
\par\noindent\hangindent=1em\hspace*{0em}\phantomsection\label{idx:74}可数可加性公理（Axiom of Countable Additivity），\hyperref[idxocc:74:38]{\pageref*{idxocc:74:38}}, \hyperref[exercise:1.12]{\pageref*{exercise:1.12}}\par
\par\addvspace{2pt}
\par\noindent\hangindent=1em\hspace*{0em}\phantomsection\label{idx:75}有限可加性公理（Axiom of Finite Additivity），\hyperref[idxocc:75:39]{\pageref*{idxocc:75:39}}, \hyperref[exercise:1.12]{\pageref*{exercise:1.12}}\par
\par\addvspace{2pt}
\par\noindent\hangindent=1em\hspace*{0em}\phantomsection\label{idx:76}Basu定理（Basu’s Theorem），\hyperref[thm:6.2.24]{\pageref*{thm:6.2.24}}\par
\par\noindent\hangindent=2em\hspace*{1em}\phantomsection\label{idx:77}逆命题（converse），\hyperref[sec:6.6.1]{\pageref*{sec:6.6.1}}\par
\par\addvspace{2pt}
\needspace{3\baselineskip}
\par\noindent\hangindent=1em\hspace*{0em}\phantomsection\label{idx:78}贝叶斯估计量（Bayes estimator）\par
\par\noindent\hangindent=2em\hspace*{1em}\phantomsection\label{idx:79}伯努利（Bernoulli），\hyperref[ex:7.2.14]{\pageref*{ex:7.2.14}}, \hyperref[ex:7.3.5]{\pageref*{ex:7.3.5}}\par
\par\noindent\hangindent=2em\hspace*{1em}\phantomsection\label{idx:80}正态（normal），\hyperref[ex:7.2.16]{\pageref*{ex:7.2.16}}, \hyperref[exercise:7.22]{\pageref*{exercise:7.22}}\par
\par\noindent\hangindent=2em\hspace*{1em}\phantomsection\label{idx:81}正态方差（normal variance），\hyperref[exercise:7.23]{\pageref*{exercise:7.23}}\par
\par\noindent\hangindent=2em\hspace*{1em}\phantomsection\label{idx:82}与无偏估计量的关系（relation to unbiased estimator），\hyperref[sec:7.5.2]{\pageref*{sec:7.5.2}}\par
\par\addvspace{2pt}
\par\noindent\hangindent=1em\hspace*{0em}\phantomsection\label{idx:83}贝叶斯风险（Bayes risk），\hyperref[idxocc:83:316]{\pageref*{idxocc:83:316}}\par
\par\addvspace{2pt}
\needspace{3\baselineskip}
\par\noindent\hangindent=1em\hspace*{0em}\phantomsection\label{idx:84}贝叶斯规则（Bayes rule）\par
\par\noindent\hangindent=2em\hspace*{1em}\phantomsection\label{idx:85}先验分布（Prior distribution），\hyperref[idxocc:85:316]{\pageref*{idxocc:85:316}}\par
\par\addvspace{2pt}
\par\noindent\hangindent=1em\hspace*{0em}\phantomsection\label{idx:86}贝叶斯公式（Bayes’ Rule），\hyperref[thm:1.3.5]{\pageref*{thm:1.3.5}}\par
\par\addvspace{2pt}
\par\noindent\hangindent=1em\hspace*{0em}\phantomsection\label{idx:87}Behrens–Fisher问题（Behrens–Fisher problem），\hyperref[exercise:8.42]{\pageref*{exercise:8.42}}, \hyperref[idxocc:87:455]{\pageref*{idxocc:87:455}}, \hyperref[idxocc:87:500]{\pageref*{idxocc:87:500}}\par
\sourcepageanchor{557}
\par\addvspace{2pt}
\par\noindent\hangindent=1em\hspace*{0em}\phantomsection\label{idx:88}伯努利试验（见\hyperref[idx:252]{分布：伯努利}）（Bernoulli trials (see Distribution Bernoulli)），\hyperref[idxocc:88:104]{\pageref*{idxocc:88:104}}\par
\par\addvspace{2pt}
\par\noindent\hangindent=1em\hspace*{0em}\phantomsection\label{idx:89}贝塔函数（另见\hyperref[idx:254]{分布：贝塔}）（Beta function (see also Distribution beta)），\hyperref[eq:3.3.17]{\pageref*{eq:3.3.17}}\par
\par\addvspace{2pt}
\par\noindent\hangindent=1em\hspace*{0em}\phantomsection\label{idx:90}偏差（Bias），\hyperref[def:7.3.2]{\pageref*{def:7.3.2}}\par
\par\addvspace{2pt}
\par\noindent\hangindent=1em\hspace*{0em}\phantomsection\label{idx:91}二项式系数（Binomial coefficients），\hyperref[def:1.2.17]{\pageref*{def:1.2.17}}\par
\par\addvspace{2pt}
\par\noindent\hangindent=1em\hspace*{0em}\phantomsection\label{idx:92}二项式定理（Binomial Theorem），\hyperref[thm:3.2.2]{\pageref*{thm:3.2.2}}\par
\par\addvspace{2pt}
\par\noindent\hangindent=1em\hspace*{0em}\phantomsection\label{idx:93}Birnbaum定理（Birnbaum’s Theorem），\hyperref[thm:6.3.6]{\pageref*{thm:6.3.6}}, \hyperref[exercise:6.33]{\pageref*{exercise:6.33}}\par
\par\addvspace{2pt}
\par\noindent\hangindent=1em\hspace*{0em}\phantomsection\label{idx:94}自助法（Bootstrap），\hyperref[sec:10.1.4]{\pageref*{sec:10.1.4}}\par
\par\noindent\hangindent=2em\hspace*{1em}\phantomsection\label{idx:95}非参数（nonparametric），\hyperref[ex:10.1.19]{\pageref*{ex:10.1.19}}\par
\par\noindent\hangindent=2em\hspace*{1em}\phantomsection\label{idx:96}参数（parametric），\hyperref[ex:5.6.6]{\pageref*{ex:5.6.6}}, \hyperref[idxocc:96:418]{\pageref*{idxocc:96:418}}\par
\par\addvspace{2pt}
\par\noindent\hangindent=1em\hspace*{0em}\phantomsection\label{idx:97}博雷尔域（见\hyperref[idx:972]{\ensuremath{\sigma}代数}）（Borel field (see Sigma algebra)），\hyperref[def:1.2.1]{\pageref*{def:1.2.1}}\par
\par\addvspace{2pt}
\par\noindent\hangindent=1em\hspace*{0em}\phantomsection\label{idx:98}博雷尔悖论（Borel Paradox），\hyperref[sec:4.9.3]{\pageref*{sec:4.9.3}}\par
\par\addvspace{2pt}
\par\noindent\hangindent=1em\hspace*{0em}\phantomsection\label{idx:99}Box–Muller变换（Box-Muller transformation），\hyperref[ex:5.6.4]{\pageref*{ex:5.6.4}}\par
\par\addvspace{2pt}
\par\noindent\hangindent=1em\hspace*{0em}\phantomsection\label{idx:100}崩溃值（Breakdown value），\hyperref[def:10.2.2]{\pageref*{def:10.2.2}}, \hyperref[exercise:10.20]{\pageref*{exercise:10.20}}\par
\par\addvspace{2pt}
\par\noindent\hangindent=1em\hspace*{0em}\phantomsection\label{idx:101}Carleman条件（Carleman’s Condition），\hyperref[idxocc:101:99]{\pageref*{idxocc:101:99}}\par
\par\addvspace{2pt}
\par\noindent\hangindent=1em\hspace*{0em}\phantomsection\label{idx:102}中心极限定理（Central Limit Theorem (CLT)），\hyperref[thm:5.5.14]{\pageref*{thm:5.5.14}}, \hyperref[thm:5.5.15]{\pageref*{thm:5.5.15}}, \hyperref[sec:5.8.1]{\pageref*{sec:5.8.1}}\par
\par\addvspace{2pt}
\par\noindent\hangindent=1em\hspace*{0em}\phantomsection\label{idx:103}特征函数（Characteristic function），\hyperref[idxocc:103:100]{\pageref*{idxocc:103:100}}\par
\par\noindent\hangindent=2em\hspace*{1em}\phantomsection\label{idx:104}收敛（convergence），\hyperref[thm:2.6.1]{\pageref*{thm:2.6.1}}\par
\par\noindent\hangindent=2em\hspace*{1em}\phantomsection\label{idx:105}在中心极限定理中的用途（use in CLT），\hyperref[thm:5.5.15]{\pageref*{thm:5.5.15}}\par
\par\addvspace{2pt}
\par\noindent\hangindent=1em\hspace*{0em}\phantomsection\label{idx:106}独立性卡方检验（Chi squared test of independence），\hyperref[exercise:10.31]{\pageref*{exercise:10.31}}\par
\par\addvspace{2pt}
\par\noindent\hangindent=1em\hspace*{0em}\phantomsection\label{idx:107}Cochran定理（Cochran’s Theorem），\hyperref[idxocc:107:465]{\pageref*{idxocc:107:465}}, \hyperref[sec:11.5.1]{\pageref*{sec:11.5.1}}\par
\par\addvspace{2pt}
\par\noindent\hangindent=1em\hspace*{0em}\phantomsection\label{idx:108}决定系数（Coefficient of determination），\hyperref[idxocc:108:481]{\pageref*{idxocc:108:481}}\par
\par\addvspace{2pt}
\par\noindent\hangindent=1em\hspace*{0em}\phantomsection\label{idx:109}巧合（Coincidences），\hyperref[exercise:3.9]{\pageref*{exercise:3.9}}\par
\par\addvspace{2pt}
\par\noindent\hangindent=1em\hspace*{0em}\phantomsection\label{idx:110}完备性（Completeness），\hyperref[def:6.2.21]{\pageref*{def:6.2.21}}\par
\par\noindent\hangindent=2em\hspace*{1em}\phantomsection\label{idx:111}辅助性（ancillarity），\hyperref[thm:6.2.24]{\pageref*{thm:6.2.24}}, \hyperref[thm:6.6.2]{\pageref*{thm:6.6.2}}\par
\par\noindent\hangindent=2em\hspace*{1em}\phantomsection\label{idx:112}Basu定理（Basu’s Theorem），\hyperref[thm:6.2.24]{\pageref*{thm:6.2.24}}\par
\par\noindent\hangindent=2em\hspace*{1em}\phantomsection\label{idx:113}最优无偏估计量（best unbiased estimator），\hyperref[thm:7.3.23]{\pageref*{thm:7.3.23}}\par
\par\noindent\hangindent=2em\hspace*{1em}\phantomsection\label{idx:114}二项（binomial），\hyperref[ex:6.2.22]{\pageref*{ex:6.2.22}}\par
\par\noindent\hangindent=2em\hspace*{1em}\phantomsection\label{idx:115}指数（exponential），\hyperref[ex:6.2.26]{\pageref*{ex:6.2.26}}, \hyperref[exercise:6.30]{\pageref*{exercise:6.30}}\par
\par\noindent\hangindent=2em\hspace*{1em}\phantomsection\label{idx:116}指数族（exponential family），\hyperref[thm:6.2.25]{\pageref*{thm:6.2.25}}\par
\par\noindent\hangindent=2em\hspace*{1em}\phantomsection\label{idx:117}独立性（independence），\hyperref[thm:6.2.24]{\pageref*{thm:6.2.24}}\par
\par\noindent\hangindent=2em\hspace*{1em}\phantomsection\label{idx:118}泊松（Poisson），\hyperref[exercise:6.18]{\pageref*{exercise:6.18}}\par
\par\noindent\hangindent=2em\hspace*{1em}\phantomsection\label{idx:119}充分性（sufficiency），\hyperref[thm:6.2.24]{\pageref*{thm:6.2.24}}\par
\par\noindent\hangindent=2em\hspace*{1em}\phantomsection\label{idx:120}均匀（uniform），\hyperref[ex:6.2.23]{\pageref*{ex:6.2.23}}\par
\par\addvspace{2pt}
\par\noindent\hangindent=1em\hspace*{0em}\phantomsection\label{idx:121}凹函数（Concave function），\hyperref[def:4.7.6]{\pageref*{def:4.7.6}}\par
\par\addvspace{2pt}
\par\noindent\hangindent=1em\hspace*{0em}\phantomsection\label{idx:122}条件概率（Conditional probability），\hyperref[def:1.3.2]{\pageref*{def:1.3.2}}\par
\par\noindent\hangindent=2em\hspace*{1em}\phantomsection\label{idx:123}乘法法则（multiplication rule），\hyperref[eq:1.3.3]{\pageref*{eq:1.3.3}}\par
\par\noindent\hangindent=2em\hspace*{1em}\phantomsection\label{idx:124}性质（properties），\hyperref[exercise:1.38]{\pageref*{exercise:1.38}}\par
\par\addvspace{2pt}
\par\noindent\hangindent=1em\hspace*{0em}\phantomsection\label{idx:125}条件性原则（Conditionality Principle），\hyperref[idxocc:125:269]{\pageref*{idxocc:125:269}}\par
\par\noindent\hangindent=2em\hspace*{1em}\phantomsection\label{idx:126}解释（interpretations），\hyperref[ex:6.3.7]{\pageref*{ex:6.3.7}}\par
\par\addvspace{2pt}
\par\noindent\hangindent=1em\hspace*{0em}\phantomsection\label{idx:127}置信系数（Confidence coefficient），\hyperref[def:9.1.5]{\pageref*{def:9.1.5}}, \hyperref[idxocc:127:387]{\pageref*{idxocc:127:387}}\par
\par\noindent\hangindent=2em\hspace*{1em}\phantomsection\label{idx:128}与风险的关系（relation to risk），\hyperref[idxocc:128:395]{\pageref*{idxocc:128:395}}\par
\par\addvspace{2pt}
\needspace{3\baselineskip}
\par\noindent\hangindent=1em\hspace*{0em}\phantomsection\label{idx:129}置信区间（Confidence interval）\par
\needspace{3\baselineskip}
\par\noindent\hangindent=2em\hspace*{1em}\phantomsection\label{idx:130}长度（length）\par
\par\noindent\hangindent=2.8em\hspace*{1.8em}\phantomsection\label{idx:131}正态（normal），\hyperref[ex:9.3.3]{\pageref*{ex:9.3.3}}\par
\par\noindent\hangindent=2.8em\hspace*{1.8em}\phantomsection\label{idx:132}最优（optimal），\hyperref[thm:9.3.2]{\pageref*{thm:9.3.2}}\par
\par\noindent\hangindent=2.8em\hspace*{1.8em}\phantomsection\label{idx:133}与错误覆盖的关系（relation to false coverage），\hyperref[thm:9.3.9]{\pageref*{thm:9.3.9}}\par
\par\noindent\hangindent=2.8em\hspace*{1.8em}\phantomsection\label{idx:134}与最高后验密度的关系（relation to HPD），\hyperref[idxocc:134:389]{\pageref*{idxocc:134:389}}\par
\par\noindent\hangindent=2.8em\hspace*{1.8em}\phantomsection\label{idx:135}单峰（unimodal），\hyperref[thm:9.3.2]{\pageref*{thm:9.3.2}}\par
\par\noindent\hangindent=2em\hspace*{1em}\phantomsection\label{idx:136}方差分析（ANOVA），\hyperref[eq:11.2.8]{\pageref*{eq:11.2.8}}\par
\par\noindent\hangindent=2em\hspace*{1em}\phantomsection\label{idx:137}近似（approximate），\hyperref[ex:10.4.4]{\pageref*{ex:10.4.4}}\par
\par\noindent\hangindent=2.8em\hspace*{1.8em}\phantomsection\label{idx:138}二项（binomial），\hyperref[ex:10.4.1]{\pageref*{ex:10.4.1}}\par
\par\noindent\hangindent=2.8em\hspace*{1.8em}\phantomsection\label{idx:139}最大似然（maximum likelihood），\hyperref[eq:10.4.2]{\pageref*{eq:10.4.2}}\par
\par\noindent\hangindent=2.8em\hspace*{1.8em}\phantomsection\label{idx:140}负二项（negative binomial），\hyperref[ex:10.4.9]{\pageref*{ex:10.4.9}}\par
\par\noindent\hangindent=2.8em\hspace*{1.8em}\phantomsection\label{idx:141}泊松（Poisson），\hyperref[ex:10.4.5]{\pageref*{ex:10.4.5}}\par
\needspace{3\baselineskip}
\par\noindent\hangindent=2em\hspace*{1em}\phantomsection\label{idx:142}二项（binomial）\par
\par\noindent\hangindent=2.8em\hspace*{1.8em}\phantomsection\label{idx:143}与F分布的关系（relation to F distribution），\hyperref[exercise:9.21]{\pageref*{exercise:9.21}}\par
\par\noindent\hangindent=2em\hspace*{1em}\phantomsection\label{idx:144}生物等效性（bioequivalence），\hyperref[exercise:9.33]{\pageref*{exercise:9.33}}\par
\par\noindent\hangindent=2em\hspace*{1em}\phantomsection\label{idx:145}构造（construction），\hyperref[ex:9.2.11]{\pageref*{ex:9.2.11}}, \hyperref[idxocc:145:381]{\pageref*{idxocc:145:381}}, \hyperref[exercise:9.17]{\pageref*{exercise:9.17}}\par
\par\noindent\hangindent=2em\hspace*{1em}\phantomsection\label{idx:146}决策论（decision theoretic），\hyperref[eq:9.3.4]{\pageref*{eq:9.3.4}}\par
\par\noindent\hangindent=2.8em\hspace*{1.8em}\phantomsection\label{idx:147}困难（difficulty），\hyperref[idxocc:147:396]{\pageref*{idxocc:147:396}}\par
\par\noindent\hangindent=2.8em\hspace*{1.8em}\phantomsection\label{idx:148}均值（mean），\hyperref[ex:9.3.13]{\pageref*{ex:9.3.13}}\par
\needspace{3\baselineskip}
\par\noindent\hangindent=2em\hspace*{1em}\phantomsection\label{idx:149}期望长度（expected length）\par
\par\noindent\hangindent=2.8em\hspace*{1.8em}\phantomsection\label{idx:150}正态（normal），\hyperref[ex:9.3.3]{\pageref*{ex:9.3.3}}\par
\par\noindent\hangindent=2em\hspace*{1em}\phantomsection\label{idx:151}指数（exponential），\hyperref[ex:9.2.13]{\pageref*{ex:9.2.13}}\par
\par\noindent\hangindent=2.8em\hspace*{1.8em}\phantomsection\label{idx:152}枢轴（pivotal），\hyperref[ex:9.2.9]{\pageref*{ex:9.2.9}}\par
\par\noindent\hangindent=2em\hspace*{1em}\phantomsection\label{idx:153}Fieller，\hyperref[sec:9.5.3]{\pageref*{sec:9.5.3}}\par
\needspace{3\baselineskip}
\par\noindent\hangindent=2em\hspace*{1em}\phantomsection\label{idx:154}长度（length）\par
\par\noindent\hangindent=2.8em\hspace*{1.8em}\phantomsection\label{idx:155}伽马（gamma），\hyperref[ex:9.3.4]{\pageref*{ex:9.3.4}}\par
\par\noindent\hangindent=2em\hspace*{1em}\phantomsection\label{idx:156}似然（likelihood），\hyperref[eq:9.2.7]{\pageref*{eq:9.2.7}}\par
\par\noindent\hangindent=2em\hspace*{1em}\phantomsection\label{idx:157}位置（location），\hyperref[exercise:9.51]{\pageref*{exercise:9.51}}\par
\needspace{3\baselineskip}
\par\noindent\hangindent=2em\hspace*{1em}\phantomsection\label{idx:158}似然比检验（LRT）\par
\par\noindent\hangindent=2.8em\hspace*{1.8em}\phantomsection\label{idx:159}二项（binomial），\hyperref[exercise:9.6]{\pageref*{exercise:9.6}}\par
\par\noindent\hangindent=2.8em\hspace*{1.8em}\phantomsection\label{idx:160}指数（exponential），\hyperref[ex:9.2.3]{\pageref*{ex:9.2.3}}\par
\needspace{3\baselineskip}
\par\noindent\hangindent=2em\hspace*{1em}\phantomsection\label{idx:161}负二项（negative binomial）\par
\par\noindent\hangindent=2.8em\hspace*{1.8em}\phantomsection\label{idx:162}与F分布的关系（relation to F distribution），\hyperref[exercise:9.22]{\pageref*{exercise:9.22}}\par
\par\noindent\hangindent=2em\hspace*{1em}\phantomsection\label{idx:163}正态（normal），\hyperref[ex:9.1.2]{\pageref*{ex:9.1.2}}, \hyperref[ex:9.2.1]{\pageref*{ex:9.2.1}}\par
\par\noindent\hangindent=2.8em\hspace*{1.8em}\phantomsection\label{idx:164}枢轴（pivotal），\hyperref[ex:9.2.10]{\pageref*{ex:9.2.10}}\par
\par\noindent\hangindent=2em\hspace*{1em}\phantomsection\label{idx:165}单侧（one-sided），\hyperref[idxocc:165:368]{\pageref*{idxocc:165:368}}\par
\par\noindent\hangindent=2.8em\hspace*{1.8em}\phantomsection\label{idx:166}二项（binomial），\hyperref[ex:9.2.5]{\pageref*{ex:9.2.5}}\par
\par\noindent\hangindent=2.8em\hspace*{1.8em}\phantomsection\label{idx:167}正态（normal），\hyperref[ex:9.2.4]{\pageref*{ex:9.2.4}}\par
\par\noindent\hangindent=2em\hspace*{1em}\phantomsection\label{idx:168}泊松（Poisson），\hyperref[ex:9.2.15]{\pageref*{ex:9.2.15}}\par
\par\noindent\hangindent=2em\hspace*{1em}\phantomsection\label{idx:169}均值之比（ratio of means），\hyperref[sec:9.5.3]{\pageref*{sec:9.5.3}}\par
\par\noindent\hangindent=2em\hspace*{1em}\phantomsection\label{idx:170}回归（regression），\hyperref[eq:11.3.37]{\pageref*{eq:11.3.37}}, \hyperref[sec:11.3.6]{\pageref*{sec:11.3.6}}\par
\par\noindent\hangindent=2em\hspace*{1em}\phantomsection\label{idx:171}与集合估计量的关系（relation to set estimator），\hyperref[idxocc:171:369]{\pageref*{idxocc:171:369}}\par
\needspace{3\baselineskip}
\par\noindent\hangindent=2em\hspace*{1em}\phantomsection\label{idx:172}最短（shortest）\par
\par\noindent\hangindent=2.8em\hspace*{1.8em}\phantomsection\label{idx:173}二项（binomial），\hyperref[ex:9.2.11]{\pageref*{ex:9.2.11}}, \hyperref[exercise:9.18]{\pageref*{exercise:9.18}}\par
\needspace{3\baselineskip}
\par\noindent\hangindent=2em\hspace*{1em}\phantomsection\label{idx:174}一致最准确（UMA）\par
\par\noindent\hangindent=2.8em\hspace*{1.8em}\phantomsection\label{idx:175}长度优势（length dominance），\hyperref[exercise:9.45]{\pageref*{exercise:9.45}}\par
\par\noindent\hangindent=2.8em\hspace*{1.8em}\phantomsection\label{idx:176}正态（normal），\hyperref[exercise:9.47]{\pageref*{exercise:9.47}}\par
\par\noindent\hangindent=2em\hspace*{1em}\phantomsection\label{idx:177}均匀（uniform），\hyperref[ex:9.1.6]{\pageref*{ex:9.1.6}}\par
\par\noindent\hangindent=2em\hspace*{1em}\phantomsection\label{idx:178}方差（variance），\hyperref[exercise:9.52]{\pageref*{exercise:9.52}}\par
\par\addvspace{2pt}
\par\noindent\hangindent=1em\hspace*{0em}\phantomsection\label{idx:179}置信区间（另见\hyperref[idx:181]{置信集}；\hyperref[idx:180]{置信程序}；\hyperref[idx:230]{可信集}）（Confidence interval (see also Confidence set; Confidence procedure; Credible Set)），\hyperref[def:9.1.1]{\pageref*{def:9.1.1}}\par
\par\addvspace{2pt}
\par\noindent\hangindent=1em\hspace*{0em}\phantomsection\label{idx:180}置信程序（另见\hyperref[idx:181]{置信集}；\hyperref[idx:129]{置信区间}）（Confidence procedure (see also Confidence set; Confidence interval)），\hyperref[sec:9.5.1]{\pageref*{sec:9.5.1}}\par
\par\addvspace{2pt}
\needspace{3\baselineskip}
\par\noindent\hangindent=1em\hspace*{0em}\phantomsection\label{idx:181}置信集（Confidence set）\par
\par\noindent\hangindent=2em\hspace*{1em}\phantomsection\label{idx:182}变量含误差（EIV），\hyperref[sec:12.2.4]{\pageref*{sec:12.2.4}}\par
\needspace{3\baselineskip}
\par\noindent\hangindent=2em\hspace*{1em}\phantomsection\label{idx:183}频率学派与贝叶斯（frequentist–Bayesian）\par
\par\noindent\hangindent=2.8em\hspace*{1.8em}\phantomsection\label{idx:184}解释（interpretations），\hyperref[sec:9.2.4]{\pageref*{sec:9.2.4}}\par
\needspace{3\baselineskip}
\par\noindent\hangindent=2em\hspace*{1em}\phantomsection\label{idx:185}枢轴（pivotal）\par
\par\noindent\hangindent=2.8em\hspace*{1.8em}\phantomsection\label{idx:186}变量含误差（EIV），\hyperref[eq:12.2.24]{\pageref*{eq:12.2.24}}\par
\par\noindent\hangindent=2em\hspace*{1em}\phantomsection\label{idx:187}与假设检验的关系（relation to hypothesis test），\hyperref[thm:9.2.2]{\pageref*{thm:9.2.2}}\par
\sourcepageanchor{558}
\par\noindent\hangindent=2em\hspace*{1em}\phantomsection\label{idx:188}大小（size），\hyperref[idxocc:188:387]{\pageref*{idxocc:188:387}}\par
\par\noindent\hangindent=2em\hspace*{1em}\phantomsection\label{idx:189}一致最准确（UMA），\hyperref[idxocc:189:391]{\pageref*{idxocc:189:391}}\par
\par\noindent\hangindent=2.8em\hspace*{1.8em}\phantomsection\label{idx:190}正态（normal），\hyperref[ex:9.3.6]{\pageref*{ex:9.3.6}}\par
\par\noindent\hangindent=2.8em\hspace*{1.8em}\phantomsection\label{idx:191}与一致最强检验的关系（relation to UMP test），\hyperref[thm:9.3.5]{\pageref*{thm:9.3.5}}\par
\par\noindent\hangindent=2em\hspace*{1em}\phantomsection\label{idx:192}无偏（unbiased），\hyperref[def:9.3.7]{\pageref*{def:9.3.7}}\par
\par\noindent\hangindent=2.8em\hspace*{1.8em}\phantomsection\label{idx:193}正态（normal），\hyperref[ex:9.3.8]{\pageref*{ex:9.3.8}}\par
\par\noindent\hangindent=2.8em\hspace*{1.8em}\phantomsection\label{idx:194}与无偏检验的关系（relation to unbiased test），\hyperref[idxocc:194:392]{\pageref*{idxocc:194:392}}\par
\par\addvspace{2pt}
\par\noindent\hangindent=1em\hspace*{0em}\phantomsection\label{idx:195}共轭族（Conjugate family），\hyperref[def:7.2.15]{\pageref*{def:7.2.15}}\par
\par\addvspace{2pt}
\par\noindent\hangindent=1em\hspace*{0em}\phantomsection\label{idx:196}相合性（见\hyperref[idx:770]{点估计量的相合性}）（Consistency (see Point estimator consistent)），\hyperref[ex:5.5.3]{\pageref*{ex:5.5.3}}\par
\par\addvspace{2pt}
\par\noindent\hangindent=1em\hspace*{0em}\phantomsection\label{idx:197}列联表（Contingency table），\hyperref[exercise:10.31]{\pageref*{exercise:10.31}}\par
\par\addvspace{2pt}
\par\noindent\hangindent=1em\hspace*{0em}\phantomsection\label{idx:198}连续性校正（Continuity correction），\hyperref[ex:3.3.2]{\pageref*{ex:3.3.2}}\par
\par\addvspace{2pt}
\par\noindent\hangindent=1em\hspace*{0em}\phantomsection\label{idx:199}对比（Contrasts），\hyperref[def:11.2.4]{\pageref*{def:11.2.4}}\par
\par\addvspace{2pt}
\needspace{3\baselineskip}
\par\noindent\hangindent=1em\hspace*{0em}\phantomsection\label{idx:200}收敛（Convergence）\par
\par\noindent\hangindent=2em\hspace*{1em}\phantomsection\label{idx:201}几乎处处（almost sure），\hyperref[def:5.5.6]{\pageref*{def:5.5.6}}\par
\par\noindent\hangindent=2em\hspace*{1em}\phantomsection\label{idx:202}几乎处处（almost surely），\hyperref[sec:5.8.4]{\pageref*{sec:5.8.4}}\par
\par\noindent\hangindent=2em\hspace*{1em}\phantomsection\label{idx:203}遍历定理（ergodic theorem），\hyperref[idxocc:203:250]{\pageref*{idxocc:203:250}}\par
\par\noindent\hangindent=2em\hspace*{1em}\phantomsection\label{idx:204}依分布（in distribution），\hyperref[def:5.5.10]{\pageref*{def:5.5.10}}, \hyperref[thm:5.5.17]{\pageref*{thm:5.5.17}}, \hyperref[sec:10.3.2]{\pageref*{sec:10.3.2}}\par
\par\noindent\hangindent=2.8em\hspace*{1.8em}\phantomsection\label{idx:205}与依概率收敛的关系（relation to convergence in probability），\hyperref[thm:5.5.12]{\pageref*{thm:5.5.12}}\par
\par\noindent\hangindent=2em\hspace*{1em}\phantomsection\label{idx:206}依概率（in probability），\hyperref[def:5.5.1]{\pageref*{def:5.5.1}}\par
\par\noindent\hangindent=2.8em\hspace*{1.8em}\phantomsection\label{idx:207}与依分布收敛的关系（relation to convergence in distribution），\hyperref[thm:5.5.12]{\pageref*{thm:5.5.12}}\par
\par\noindent\hangindent=2.8em\hspace*{1.8em}\phantomsection\label{idx:208}收敛到常数（to a constant），\hyperref[thm:5.5.13]{\pageref*{thm:5.5.13}}\par
\par\noindent\hangindent=2em\hspace*{1em}\phantomsection\label{idx:209}强大数定律（strong law），\hyperref[thm:5.5.9]{\pageref*{thm:5.5.9}}\par
\par\noindent\hangindent=2em\hspace*{1em}\phantomsection\label{idx:210}弱大数定律（weak law），\hyperref[thm:5.5.2]{\pageref*{thm:5.5.2}}\par
\par\noindent\hangindent=2em\hspace*{1em}\phantomsection\label{idx:211}以概率1（with probability one），\hyperref[def:5.5.6]{\pageref*{def:5.5.6}}\par
\par\addvspace{2pt}
\par\noindent\hangindent=1em\hspace*{0em}\phantomsection\label{idx:212}凸函数（Convex function），\hyperref[def:4.7.6]{\pageref*{def:4.7.6}}\par
\par\addvspace{2pt}
\par\noindent\hangindent=1em\hspace*{0em}\phantomsection\label{idx:213}卷积（Convolution），\hyperref[thm:5.2.9]{\pageref*{thm:5.2.9}}\par
\par\addvspace{2pt}
\needspace{3\baselineskip}
\par\noindent\hangindent=1em\hspace*{0em}\phantomsection\label{idx:214}相关（Correlation）\par
\par\noindent\hangindent=2em\hspace*{1em}\phantomsection\label{idx:215}二元（binary），\hyperref[exercise:5.57]{\pageref*{exercise:5.57}}\par
\par\noindent\hangindent=2em\hspace*{1em}\phantomsection\label{idx:216}系数（coefficient），\hyperref[def:4.5.2]{\pageref*{def:4.5.2}}\par
\par\addvspace{2pt}
\par\noindent\hangindent=1em\hspace*{0em}\phantomsection\label{idx:217}相关（另见\hyperref[idx:220]{协方差}）（Correlation (see also Covariance)），\hyperref[def:4.5.2]{\pageref*{def:4.5.2}}, \hyperref[thm:4.5.7]{\pageref*{thm:4.5.7}}\par
\par\addvspace{2pt}
\par\noindent\hangindent=1em\hspace*{0em}\phantomsection\label{idx:218}计数公式（Counting formulas），\hyperref[tab:1.2.1]{\pageref*{tab:1.2.1}}\par
\par\addvspace{2pt}
\par\noindent\hangindent=1em\hspace*{0em}\phantomsection\label{idx:219}计数方法（Counting methods），\hyperref[sec:1.2.3]{\pageref*{sec:1.2.3}}\par
\par\addvspace{2pt}
\needspace{3\baselineskip}
\par\noindent\hangindent=1em\hspace*{0em}\phantomsection\label{idx:220}协方差（Covariance）\par
\par\noindent\hangindent=2em\hspace*{1em}\phantomsection\label{idx:221}独立性（independence），\hyperref[thm:4.5.5]{\pageref*{thm:4.5.5}}\par
\par\noindent\hangindent=2em\hspace*{1em}\phantomsection\label{idx:222}线性函数（linear functions），\hyperref[item:12.11.a]{\pageref*{item:12.11.a}}\par
\par\addvspace{2pt}
\par\noindent\hangindent=1em\hspace*{0em}\phantomsection\label{idx:223}协方差（另见\hyperref[idx:214]{相关}）（Covariance (see also Correlation)），\hyperref[def:4.5.1]{\pageref*{def:4.5.1}}\par
\par\addvspace{2pt}
\par\noindent\hangindent=1em\hspace*{0em}\phantomsection\label{idx:224}覆盖概率（Coverage probability），\hyperref[def:9.1.4]{\pageref*{def:9.1.4}}\par
\par\noindent\hangindent=2em\hspace*{1em}\phantomsection\label{idx:225}频率学派与贝叶斯解释（frequentist–Bayesian interpretation），\hyperref[idxocc:225:384]{\pageref*{idxocc:225:384}}\par
\par\noindent\hangindent=2em\hspace*{1em}\phantomsection\label{idx:226}与风险的关系（relation to risk），\hyperref[idxocc:226:395]{\pageref*{idxocc:226:395}}\par
\par\addvspace{2pt}
\par\noindent\hangindent=1em\hspace*{0em}\phantomsection\label{idx:227}Cox悖论（Cox’s paradox），\hyperref[exercise:9.49]{\pageref*{exercise:9.49}}\par
\par\addvspace{2pt}
\par\noindent\hangindent=1em\hspace*{0em}\phantomsection\label{idx:228}Cramér–Rao下界（Cramér–Rao Lower Bound），\hyperref[thm:7.3.9]{\pageref*{thm:7.3.9}}\par
\par\noindent\hangindent=2em\hspace*{1em}\phantomsection\label{idx:229}必要条件（necessary condition），\hyperref[cor:7.3.15]{\pageref*{cor:7.3.15}}\par
\par\addvspace{2pt}
\needspace{3\baselineskip}
\par\noindent\hangindent=1em\hspace*{0em}\phantomsection\label{idx:230}可信集（Credible set）\par
\par\noindent\hangindent=2em\hspace*{1em}\phantomsection\label{idx:231}正态（normal），\hyperref[ex:9.2.18]{\pageref*{ex:9.2.18}}\par
\par\noindent\hangindent=2em\hspace*{1em}\phantomsection\label{idx:232}泊松（Poisson），\hyperref[ex:9.2.16]{\pageref*{ex:9.2.16}}\par
\needspace{3\baselineskip}
\par\noindent\hangindent=2em\hspace*{1em}\phantomsection\label{idx:233}与置信集的关系（relation to confidence set）\par
\par\noindent\hangindent=2.8em\hspace*{1.8em}\phantomsection\label{idx:234}正态（normal），\hyperref[idxocc:234:386]{\pageref*{idxocc:234:386}}, \hyperref[exercise:9.32]{\pageref*{exercise:9.32}}\par
\par\addvspace{2pt}
\par\noindent\hangindent=1em\hspace*{0em}\phantomsection\label{idx:235}可信集（另见\hyperref[idx:181]{置信集}；\hyperref[idx:129]{置信区间}；\hyperref[idx:180]{置信程序}）（Credible set (see also Confidence set, Confidence interval, Confidence procedure），\hyperref[eq:9.2.18]{\pageref*{eq:9.2.18}}\par
\par\addvspace{2pt}
\par\noindent\hangindent=1em\hspace*{0em}\phantomsection\label{idx:236}直积集（Cross product set），\hyperref[idxocc:236:157]{\pageref*{idxocc:236:157}}\par
\par\addvspace{2pt}
\par\noindent\hangindent=1em\hspace*{0em}\phantomsection\label{idx:237}累积分布函数（Cumulative distribution function (cdf)），\hyperref[def:1.5.1]{\pageref*{def:1.5.1}}\par
\par\noindent\hangindent=2em\hspace*{1em}\phantomsection\label{idx:238}独立性（independence），\hyperref[exercise:4.9]{\pageref*{exercise:4.9}}\par
\par\noindent\hangindent=2em\hspace*{1em}\phantomsection\label{idx:239}联合（joint），\hyperref[idxocc:239:152]{\pageref*{idxocc:239:152}}\par
\par\noindent\hangindent=2em\hspace*{1em}\phantomsection\label{idx:240}混合（mixed），\hyperref[ex:1.5.6]{\pageref*{ex:1.5.6}}\par
\par\noindent\hangindent=2em\hspace*{1em}\phantomsection\label{idx:241}单调变换（monotone transformation），\hyperref[thm:2.1.3]{\pageref*{thm:2.1.3}}\par
\par\noindent\hangindent=2em\hspace*{1em}\phantomsection\label{idx:242}性质（properties），\hyperref[thm:1.5.3]{\pageref*{thm:1.5.3}}\par
\par\addvspace{2pt}
\par\noindent\hangindent=1em\hspace*{0em}\phantomsection\label{idx:243}临界点（Cutoff point），\hyperref[idxocc:243:344]{\pageref*{idxocc:243:344}}\par
\par\addvspace{2pt}
\par\noindent\hangindent=1em\hspace*{0em}\phantomsection\label{idx:244}决策论（Decision theory），\hyperref[sec:7.3.4]{\pageref*{sec:7.3.4}}, \hyperref[sec:8.3.5]{\pageref*{sec:8.3.5}}\par
\par\addvspace{2pt}
\par\noindent\hangindent=1em\hspace*{0em}\phantomsection\label{idx:245}自由度（Degrees of freedom），\hyperref[eq:3.3.10]{\pageref*{eq:3.3.10}}\par
\par\addvspace{2pt}
\par\noindent\hangindent=1em\hspace*{0em}\phantomsection\label{idx:246}Delta方法（Delta method），\hyperref[sec:5.5.4]{\pageref*{sec:5.5.4}}\par
\par\addvspace{2pt}
\par\noindent\hangindent=1em\hspace*{0em}\phantomsection\label{idx:247}德摩根律（见\hyperref[idx:959]{集合运算}）（DeMorgan’s Laws (see Set operation)），\hyperref[idxocc:247:36]{\pageref*{idxocc:247:36}}, \hyperref[exercise:1.9]{\pageref*{exercise:1.9}}\par
\par\addvspace{2pt}
\par\noindent\hangindent=1em\hspace*{0em}\phantomsection\label{idx:248}行列式（Determinant），\hyperref[idxocc:248:160]{\pageref*{idxocc:248:160}}, \hyperref[idxocc:248:182]{\pageref*{idxocc:248:182}}\par
\par\addvspace{2pt}
\par\noindent\hangindent=1em\hspace*{0em}\phantomsection\label{idx:249}不交集（Disjoint sets），\hyperref[def:1.1.5]{\pageref*{def:1.1.5}}\par
\par\noindent\hangindent=2em\hspace*{1em}\phantomsection\label{idx:250}独立性（independence），\hyperref[exercise:1.39]{\pageref*{exercise:1.39}}\par
\par\addvspace{2pt}
\needspace{3\baselineskip}
\par\noindent\hangindent=1em\hspace*{0em}\phantomsection\label{idx:251}分布（Distribution）\par
\par\noindent\hangindent=2em\hspace*{1em}\phantomsection\label{idx:252}伯努利（Bernoulli），\hyperref[eq:3.2.3]{\pageref*{eq:3.2.3}}\par
\par\noindent\hangindent=2.8em\hspace*{1.8em}\phantomsection\label{idx:253}极限（limit），\hyperref[exercise:5.44]{\pageref*{exercise:5.44}}\par
\par\noindent\hangindent=2em\hspace*{1em}\phantomsection\label{idx:254}贝塔（beta），\hyperref[eq:3.3.16]{\pageref*{eq:3.3.16}}\par
\par\noindent\hangindent=2.8em\hspace*{1.8em}\phantomsection\label{idx:255}生成（generating），\hyperref[eq:5.6.5]{\pageref*{eq:5.6.5}}, \hyperref[ex:5.6.7]{\pageref*{ex:5.6.7}}, \hyperref[ex:5.6.9]{\pageref*{ex:5.6.9}}, \hyperref[exercise:5.61]{\pageref*{exercise:5.61}}\par
\par\noindent\hangindent=2.8em\hspace*{1.8em}\phantomsection\label{idx:256}乘积（product），\hyperref[ex:4.3.3]{\pageref*{ex:4.3.3}}\par
\par\noindent\hangindent=2.8em\hspace*{1.8em}\phantomsection\label{idx:257}与二项的关系（relation to binomial），\hyperref[exercise:2.40]{\pageref*{exercise:2.40}}\par
\par\noindent\hangindent=2.8em\hspace*{1.8em}\phantomsection\label{idx:258}与F的关系（relation to F），\hyperref[thm:5.3.8]{\pageref*{thm:5.3.8}}\par
\par\noindent\hangindent=2.8em\hspace*{1.8em}\phantomsection\label{idx:259}与伽马的关系（relation to gamma），\hyperref[exercise:4.24]{\pageref*{exercise:4.24}}\par
\par\noindent\hangindent=2em\hspace*{1em}\phantomsection\label{idx:260}贝塔二项（beta–binomial），\hyperref[exercise:4.34]{\pageref*{exercise:4.34}}\par
\par\noindent\hangindent=2em\hspace*{1em}\phantomsection\label{idx:261}二项（binomial），\hyperref[ex:2.1.1]{\pageref*{ex:2.1.1}}, \hyperref[idxocc:261:104]{\pageref*{idxocc:261:104}}\par
\par\noindent\hangindent=2.8em\hspace*{1.8em}\phantomsection\label{idx:262}指数族（exponential family），\hyperref[ex:3.4.1]{\pageref*{ex:3.4.1}}\par
\par\noindent\hangindent=2.8em\hspace*{1.8em}\phantomsection\label{idx:263}生成（generating），\hyperref[ex:5.6.5]{\pageref*{ex:5.6.5}}\par
\par\noindent\hangindent=2.8em\hspace*{1.8em}\phantomsection\label{idx:264}均值（mean），\hyperref[ex:2.2.3]{\pageref*{ex:2.2.3}}\par
\par\noindent\hangindent=2.8em\hspace*{1.8em}\phantomsection\label{idx:265}矩母函数（mgf），\hyperref[ex:2.3.9]{\pageref*{ex:2.3.9}}\par
\par\noindent\hangindent=2.8em\hspace*{1.8em}\phantomsection\label{idx:266}矩（moments），\hyperref[ex:3.4.3]{\pageref*{ex:3.4.3}}\par
\par\noindent\hangindent=2.8em\hspace*{1.8em}\phantomsection\label{idx:267}递推关系（recursion relation），\hyperref[eq:3.2.7]{\pageref*{eq:3.2.7}}\par
\par\noindent\hangindent=2.8em\hspace*{1.8em}\phantomsection\label{idx:268}与贝塔的关系（relation to beta），\hyperref[exercise:2.40]{\pageref*{exercise:2.40}}\par
\par\noindent\hangindent=2.8em\hspace*{1.8em}\phantomsection\label{idx:269}与多项的关系（relation to multinomial），\hyperref[idxocc:269:179]{\pageref*{idxocc:269:179}}\par
\par\noindent\hangindent=2.8em\hspace*{1.8em}\phantomsection\label{idx:270}与负二项的关系（relation to negative binomial），\hyperref[exercise:3.12]{\pageref*{exercise:3.12}}\par
\par\noindent\hangindent=2.8em\hspace*{1.8em}\phantomsection\label{idx:271}与泊松的关系（relation to Poisson），\hyperref[exercise:4.15]{\pageref*{exercise:4.15}}\par
\par\noindent\hangindent=2.8em\hspace*{1.8em}\phantomsection\label{idx:272}方差（variance），\hyperref[ex:2.3.5]{\pageref*{ex:2.3.5}}\par
\par\noindent\hangindent=2em\hspace*{1em}\phantomsection\label{idx:273}二元正态（bivariate normal），\hyperref[def:4.5.10]{\pageref*{def:4.5.10}}, \hyperref[exercise:4.45]{\pageref*{exercise:4.45}}\par
\par\noindent\hangindent=2.8em\hspace*{1.8em}\phantomsection\label{idx:274}相关（correlation），\hyperref[item:10.17.d]{\pageref*{item:10.17.d}}\par
\par\noindent\hangindent=2em\hspace*{1em}\phantomsection\label{idx:275}二元t（bivariate t），\hyperref[idxocc:275:479]{\pageref*{idxocc:275:479}}\par
\par\noindent\hangindent=2em\hspace*{1em}\phantomsection\label{idx:276}柯西（Cauchy），\hyperref[eq:3.3.19]{\pageref*{eq:3.3.19}}, \hyperref[exercise:3.39]{\pageref*{exercise:3.39}}\par
\par\noindent\hangindent=2.8em\hspace*{1.8em}\phantomsection\label{idx:277}生成（generation），\hyperref[exercise:2.12]{\pageref*{exercise:2.12}}\par
\par\noindent\hangindent=2.8em\hspace*{1.8em}\phantomsection\label{idx:278}均值（mean），\hyperref[ex:2.2.4]{\pageref*{ex:2.2.4}}\par
\par\noindent\hangindent=2.8em\hspace*{1.8em}\phantomsection\label{idx:279}与正态的关系（relation to normal），\hyperref[ex:4.3.6]{\pageref*{ex:4.3.6}}, \hyperref[exercise:4.28]{\pageref*{exercise:4.28}}\par
\par\noindent\hangindent=2em\hspace*{1em}\phantomsection\label{idx:280}卡方（chi squared），\hyperref[eq:3.3.10]{\pageref*{eq:3.3.10}}\par
\par\noindent\hangindent=2.8em\hspace*{1.8em}\phantomsection\label{idx:281}生成（generating），\hyperref[eq:5.6.5]{\pageref*{eq:5.6.5}}\par
\par\noindent\hangindent=2.8em\hspace*{1.8em}\phantomsection\label{idx:282}性质（properties），\hyperref[lem:5.3.2]{\pageref*{lem:5.3.2}}\par
\par\noindent\hangindent=2.8em\hspace*{1.8em}\phantomsection\label{idx:283}与负二项的关系（relation to negative binomial），\hyperref[exercise:2.38]{\pageref*{exercise:2.38}}\par
\par\noindent\hangindent=2.8em\hspace*{1.8em}\phantomsection\label{idx:284}与正态的关系（relation to normal），\hyperref[ex:2.1.9]{\pageref*{ex:2.1.9}}\par
\sourcepageanchor{559}
\needspace{3\baselineskip}
\par\noindent\hangindent=2em\hspace*{1em}\phantomsection\label{idx:285}离散（discrete）\par
\par\noindent\hangindent=2.8em\hspace*{1.8em}\phantomsection\label{idx:286}生成（generating），\hyperref[ex:5.6.5]{\pageref*{ex:5.6.5}}\par
\par\noindent\hangindent=2em\hspace*{1em}\phantomsection\label{idx:287}离散均匀（discrete uniform），\hyperref[eq:3.2.1]{\pageref*{eq:3.2.1}}\par
\par\noindent\hangindent=2em\hspace*{1em}\phantomsection\label{idx:288}双指数（double exponential），\hyperref[eq:3.3.23]{\pageref*{eq:3.3.23}}\par
\par\noindent\hangindent=2em\hspace*{1em}\phantomsection\label{idx:289}指数（exponential），\hyperref[eq:3.3.11]{\pageref*{eq:3.3.11}}\par
\par\noindent\hangindent=2.8em\hspace*{1.8em}\phantomsection\label{idx:290}生成（generating），\hyperref[ex:5.6.3]{\pageref*{ex:5.6.3}}\par
\par\noindent\hangindent=2.8em\hspace*{1.8em}\phantomsection\label{idx:291}风险率函数（hazard function），\hyperref[idxocc:291:139]{\pageref*{idxocc:291:139}}\par
\par\noindent\hangindent=2.8em\hspace*{1.8em}\phantomsection\label{idx:292}均值（mean），\hyperref[ex:2.2.2]{\pageref*{ex:2.2.2}}\par
\par\noindent\hangindent=2.8em\hspace*{1.8em}\phantomsection\label{idx:293}无记忆性（memoryless），\hyperref[idxocc:293:114]{\pageref*{idxocc:293:114}}\par
\par\noindent\hangindent=2.8em\hspace*{1.8em}\phantomsection\label{idx:294}递推关系（recursion relation），\hyperref[eq:2.4.7]{\pageref*{eq:2.4.7}}\par
\par\noindent\hangindent=2.8em\hspace*{1.8em}\phantomsection\label{idx:295}与伽马的关系（relation to gamma），\hyperref[exercise:4.38]{\pageref*{exercise:4.38}}\par
\par\noindent\hangindent=2.8em\hspace*{1.8em}\phantomsection\label{idx:296}与均匀的关系（relation to uniform），\hyperref[ex:2.1.4]{\pageref*{ex:2.1.4}}\par
\par\noindent\hangindent=2.8em\hspace*{1.8em}\phantomsection\label{idx:297}方差（variance），\hyperref[ex:2.3.3]{\pageref*{ex:2.3.3}}\par
\needspace{3\baselineskip}
\par\noindent\hangindent=2em\hspace*{1em}\phantomsection\label{idx:298}指数族（exponential family）\par
\par\noindent\hangindent=2.8em\hspace*{1.8em}\phantomsection\label{idx:299}曲线（curved），\hyperref[idxocc:299:208]{\pageref*{idxocc:299:208}}\par
\par\noindent\hangindent=2em\hspace*{1em}\phantomsection\label{idx:300}极值（extreme value），\hyperref[idxocc:300:139]{\pageref*{idxocc:300:139}}\par
\par\noindent\hangindent=2em\hspace*{1em}\phantomsection\label{idx:301}F，\hyperref[def:5.3.6]{\pageref*{def:5.3.6}}\par
\par\noindent\hangindent=2.8em\hspace*{1.8em}\phantomsection\label{idx:302}倒数（reciprocal），\hyperref[thm:5.3.8]{\pageref*{thm:5.3.8}}\par
\par\noindent\hangindent=2.8em\hspace*{1.8em}\phantomsection\label{idx:303}与贝塔的关系（relation to beta），\hyperref[thm:5.3.8]{\pageref*{thm:5.3.8}}\par
\par\noindent\hangindent=2.8em\hspace*{1.8em}\phantomsection\label{idx:304}与t的关系（relation to t），\hyperref[thm:5.3.8]{\pageref*{thm:5.3.8}}\par
\par\noindent\hangindent=2em\hspace*{1em}\phantomsection\label{idx:305}折叠正态（folded normal），\hyperref[exercise:3.20]{\pageref*{exercise:3.20}}\par
\par\noindent\hangindent=2em\hspace*{1em}\phantomsection\label{idx:306}伽马（gamma），\hyperref[eq:3.3.6]{\pageref*{eq:3.3.6}}\par
\par\noindent\hangindent=2.8em\hspace*{1.8em}\phantomsection\label{idx:307}生成（generating），\hyperref[eq:5.6.5]{\pageref*{eq:5.6.5}}, \hyperref[exercise:5.61]{\pageref*{exercise:5.61}}\par
\par\noindent\hangindent=2.8em\hspace*{1.8em}\phantomsection\label{idx:308}矩母函数（mgf），\hyperref[ex:2.3.8]{\pageref*{ex:2.3.8}}\par
\par\noindent\hangindent=2.8em\hspace*{1.8em}\phantomsection\label{idx:309}矩（moments），\hyperref[exercise:3.17]{\pageref*{exercise:3.17}}\par
\par\noindent\hangindent=2.8em\hspace*{1.8em}\phantomsection\label{idx:310}递推关系（recursion relation），\hyperref[thm:3.6.4]{\pageref*{thm:3.6.4}}\par
\par\noindent\hangindent=2.8em\hspace*{1.8em}\phantomsection\label{idx:311}与贝塔的关系（relation to beta），\hyperref[exercise:4.24]{\pageref*{exercise:4.24}}\par
\par\noindent\hangindent=2.8em\hspace*{1.8em}\phantomsection\label{idx:312}与指数的关系（relation to exponential），\hyperref[exercise:4.38]{\pageref*{exercise:4.38}}\par
\par\noindent\hangindent=2.8em\hspace*{1.8em}\phantomsection\label{idx:313}与逆伽马的关系（relation to inverted gamma），\hyperref[ex:2.1.6]{\pageref*{ex:2.1.6}}\par
\par\noindent\hangindent=2.8em\hspace*{1.8em}\phantomsection\label{idx:314}与负二项的关系（relation to negative binomial），\hyperref[exercise:3.18]{\pageref*{exercise:3.18}}\par
\par\noindent\hangindent=2.8em\hspace*{1.8em}\phantomsection\label{idx:315}与泊松的关系（relation to Poisson），\hyperref[ex:3.3.1]{\pageref*{ex:3.3.1}}, \hyperref[exercise:3.19]{\pageref*{exercise:3.19}}\par
\par\noindent\hangindent=2.8em\hspace*{1.8em}\phantomsection\label{idx:316}和（sum），\hyperref[ex:4.6.8]{\pageref*{ex:4.6.8}}\par
\par\noindent\hangindent=2em\hspace*{1em}\phantomsection\label{idx:317}高斯（Gaussian），\hyperref[idxocc:317:114]{\pageref*{idxocc:317:114}}\par
\par\noindent\hangindent=2em\hspace*{1em}\phantomsection\label{idx:318}几何（geometric），\hyperref[idxocc:318:111]{\pageref*{idxocc:318:111}}\par
\par\noindent\hangindent=2.8em\hspace*{1.8em}\phantomsection\label{idx:319}累积分布函数（cdf），\hyperref[ex:1.5.4]{\pageref*{ex:1.5.4}}\par
\par\noindent\hangindent=2.8em\hspace*{1.8em}\phantomsection\label{idx:320}无记忆性（memoryless），\hyperref[eq:3.2.11]{\pageref*{eq:3.2.11}}\par
\par\noindent\hangindent=2.8em\hspace*{1.8em}\phantomsection\label{idx:321}概率质量函数（pmf），\hyperref[ex:1.6.2]{\pageref*{ex:1.6.2}}\par
\par\noindent\hangindent=2em\hspace*{1em}\phantomsection\label{idx:322}耿贝尔（Gumbel），\hyperref[idxocc:322:139]{\pageref*{idxocc:322:139}}\par
\par\noindent\hangindent=2em\hspace*{1em}\phantomsection\label{idx:323}超几何（hypergeometric），\hyperref[eq:3.2.2]{\pageref*{eq:3.2.2}}\par
\par\noindent\hangindent=2em\hspace*{1em}\phantomsection\label{idx:324}逆高斯（inverse Gaussian），\hyperref[exercise:6.27]{\pageref*{exercise:6.27}}, \hyperref[exercise:7.15]{\pageref*{exercise:7.15}}\par
\needspace{3\baselineskip}
\par\noindent\hangindent=2em\hspace*{1em}\phantomsection\label{idx:325}逆伽马（inverted gamma）\par
\par\noindent\hangindent=2.8em\hspace*{1.8em}\phantomsection\label{idx:326}与伽马的关系（relation to gamma），\hyperref[ex:2.1.6]{\pageref*{ex:2.1.6}}\par
\par\noindent\hangindent=2em\hspace*{1em}\phantomsection\label{idx:327}对数级数（logarithmic series），\hyperref[exercise:3.14]{\pageref*{exercise:3.14}}\par
\needspace{3\baselineskip}
\par\noindent\hangindent=2em\hspace*{1em}\phantomsection\label{idx:328}逻辑（logistic）\par
\par\noindent\hangindent=2.8em\hspace*{1.8em}\phantomsection\label{idx:329}累积分布函数（cdf），\hyperref[ex:1.5.5]{\pageref*{ex:1.5.5}}\par
\par\noindent\hangindent=2.8em\hspace*{1.8em}\phantomsection\label{idx:330}风险率函数（hazard function），\hyperref[idxocc:330:139]{\pageref*{idxocc:330:139}}\par
\par\noindent\hangindent=2.8em\hspace*{1.8em}\phantomsection\label{idx:331}概率密度函数（pdf），\hyperref[ex:1.6.4]{\pageref*{ex:1.6.4}}\par
\par\noindent\hangindent=2em\hspace*{1em}\phantomsection\label{idx:332}对数正态（lognormal），\hyperref[eq:3.3.21]{\pageref*{eq:3.3.21}}\par
\par\noindent\hangindent=2.8em\hspace*{1.8em}\phantomsection\label{idx:333}矩母函数（mgf），\hyperref[exercise:2.36]{\pageref*{exercise:2.36}}\par
\par\noindent\hangindent=2.8em\hspace*{1.8em}\phantomsection\label{idx:334}矩（moments），\hyperref[ex:2.3.10]{\pageref*{ex:2.3.10}}\par
\par\noindent\hangindent=2em\hspace*{1em}\phantomsection\label{idx:335}麦克斯韦（Maxwell），\hyperref[idxocc:335:139]{\pageref*{idxocc:335:139}}\par
\par\noindent\hangindent=2em\hspace*{1em}\phantomsection\label{idx:336}混合（mixture），\hyperref[def:4.4.4]{\pageref*{def:4.4.4}}\par
\par\noindent\hangindent=2.8em\hspace*{1.8em}\phantomsection\label{idx:337}伯努利–贝塔（Bernoulli–beta），\hyperref[ex:4.4.6]{\pageref*{ex:4.4.6}}, \hyperref[ex:4.4.8]{\pageref*{ex:4.4.8}}\par
\par\noindent\hangindent=2.8em\hspace*{1.8em}\phantomsection\label{idx:338}二项–贝塔（binomial–beta），\hyperref[exercise:4.34]{\pageref*{exercise:4.34}}\par
\par\noindent\hangindent=2.8em\hspace*{1.8em}\phantomsection\label{idx:339}二项–负二项（binomial–negative binomial），\hyperref[idxocc:339:167]{\pageref*{idxocc:339:167}}\par
\par\noindent\hangindent=2.8em\hspace*{1.8em}\phantomsection\label{idx:340}二项–泊松（binomial–Poisson），\hyperref[ex:4.4.1]{\pageref*{ex:4.4.1}}\par
\par\noindent\hangindent=2.8em\hspace*{1.8em}\phantomsection\label{idx:341}二项–泊松–指数（binomial–Poisson–exponential），\hyperref[ex:4.4.5]{\pageref*{ex:4.4.5}}\par
\par\noindent\hangindent=2.8em\hspace*{1.8em}\phantomsection\label{idx:342}卡方–泊松（chi squared–Poisson），\hyperref[eq:4.4.3]{\pageref*{eq:4.4.3}}\par
\par\noindent\hangindent=2.8em\hspace*{1.8em}\phantomsection\label{idx:343}指数–伽马（exponential-gamma），\hyperref[exercise:4.38]{\pageref*{exercise:4.38}}\par
\par\noindent\hangindent=2.8em\hspace*{1.8em}\phantomsection\label{idx:344}对数级数–泊松（logarithmic series–Poisson），\hyperref[exercise:4.33]{\pageref*{exercise:4.33}}\par
\par\noindent\hangindent=2.8em\hspace*{1.8em}\phantomsection\label{idx:345}正态–伽马（normal–gamma），\hyperref[exercise:5.20]{\pageref*{exercise:5.20}}\par
\par\noindent\hangindent=2.8em\hspace*{1.8em}\phantomsection\label{idx:346}泊松–伽马（Poisson–gamma），\hyperref[idxocc:346:167]{\pageref*{idxocc:346:167}}, \hyperref[exercise:4.32]{\pageref*{exercise:4.32}}\par
\par\noindent\hangindent=2.8em\hspace*{1.8em}\phantomsection\label{idx:347}t–F，\hyperref[exercise:5.20]{\pageref*{exercise:5.20}}\par
\par\noindent\hangindent=2em\hspace*{1em}\phantomsection\label{idx:348}多项（multinomial），\hyperref[def:4.6.2]{\pageref*{def:4.6.2}}\par
\par\noindent\hangindent=2.8em\hspace*{1.8em}\phantomsection\label{idx:349}协方差（covariance），\hyperref[idxocc:349:180]{\pageref*{idxocc:349:180}}, \hyperref[exercise:4.39]{\pageref*{exercise:4.39}}\par
\par\noindent\hangindent=2.8em\hspace*{1.8em}\phantomsection\label{idx:350}与二项的关系（relation to binomial），\hyperref[idxocc:350:179]{\pageref*{idxocc:350:179}}\par
\par\noindent\hangindent=2em\hspace*{1em}\phantomsection\label{idx:351}负二项（negative binomial），\hyperref[eq:3.2.10]{\pageref*{eq:3.2.10}}, \hyperref[exercise:4.32]{\pageref*{exercise:4.32}}\par
\par\noindent\hangindent=2.8em\hspace*{1.8em}\phantomsection\label{idx:352}二次方差（quadratic variance），\hyperref[idxocc:352:110]{\pageref*{idxocc:352:110}}\par
\par\noindent\hangindent=2.8em\hspace*{1.8em}\phantomsection\label{idx:353}与二项的关系（relation to binomial），\hyperref[exercise:3.12]{\pageref*{exercise:3.12}}\par
\par\noindent\hangindent=2.8em\hspace*{1.8em}\phantomsection\label{idx:354}与卡方的关系（relation to chi squared），\hyperref[exercise:2.38]{\pageref*{exercise:2.38}}\par
\par\noindent\hangindent=2.8em\hspace*{1.8em}\phantomsection\label{idx:355}与伽马的关系（relation to gamma），\hyperref[exercise:3.18]{\pageref*{exercise:3.18}}\par
\par\noindent\hangindent=2.8em\hspace*{1.8em}\phantomsection\label{idx:356}与泊松的关系（relation to Poisson），\hyperref[idxocc:356:110]{\pageref*{idxocc:356:110}}\par
\par\noindent\hangindent=2em\hspace*{1em}\phantomsection\label{idx:357}非中心卡方（noncentral chi squared），\hyperref[eq:4.4.3]{\pageref*{eq:4.4.3}}\par
\par\noindent\hangindent=2em\hspace*{1em}\phantomsection\label{idx:358}正态（normal），\hyperref[idxocc:358:114]{\pageref*{idxocc:358:114}}\par
\par\noindent\hangindent=2.8em\hspace*{1.8em}\phantomsection\label{idx:359}二元（bivariate），\hyperref[exercise:4.47]{\pageref*{exercise:4.47}}\par
\par\noindent\hangindent=2.8em\hspace*{1.8em}\phantomsection\label{idx:360}指数族（exponential family），\hyperref[ex:3.4.4]{\pageref*{ex:3.4.4}}\par
\par\noindent\hangindent=2.8em\hspace*{1.8em}\phantomsection\label{idx:361}生成（generating），\hyperref[eq:5.6.5]{\pageref*{eq:5.6.5}}\par
\par\noindent\hangindent=2.8em\hspace*{1.8em}\phantomsection\label{idx:362}线性组合（linear combination），\hyperref[cor:4.6.10]{\pageref*{cor:4.6.10}}, \hyperref[lem:5.3.3]{\pageref*{lem:5.3.3}}\par
\par\noindent\hangindent=2.8em\hspace*{1.8em}\phantomsection\label{idx:363}比值（ratio），\hyperref[ex:4.3.6]{\pageref*{ex:4.3.6}}\par
\par\noindent\hangindent=2.8em\hspace*{1.8em}\phantomsection\label{idx:364}递推关系（recursion relation），\hyperref[idxocc:364:90]{\pageref*{idxocc:364:90}}\par
\par\noindent\hangindent=2.8em\hspace*{1.8em}\phantomsection\label{idx:365}与柯西的关系（relation to Cauchy），\hyperref[ex:4.3.6]{\pageref*{ex:4.3.6}}, \hyperref[exercise:4.28]{\pageref*{exercise:4.28}}\par
\par\noindent\hangindent=2.8em\hspace*{1.8em}\phantomsection\label{idx:366}与卡方的关系（relation to chi squared），\hyperref[ex:2.1.9]{\pageref*{ex:2.1.9}}\par
\par\noindent\hangindent=2.8em\hspace*{1.8em}\phantomsection\label{idx:367}与t的关系（relation to t），\hyperref[exercise:5.18]{\pageref*{exercise:5.18}}\par
\par\noindent\hangindent=2.8em\hspace*{1.8em}\phantomsection\label{idx:368}标准（standard），\hyperref[idxocc:368:115]{\pageref*{idxocc:368:115}}\par
\par\noindent\hangindent=2.8em\hspace*{1.8em}\phantomsection\label{idx:369}和（sum），\hyperref[ex:4.2.13]{\pageref*{ex:4.2.13}}, \hyperref[ex:4.3.4]{\pageref*{ex:4.3.4}}\par
\par\noindent\hangindent=2em\hspace*{1em}\phantomsection\label{idx:370}帕累托（Pareto），\hyperref[exercise:3.23]{\pageref*{exercise:3.23}}\par
\par\noindent\hangindent=2em\hspace*{1em}\phantomsection\label{idx:371}泊松（Poisson），\hyperref[eq:3.2.5]{\pageref*{eq:3.2.5}}\par
\par\noindent\hangindent=2.8em\hspace*{1.8em}\phantomsection\label{idx:372}矩母函数（mgf），\hyperref[lem:2.3.14]{\pageref*{lem:2.3.14}}\par
\par\noindent\hangindent=2.8em\hspace*{1.8em}\phantomsection\label{idx:373}正态近似（normal approximation），\hyperref[ex:3.4.9]{\pageref*{ex:3.4.9}}\par
\par\noindent\hangindent=2.8em\hspace*{1.8em}\phantomsection\label{idx:374}基本假设（postulates），\hyperref[thm:3.8.1]{\pageref*{thm:3.8.1}}\par
\par\noindent\hangindent=2.8em\hspace*{1.8em}\phantomsection\label{idx:375}递推关系（recursion relation），\hyperref[eq:3.2.6]{\pageref*{eq:3.2.6}}, \hyperref[eq:3.6.2]{\pageref*{eq:3.6.2}}\par
\par\noindent\hangindent=2.8em\hspace*{1.8em}\phantomsection\label{idx:376}与二项的关系（relation to binomial），\hyperref[exercise:4.15]{\pageref*{exercise:4.15}}\par
\par\noindent\hangindent=2.8em\hspace*{1.8em}\phantomsection\label{idx:377}与伽马的关系（relation to gamma），\hyperref[ex:3.3.1]{\pageref*{ex:3.3.1}}, \hyperref[exercise:3.19]{\pageref*{exercise:3.19}}\par
\par\noindent\hangindent=2.8em\hspace*{1.8em}\phantomsection\label{idx:378}与负二项的关系（relation to negative binomial），\hyperref[idxocc:378:110]{\pageref*{idxocc:378:110}}, \hyperref[exercise:3.15]{\pageref*{exercise:3.15}}\par
\par\noindent\hangindent=2.8em\hspace*{1.8em}\phantomsection\label{idx:379}和（sum），\hyperref[thm:4.3.2]{\pageref*{thm:4.3.2}}\par
\par\noindent\hangindent=2em\hspace*{1em}\phantomsection\label{idx:380}瑞利（Rayleigh），\hyperref[idxocc:380:139]{\pageref*{idxocc:380:139}}\par
\par\noindent\hangindent=2em\hspace*{1em}\phantomsection\label{idx:381}t，\hyperref[def:5.3.4]{\pageref*{def:5.3.4}}\par
\par\noindent\hangindent=2.8em\hspace*{1.8em}\phantomsection\label{idx:382}非中心（noncentral），\hyperref[exercise:8.35]{\pageref*{exercise:8.35}}\par
\par\noindent\hangindent=2.8em\hspace*{1.8em}\phantomsection\label{idx:383}与F的关系（relation to F），\hyperref[thm:5.3.8]{\pageref*{thm:5.3.8}}\par
\par\noindent\hangindent=2.8em\hspace*{1.8em}\phantomsection\label{idx:384}与正态的关系（relation to normal），\hyperref[exercise:5.18]{\pageref*{exercise:5.18}}\par
\par\noindent\hangindent=2em\hspace*{1em}\phantomsection\label{idx:385}截断（truncated），\hyperref[exercise:3.13]{\pageref*{exercise:3.13}}\par
\par\noindent\hangindent=2em\hspace*{1em}\phantomsection\label{idx:386}Tukey，\hyperref[exercise:10.23]{\pageref*{exercise:10.23}}\par
\par\noindent\hangindent=2em\hspace*{1em}\phantomsection\label{idx:387}均匀（uniform），\hyperref[eq:3.3.1]{\pageref*{eq:3.3.1}}\par
\par\noindent\hangindent=2.8em\hspace*{1.8em}\phantomsection\label{idx:388}最大值（maximum），\hyperref[ex:5.5.11]{\pageref*{ex:5.5.11}}\par
\sourcepageanchor{560}
\par\noindent\hangindent=2.8em\hspace*{1.8em}\phantomsection\label{idx:389}次序统计量（order statistic），\hyperref[ex:5.4.5]{\pageref*{ex:5.4.5}}\par
\par\noindent\hangindent=2.8em\hspace*{1.8em}\phantomsection\label{idx:390}与指数的关系（relation to exponential），\hyperref[ex:2.1.4]{\pageref*{ex:2.1.4}}\par
\par\noindent\hangindent=2em\hspace*{1em}\phantomsection\label{idx:391}威布尔（Weibull），\hyperref[eq:3.3.12]{\pageref*{eq:3.3.12}}, \hyperref[idxocc:391:139]{\pageref*{idxocc:391:139}}\par
\par\noindent\hangindent=2.8em\hspace*{1.8em}\phantomsection\label{idx:392}风险率函数（hazard function），\hyperref[idxocc:392:139]{\pageref*{idxocc:392:139}}\par
\par\addvspace{2pt}
\par\noindent\hangindent=1em\hspace*{0em}\phantomsection\label{idx:393}点记号（Dot notation），\hyperref[idxocc:393:457]{\pageref*{idxocc:393:457}}\par
\par\addvspace{2pt}
\par\noindent\hangindent=1em\hspace*{0em}\phantomsection\label{idx:394}效率（另见\hyperref[idx:763]{点估计}）（Efficiency (see also Point estimation)），\hyperref[sec:10.1.2]{\pageref*{sec:10.1.2}}\par
\par\addvspace{2pt}
\par\noindent\hangindent=1em\hspace*{0em}\phantomsection\label{idx:395}EM算法（见\hyperref[idx:7]{算法：EM}）（EM Algorithm (see Algorithm, EM)），\hyperref[sec:7.2.4]{\pageref*{sec:7.2.4}}, \hyperref[sec:7.5.4]{\pageref*{sec:7.5.4}}\par
\par\addvspace{2pt}
\par\noindent\hangindent=1em\hspace*{0em}\phantomsection\label{idx:396}经验贝叶斯估计量（Empirical Bayes estimator），\hyperref[idxocc:396:331]{\pageref*{idxocc:396:331}}\par
\par\addvspace{2pt}
\par\noindent\hangindent=1em\hspace*{0em}\phantomsection\label{idx:397}空集（Empty set），\hyperref[idxocc:397:33]{\pageref*{idxocc:397:33}}\par
\par\addvspace{2pt}
\par\noindent\hangindent=1em\hspace*{0em}\phantomsection\label{idx:398}熵（Entropy），\hyperref[idxocc:398:196]{\pageref*{idxocc:398:196}}\par
\par\addvspace{2pt}
\needspace{3\baselineskip}
\par\noindent\hangindent=1em\hspace*{0em}\phantomsection\label{idx:399}等变性（Equivariance）\par
\par\noindent\hangindent=2em\hspace*{1em}\phantomsection\label{idx:400}二项（binomial），\hyperref[ex:6.4.1]{\pageref*{ex:6.4.1}}\par
\par\noindent\hangindent=2em\hspace*{1em}\phantomsection\label{idx:401}测量（measurement），\hyperref[idxocc:401:272]{\pageref*{idxocc:401:272}}, \hyperref[exercise:6.39]{\pageref*{exercise:6.39}}\par
\par\noindent\hangindent=2.8em\hspace*{1.8em}\phantomsection\label{idx:402}位置（location），\hyperref[exercise:6.41]{\pageref*{exercise:6.41}}\par
\par\addvspace{2pt}
\par\noindent\hangindent=1em\hspace*{0em}\phantomsection\label{idx:403}等变性（另见\hyperref[idx:587]{不变性}）（Equivariance (see also Invariance)），\hyperref[sec:6.4]{\pageref*{sec:6.4}}\par
\par\addvspace{2pt}
\par\noindent\hangindent=1em\hspace*{0em}\phantomsection\label{idx:404}等变性原则（Equivariance Principle），\hyperref[idxocc:404:272]{\pageref*{idxocc:404:272}}\par
\par\noindent\hangindent=2em\hspace*{1em}\phantomsection\label{idx:405}解释（interpretations），\hyperref[ex:6.4.6]{\pageref*{ex:6.4.6}}\par
\par\addvspace{2pt}
\par\noindent\hangindent=1em\hspace*{0em}\phantomsection\label{idx:406}遍历定理（Ergodic theorem），\hyperref[idxocc:406:250]{\pageref*{idxocc:406:250}}\par
\par\addvspace{2pt}
\par\noindent\hangindent=1em\hspace*{0em}\phantomsection\label{idx:407}变量含误差的回归（Errors in variables regression (EIV)），\hyperref[sec:12.2]{\pageref*{sec:12.2}}\par
\par\noindent\hangindent=2em\hspace*{1em}\phantomsection\label{idx:408}置信集（confidence set），\hyperref[sec:12.2.4]{\pageref*{sec:12.2.4}}\par
\par\noindent\hangindent=2em\hspace*{1em}\phantomsection\label{idx:409}相合性（consistency），\hyperref[idxocc:409:505]{\pageref*{idxocc:409:505}}, \hyperref[idxocc:409:507]{\pageref*{idxocc:409:507}}, \hyperref[sec:12.6.2]{\pageref*{sec:12.6.2}}\par
\par\noindent\hangindent=2em\hspace*{1em}\phantomsection\label{idx:410}函数关系（functional relationship），\hyperref[eq:12.2.5]{\pageref*{eq:12.2.5}}, \hyperref[sec:12.6.1]{\pageref*{sec:12.6.1}}\par
\par\noindent\hangindent=2em\hspace*{1em}\phantomsection\label{idx:411}可识别性（identifiability），\hyperref[idxocc:411:501]{\pageref*{idxocc:411:501}}, \hyperref[idxocc:411:506]{\pageref*{idxocc:411:506}}\par
\needspace{3\baselineskip}
\par\noindent\hangindent=2em\hspace*{1em}\phantomsection\label{idx:412}最大似然估计量（MLE）\par
\par\noindent\hangindent=2.8em\hspace*{1.8em}\phantomsection\label{idx:413}函数型（functional），\hyperref[eq:12.2.16]{\pageref*{eq:12.2.16}}\par
\par\noindent\hangindent=2.8em\hspace*{1.8em}\phantomsection\label{idx:414}与最小二乘的关系（relation to least squares），\hyperref[exercise:12.4]{\pageref*{exercise:12.4}}\par
\par\noindent\hangindent=2.8em\hspace*{1.8em}\phantomsection\label{idx:415}结构型（structural），\hyperref[idxocc:415:505]{\pageref*{idxocc:415:505}}\par
\par\noindent\hangindent=2em\hspace*{1em}\phantomsection\label{idx:416}与随机完全区组的关系（relation to RCB），\hyperref[idxocc:416:506]{\pageref*{idxocc:416:506}}\par
\par\noindent\hangindent=2em\hspace*{1em}\phantomsection\label{idx:417}结构关系（structural relationship），\hyperref[eq:12.2.6]{\pageref*{eq:12.2.6}}, \hyperref[sec:12.6.1]{\pageref*{sec:12.6.1}}\par
\par\addvspace{2pt}
\par\noindent\hangindent=1em\hspace*{0em}\phantomsection\label{idx:418}事件（Event），\hyperref[def:1.1.2]{\pageref*{def:1.1.2}}\par
\par\addvspace{2pt}
\par\noindent\hangindent=1em\hspace*{0em}\phantomsection\label{idx:419}证据（Evidence），\hyperref[idxocc:419:269]{\pageref*{idxocc:419:269}}\par
\par\addvspace{2pt}
\par\noindent\hangindent=1em\hspace*{0em}\phantomsection\label{idx:420}交换悖论（Exchange paradox），\hyperref[sec:4.9.1]{\pageref*{sec:4.9.1}}\par
\par\addvspace{2pt}
\par\noindent\hangindent=1em\hspace*{0em}\phantomsection\label{idx:421}可交换（Exchangeable），\hyperref[exercise:5.4]{\pageref*{exercise:5.4}}\par
\par\addvspace{2pt}
\par\noindent\hangindent=1em\hspace*{0em}\phantomsection\label{idx:422}期望值（Expected value），\hyperref[def:2.2.1]{\pageref*{def:2.2.1}}\par
\par\noindent\hangindent=2em\hspace*{1em}\phantomsection\label{idx:423}二元（bivariate），\hyperref[idxocc:423:147]{\pageref*{idxocc:423:147}}, \hyperref[idxocc:423:149]{\pageref*{idxocc:423:149}}\par
\par\noindent\hangindent=2.8em\hspace*{1.8em}\phantomsection\label{idx:424}性质（properties），\hyperref[exercise:4.2]{\pageref*{exercise:4.2}}\par
\par\noindent\hangindent=2em\hspace*{1em}\phantomsection\label{idx:425}条件（conditional），\hyperref[idxocc:425:154]{\pageref*{idxocc:425:154}}\par
\par\noindent\hangindent=2.8em\hspace*{1.8em}\phantomsection\label{idx:426}最小化（minimizing），\hyperref[exercise:4.13]{\pageref*{exercise:4.13}}\par
\par\noindent\hangindent=2em\hspace*{1em}\phantomsection\label{idx:427}独立性（independence），\hyperref[thm:4.2.10]{\pageref*{thm:4.2.10}}, \hyperref[thm:4.6.6]{\pageref*{thm:4.6.6}}\par
\par\noindent\hangindent=2em\hspace*{1em}\phantomsection\label{idx:428}迭代（iterated），\hyperref[thm:4.4.3]{\pageref*{thm:4.4.3}}\par
\par\noindent\hangindent=2em\hspace*{1em}\phantomsection\label{idx:429}联合（joint），\hyperref[eq:4.6.3]{\pageref*{eq:4.6.3}}\par
\par\noindent\hangindent=2em\hspace*{1em}\phantomsection\label{idx:430}最小化性质（minimizing property），\hyperref[ex:2.2.6]{\pageref*{ex:2.2.6}}\par
\par\noindent\hangindent=2em\hspace*{1em}\phantomsection\label{idx:431}极大值与极小值的（of maxima and minima），\hyperref[exercise:2.15]{\pageref*{exercise:2.15}}\par
\par\noindent\hangindent=2em\hspace*{1em}\phantomsection\label{idx:432}性质（properties），\hyperref[thm:2.2.5]{\pageref*{thm:2.2.5}}\par
\par\noindent\hangindent=2em\hspace*{1em}\phantomsection\label{idx:433}和（sum），\hyperref[lem:5.2.5]{\pageref*{lem:5.2.5}}\par
\par\noindent\hangindent=2em\hspace*{1em}\phantomsection\label{idx:434}两种计算方法（two calculation methods），\hyperref[eq:2.2.5]{\pageref*{eq:2.2.5}}\par
\par\addvspace{2pt}
\par\noindent\hangindent=1em\hspace*{0em}\phantomsection\label{idx:435}指数族（Exponential family），\hyperref[eq:3.4.1]{\pageref*{eq:3.4.1}}, \hyperref[exercise:3.28]{\pageref*{exercise:3.28}}\par
\par\noindent\hangindent=2em\hspace*{1em}\phantomsection\label{idx:436}曲线（curved），\hyperref[def:3.4.7]{\pageref*{def:3.4.7}}, \hyperref[idxocc:436:265]{\pageref*{idxocc:436:265}}\par
\par\noindent\hangindent=2.8em\hspace*{1.8em}\phantomsection\label{idx:437}正态近似（normal approximations），\hyperref[exercise:3.34]{\pageref*{exercise:3.34}}\par
\par\noindent\hangindent=2em\hspace*{1em}\phantomsection\label{idx:438}对数正态（lognormal），\hyperref[sec:3.8.3]{\pageref*{sec:3.8.3}}\par
\par\noindent\hangindent=2em\hspace*{1em}\phantomsection\label{idx:439}矩（moments），\hyperref[thm:3.4.2]{\pageref*{thm:3.4.2}}\par
\par\noindent\hangindent=2em\hspace*{1em}\phantomsection\label{idx:440}自然参数（natural parameter），\hyperref[idxocc:440:125]{\pageref*{idxocc:440:125}}\par
\par\noindent\hangindent=2em\hspace*{1em}\phantomsection\label{idx:441}和（sums），\hyperref[thm:5.2.11]{\pageref*{thm:5.2.11}}\par
\par\addvspace{2pt}
\par\noindent\hangindent=1em\hspace*{0em}\phantomsection\label{idx:442}阶乘（Factorial），\hyperref[def:1.2.16]{\pageref*{def:1.2.16}}\par
\par\addvspace{2pt}
\par\noindent\hangindent=1em\hspace*{0em}\phantomsection\label{idx:443}因子分解定理（Factorization Theorem），\hyperref[thm:6.2.6]{\pageref*{thm:6.2.6}}\par
\par\addvspace{2pt}
\par\noindent\hangindent=1em\hspace*{0em}\phantomsection\label{idx:444}信仰推断（Fiducial inference），\hyperref[idxocc:444:268]{\pageref*{idxocc:444:268}}\par
\par\addvspace{2pt}
\par\noindent\hangindent=1em\hspace*{0em}\phantomsection\label{idx:445}Fieller定理（Fieller’s Theorem），\hyperref[sec:9.5.3]{\pageref*{sec:9.5.3}}\par
\par\addvspace{2pt}
\par\noindent\hangindent=1em\hspace*{0em}\phantomsection\label{idx:446}强制二元选择（Forced binary choice），\hyperref[exercise:4.37]{\pageref*{exercise:4.37}}\par
\par\addvspace{2pt}
\par\noindent\hangindent=1em\hspace*{0em}\phantomsection\label{idx:447}泛函（Functional），\hyperref[idxocc:447:449]{\pageref*{idxocc:447:449}}\par
\par\addvspace{2pt}
\par\noindent\hangindent=1em\hspace*{0em}\phantomsection\label{idx:448}微积分基本定理（Fundamental Theorem of Calculus），\hyperref[eq:1.6.2]{\pageref*{eq:1.6.2}}, \hyperref[eq:4.1.3]{\pageref*{eq:4.1.3}}\par
\par\addvspace{2pt}
\par\noindent\hangindent=1em\hspace*{0em}\phantomsection\label{idx:449}计数基本定理（Fundamental Theorem of Counting），\hyperref[thm:1.2.14]{\pageref*{thm:1.2.14}}\par
\par\addvspace{2pt}
\par\noindent\hangindent=1em\hspace*{0em}\phantomsection\label{idx:450}伽马函数（另见\hyperref[idx:306]{分布：伽马}）（Gamma function (See also Distribution gamma)），\hyperref[eq:3.3.2]{\pageref*{eq:3.3.2}}\par
\par\addvspace{2pt}
\par\noindent\hangindent=1em\hspace*{0em}\phantomsection\label{idx:451}Gauss–Markov定理（Gauss–Markov Theorem），\hyperref[idxocc:451:474]{\pageref*{idxocc:451:474}}\par
\par\addvspace{2pt}
\par\noindent\hangindent=1em\hspace*{0em}\phantomsection\label{idx:452}广义均值（generalized means），\hyperref[exercise:7.16]{\pageref*{exercise:7.16}}\par
\par\addvspace{2pt}
\needspace{3\baselineskip}
\par\noindent\hangindent=1em\hspace*{0em}\phantomsection\label{idx:453}母函数（Generating function）\par
\par\noindent\hangindent=2em\hspace*{1em}\phantomsection\label{idx:454}累积量（cumulant），\hyperref[idxocc:454:99]{\pageref*{idxocc:454:99}}\par
\par\noindent\hangindent=2em\hspace*{1em}\phantomsection\label{idx:455}阶乘矩（factorial moment），\hyperref[idxocc:455:99]{\pageref*{idxocc:455:99}}\par
\par\noindent\hangindent=2em\hspace*{1em}\phantomsection\label{idx:456}概率（probability），\hyperref[idxocc:456:99]{\pageref*{idxocc:456:99}}\par
\par\addvspace{2pt}
\par\noindent\hangindent=1em\hspace*{0em}\phantomsection\label{idx:457}几何平均数（Geometric mean），\hyperref[ex:4.7.8]{\pageref*{ex:4.7.8}}\par
\par\addvspace{2pt}
\needspace{3\baselineskip}
\par\noindent\hangindent=1em\hspace*{0em}\phantomsection\label{idx:458}等比级数（Geometric series）\par
\par\noindent\hangindent=2em\hspace*{1em}\phantomsection\label{idx:459}求导（differentiation），\hyperref[ex:2.4.7]{\pageref*{ex:2.4.7}}\par
\par\noindent\hangindent=2em\hspace*{1em}\phantomsection\label{idx:460}部分和（partial sum），\hyperref[eq:1.5.4]{\pageref*{eq:1.5.4}}\par
\par\addvspace{2pt}
\par\noindent\hangindent=1em\hspace*{0em}\phantomsection\label{idx:461}Gosset，W. S.（Gosset, W. S.），\hyperref[idxocc:461:212]{\pageref*{idxocc:461:212}}\par
\par\addvspace{2pt}
\par\noindent\hangindent=1em\hspace*{0em}\phantomsection\label{idx:462}总均值（Grand mean），\hyperref[idxocc:462:457]{\pageref*{idxocc:462:457}}\par
\par\addvspace{2pt}
\needspace{3\baselineskip}
\par\noindent\hangindent=1em\hspace*{0em}\phantomsection\label{idx:463}群（Group）\par
\par\noindent\hangindent=2em\hspace*{1em}\phantomsection\label{idx:464}位置（location），\hyperref[ex:6.4.6]{\pageref*{ex:6.4.6}}, \hyperref[ex:7.3.6]{\pageref*{ex:7.3.6}}\par
\par\noindent\hangindent=2em\hspace*{1em}\phantomsection\label{idx:465}位置尺度（location–scale），\hyperref[exercise:6.42]{\pageref*{exercise:6.42}}\par
\par\noindent\hangindent=2em\hspace*{1em}\phantomsection\label{idx:466}变换的（of transformations），\hyperref[def:6.4.2]{\pageref*{def:6.4.2}}\par
\par\noindent\hangindent=2em\hspace*{1em}\phantomsection\label{idx:467}尺度（scale），\hyperref[exercise:6.43]{\pageref*{exercise:6.43}}\par
\par\addvspace{2pt}
\par\noindent\hangindent=1em\hspace*{0em}\phantomsection\label{idx:468}调和平均数（Harmonic mean），\hyperref[ex:4.7.8]{\pageref*{ex:4.7.8}}\par
\par\addvspace{2pt}
\par\noindent\hangindent=1em\hspace*{0em}\phantomsection\label{idx:469}风险率函数（Hazard function），\hyperref[idxocc:469:114]{\pageref*{idxocc:469:114}}, \hyperref[exercise:3.25]{\pageref*{exercise:3.25}}\par
\par\addvspace{2pt}
\par\noindent\hangindent=1em\hspace*{0em}\phantomsection\label{idx:470}分层贝叶斯估计量（Hierarchical Bayes estimator），\hyperref[idxocc:470:332]{\pageref*{idxocc:470:332}}\par
\par\addvspace{2pt}
\par\noindent\hangindent=1em\hspace*{0em}\phantomsection\label{idx:471}最高后验密度区域（Highest posterior density (HPD) region），\hyperref[cor:9.3.10]{\pageref*{cor:9.3.10}}\par
\par\addvspace{2pt}
\par\noindent\hangindent=1em\hspace*{0em}\phantomsection\label{idx:472}同方差性（Homoscedasticity），\hyperref[item:11.2.assumption.iii]{\pageref*{item:11.2.assumption.iii}}\par
\par\addvspace{2pt}
\par\noindent\hangindent=1em\hspace*{0em}\phantomsection\label{idx:473}Hwang引理（Hwang’s Lemma），\hyperref[thm:3.6.8]{\pageref*{thm:3.6.8}}\par
\par\addvspace{2pt}
\par\noindent\hangindent=1em\hspace*{0em}\phantomsection\label{idx:474}假设（Hypothesis），\hyperref[def:8.1.1]{\pageref*{def:8.1.1}}\par
\par\noindent\hangindent=2em\hspace*{1em}\phantomsection\label{idx:475}备择（alternative），\hyperref[def:8.1.2]{\pageref*{def:8.1.2}}\par
\par\noindent\hangindent=2em\hspace*{1em}\phantomsection\label{idx:476}复合（composite），\hyperref[idxocc:476:348]{\pageref*{idxocc:476:348}}\par
\par\noindent\hangindent=2em\hspace*{1em}\phantomsection\label{idx:477}零（null），\hyperref[def:8.1.2]{\pageref*{def:8.1.2}}\par
\par\noindent\hangindent=2em\hspace*{1em}\phantomsection\label{idx:478}单侧（one-sided），\hyperref[idxocc:478:348]{\pageref*{idxocc:478:348}}\par
\par\noindent\hangindent=2em\hspace*{1em}\phantomsection\label{idx:479}研究（research），\hyperref[idxocc:479:344]{\pageref*{idxocc:479:344}}\par
\par\noindent\hangindent=2em\hspace*{1em}\phantomsection\label{idx:480}简单（simple），\hyperref[idxocc:480:345]{\pageref*{idxocc:480:345}}, \hyperref[idxocc:480:348]{\pageref*{idxocc:480:348}}\par
\par\noindent\hangindent=2em\hspace*{1em}\phantomsection\label{idx:481}双侧（two-sided），\hyperref[idxocc:481:348]{\pageref*{idxocc:481:348}}\par
\par\addvspace{2pt}
\par\noindent\hangindent=1em\hspace*{0em}\phantomsection\label{idx:482}假设检验（Hypothesis test），\hyperref[def:8.1.3]{\pageref*{def:8.1.3}}\par
\par\noindent\hangindent=2em\hspace*{1em}\phantomsection\label{idx:483}滥用（abuse），\hyperref[exercise:8.16]{\pageref*{exercise:8.16}}\par
\par\noindent\hangindent=2em\hspace*{1em}\phantomsection\label{idx:484}接受域（acceptance region），\hyperref[idxocc:484:333]{\pageref*{idxocc:484:333}}\par
\par\noindent\hangindent=2em\hspace*{1em}\phantomsection\label{idx:485}渐近（asymptotic），\hyperref[exercise:10.32]{\pageref*{exercise:10.32}}\par
\par\noindent\hangindent=2.8em\hspace*{1.8em}\phantomsection\label{idx:486}伯努利（Bernoulli），\hyperref[ex:10.3.5]{\pageref*{ex:10.3.5}}\par
\par\noindent\hangindent=2em\hspace*{1em}\phantomsection\label{idx:487}贝叶斯（Bayes），\hyperref[sec:8.5.2]{\pageref*{sec:8.5.2}}\par
\sourcepageanchor{561}
\par\noindent\hangindent=2.8em\hspace*{1.8em}\phantomsection\label{idx:488}正态（normal），\hyperref[ex:8.2.7]{\pageref*{ex:8.2.7}}\par
\par\noindent\hangindent=2.8em\hspace*{1.8em}\phantomsection\label{idx:489}p值（p-value），\hyperref[exercise:8.53]{\pageref*{exercise:8.53}}\par
\par\noindent\hangindent=2em\hspace*{1em}\phantomsection\label{idx:490}生物等效性（bioequivalence），\hyperref[exercise:8.47]{\pageref*{exercise:8.47}}\par
\par\noindent\hangindent=2em\hspace*{1em}\phantomsection\label{idx:491}列联表（contingency table），\hyperref[exercise:10.31]{\pageref*{exercise:10.31}}\par
\par\noindent\hangindent=2em\hspace*{1em}\phantomsection\label{idx:492}决策论（decision theoretic），\hyperref[sec:8.3.5]{\pageref*{sec:8.3.5}}\par
\par\noindent\hangindent=2.8em\hspace*{1.8em}\phantomsection\label{idx:493}均值（mean），\hyperref[ex:8.3.31]{\pageref*{ex:8.3.31}}\par
\par\noindent\hangindent=2em\hspace*{1em}\phantomsection\label{idx:494}交并（intersection–union），\hyperref[eq:8.2.6]{\pageref*{eq:8.2.6}}\par
\par\noindent\hangindent=2.8em\hspace*{1.8em}\phantomsection\label{idx:495}伯努利（Bernoulli），\hyperref[ex:8.2.9]{\pageref*{ex:8.2.9}}, \hyperref[ex:8.3.25]{\pageref*{ex:8.3.25}}\par
\par\noindent\hangindent=2.8em\hspace*{1.8em}\phantomsection\label{idx:496}水平（level），\hyperref[thm:8.3.23]{\pageref*{thm:8.3.23}}\par
\par\noindent\hangindent=2.8em\hspace*{1.8em}\phantomsection\label{idx:497}正态（normal），\hyperref[ex:8.2.9]{\pageref*{ex:8.2.9}}, \hyperref[ex:8.3.25]{\pageref*{ex:8.3.25}}\par
\par\noindent\hangindent=2em\hspace*{1em}\phantomsection\label{idx:498}检验水平（level of test），\hyperref[def:8.3.6]{\pageref*{def:8.3.6}}\par
\par\noindent\hangindent=2em\hspace*{1em}\phantomsection\label{idx:499}McNemar检验（McNemar’s test），\hyperref[exercise:10.32]{\pageref*{exercise:10.32}}\par
\par\noindent\hangindent=2em\hspace*{1em}\phantomsection\label{idx:500}极小极大（minimax），\hyperref[exercise:8.58]{\pageref*{exercise:8.58}}\par
\par\noindent\hangindent=2em\hspace*{1em}\phantomsection\label{idx:501}功效（power），\hyperref[def:8.3.1]{\pageref*{def:8.3.1}}\par
\par\noindent\hangindent=2.8em\hspace*{1.8em}\phantomsection\label{idx:502}二项（binomial），\hyperref[ex:8.3.2]{\pageref*{ex:8.3.2}}\par
\par\noindent\hangindent=2.8em\hspace*{1.8em}\phantomsection\label{idx:503}单调（monotone），\hyperref[sec:8.5.1]{\pageref*{sec:8.5.1}}\par
\par\noindent\hangindent=2.8em\hspace*{1.8em}\phantomsection\label{idx:504}正态（normal），\hyperref[ex:8.3.3]{\pageref*{ex:8.3.3}}\par
\par\noindent\hangindent=2.8em\hspace*{1.8em}\phantomsection\label{idx:505}与风险的关系（relation to risk），\hyperref[eq:8.3.12]{\pageref*{eq:8.3.12}}\par
\par\noindent\hangindent=2em\hspace*{1em}\phantomsection\label{idx:506}拒绝域（rejection region），\hyperref[idxocc:506:333]{\pageref*{idxocc:506:333}}\par
\par\noindent\hangindent=2em\hspace*{1em}\phantomsection\label{idx:507}与置信集的关系（relation to confidence set），\hyperref[thm:9.2.2]{\pageref*{thm:9.2.2}}\par
\needspace{3\baselineskip}
\par\noindent\hangindent=2em\hspace*{1em}\phantomsection\label{idx:508}样本量（sample size）\par
\par\noindent\hangindent=2.8em\hspace*{1.8em}\phantomsection\label{idx:509}渐近（asymptotic），\hyperref[idxocc:509:427]{\pageref*{idxocc:509:427}}\par
\par\noindent\hangindent=2.8em\hspace*{1.8em}\phantomsection\label{idx:510}正态（normal），\hyperref[ex:8.3.4]{\pageref*{ex:8.3.4}}\par
\par\noindent\hangindent=2em\hspace*{1em}\phantomsection\label{idx:511}检验大小（size of test），\hyperref[def:8.3.5]{\pageref*{def:8.3.5}}\par
\par\noindent\hangindent=2em\hspace*{1em}\phantomsection\label{idx:512}充分性（sufficiency），\hyperref[cor:8.3.13]{\pageref*{cor:8.3.13}}\par
\par\noindent\hangindent=2em\hspace*{1em}\phantomsection\label{idx:513}第一类错误（type I error），\hyperref[tab:8.3.1]{\pageref*{tab:8.3.1}}, \hyperref[ex:8.3.4]{\pageref*{ex:8.3.4}}\par
\par\noindent\hangindent=2em\hspace*{1em}\phantomsection\label{idx:514}第二类错误（type II error），\hyperref[tab:8.3.1]{\pageref*{tab:8.3.1}}, \hyperref[ex:8.3.4]{\pageref*{ex:8.3.4}}\par
\par\noindent\hangindent=2em\hspace*{1em}\phantomsection\label{idx:515}一致最强（UMP），\hyperref[def:8.3.11]{\pageref*{def:8.3.11}}\par
\par\noindent\hangindent=2.8em\hspace*{1.8em}\phantomsection\label{idx:516}二项（binomial），\hyperref[ex:8.3.14]{\pageref*{ex:8.3.14}}\par
\par\noindent\hangindent=2.8em\hspace*{1.8em}\phantomsection\label{idx:517}正态（normal），\hyperref[ex:8.3.15]{\pageref*{ex:8.3.15}}, \hyperref[ex:8.3.18]{\pageref*{ex:8.3.18}}, \hyperref[exercise:8.37]{\pageref*{exercise:8.37}}\par
\par\noindent\hangindent=2.8em\hspace*{1.8em}\phantomsection\label{idx:518}与一致最准确集的关系（relation to UMA set），\hyperref[thm:9.3.5]{\pageref*{thm:9.3.5}}\par
\par\noindent\hangindent=2.8em\hspace*{1.8em}\phantomsection\label{idx:519}均匀（uniform），\hyperref[exercise:8.33]{\pageref*{exercise:8.33}}\par
\par\noindent\hangindent=2em\hspace*{1em}\phantomsection\label{idx:520}无偏（unbiased），\hyperref[def:8.3.9]{\pageref*{def:8.3.9}}\par
\par\noindent\hangindent=2.8em\hspace*{1.8em}\phantomsection\label{idx:521}与无偏集的关系（relation to unbiased set），\hyperref[idxocc:521:392]{\pageref*{idxocc:521:392}}\par
\par\noindent\hangindent=2em\hspace*{1em}\phantomsection\label{idx:522}并交（union–intersection），\hyperref[eq:8.2.4]{\pageref*{eq:8.2.4}}\par
\par\noindent\hangindent=2.8em\hspace*{1.8em}\phantomsection\label{idx:523}方差分析（ANOVA），\hyperref[thm:11.2.5]{\pageref*{thm:11.2.5}}\par
\par\noindent\hangindent=2.8em\hspace*{1.8em}\phantomsection\label{idx:524}水平（level），\hyperref[thm:8.3.21]{\pageref*{thm:8.3.21}}\par
\par\noindent\hangindent=2.8em\hspace*{1.8em}\phantomsection\label{idx:525}正态（normal），\hyperref[ex:8.2.8]{\pageref*{ex:8.2.8}}, \hyperref[ex:8.3.8]{\pageref*{ex:8.3.8}}\par
\par\addvspace{2pt}
\par\noindent\hangindent=1em\hspace*{0em}\phantomsection\label{idx:526}可识别性（Identifiability），\hyperref[def:11.2.2]{\pageref*{def:11.2.2}}\par
\par\noindent\hangindent=2em\hspace*{1em}\phantomsection\label{idx:527}方差分析（ANOVA），\hyperref[def:11.2.2]{\pageref*{def:11.2.2}}\par
\par\noindent\hangindent=2em\hspace*{1em}\phantomsection\label{idx:528}变量含误差（EIV），\hyperref[idxocc:528:501]{\pageref*{idxocc:528:501}}, \hyperref[idxocc:528:506]{\pageref*{idxocc:528:506}}\par
\par\addvspace{2pt}
\needspace{3\baselineskip}
\par\noindent\hangindent=1em\hspace*{0em}\phantomsection\label{idx:529}恒等式（Identity）\par
\par\noindent\hangindent=2em\hspace*{1em}\phantomsection\label{idx:530}贝塔（beta），\hyperref[idxocc:530:142]{\pageref*{idxocc:530:142}}\par
\par\noindent\hangindent=2em\hspace*{1em}\phantomsection\label{idx:531}二项式系数（binomial coefficients），\hyperref[exercise:1.27]{\pageref*{exercise:1.27}}, \hyperref[idxocc:531:103]{\pageref*{idxocc:531:103}}\par
\par\noindent\hangindent=2em\hspace*{1em}\phantomsection\label{idx:532}卡方（chi squared），\hyperref[thm:3.6.7]{\pageref*{thm:3.6.7}}\par
\par\noindent\hangindent=2em\hspace*{1em}\phantomsection\label{idx:533}条件协方差（conditional covariance），\hyperref[exercise:4.59]{\pageref*{exercise:4.59}}\par
\par\noindent\hangindent=2em\hspace*{1em}\phantomsection\label{idx:534}协方差（covariance），\hyperref[item:12.11.a]{\pageref*{item:12.11.a}}\par
\needspace{3\baselineskip}
\par\noindent\hangindent=2em\hspace*{1em}\phantomsection\label{idx:535}期望（expectation）\par
\par\noindent\hangindent=2.8em\hspace*{1.8em}\phantomsection\label{idx:536}连续（continuous），\hyperref[exercise:2.14]{\pageref*{exercise:2.14}}\par
\par\noindent\hangindent=2.8em\hspace*{1.8em}\phantomsection\label{idx:537}离散（discrete），\hyperref[exercise:2.14]{\pageref*{exercise:2.14}}\par
\par\noindent\hangindent=2em\hspace*{1em}\phantomsection\label{idx:538}指数族（exponential family），\hyperref[thm:3.4.2]{\pageref*{thm:3.4.2}}\par
\par\noindent\hangindent=2em\hspace*{1em}\phantomsection\label{idx:539}伽马（gamma），\hyperref[idxocc:539:142]{\pageref*{idxocc:539:142}}\par
\par\noindent\hangindent=2em\hspace*{1em}\phantomsection\label{idx:540}容斥（inclusion-exclusion），\hyperref[exercise:1.42]{\pageref*{exercise:1.42}}, \hyperref[idxocc:540:68]{\pageref*{idxocc:540:68}}\par
\par\noindent\hangindent=2em\hspace*{1em}\phantomsection\label{idx:541}分部积分（integration–by–parts），\hyperref[lem:3.6.5]{\pageref*{lem:3.6.5}}\par
\par\noindent\hangindent=2em\hspace*{1em}\phantomsection\label{idx:542}拉格朗日（Lagrange），\hyperref[exercise:5.9]{\pageref*{exercise:5.9}}\par
\par\noindent\hangindent=2em\hspace*{1em}\phantomsection\label{idx:543}负二项（negative binomial），\hyperref[eq:3.6.6]{\pageref*{eq:3.6.6}}\par
\par\noindent\hangindent=2em\hspace*{1em}\phantomsection\label{idx:544}泊松（Poisson），\hyperref[eq:3.6.5]{\pageref*{eq:3.6.5}}\par
\par\addvspace{2pt}
\par\noindent\hangindent=1em\hspace*{0em}\phantomsection\label{idx:545}独立同分布（Iid），\hyperref[def:5.1.1]{\pageref*{def:5.1.1}}\par
\par\addvspace{2pt}
\needspace{3\baselineskip}
\par\noindent\hangindent=1em\hspace*{0em}\phantomsection\label{idx:546}独立性（Independence）\par
\par\noindent\hangindent=2em\hspace*{1em}\phantomsection\label{idx:547}完备性（completeness），\hyperref[thm:6.2.24]{\pageref*{thm:6.2.24}}\par
\par\noindent\hangindent=2em\hspace*{1em}\phantomsection\label{idx:548}协方差（covariance），\hyperref[thm:4.5.5]{\pageref*{thm:4.5.5}}\par
\par\noindent\hangindent=2em\hspace*{1em}\phantomsection\label{idx:549}不交（disjoint），\hyperref[exercise:1.39]{\pageref*{exercise:1.39}}\par
\par\noindent\hangindent=2em\hspace*{1em}\phantomsection\label{idx:550}事件（events），\hyperref[def:1.3.7]{\pageref*{def:1.3.7}}\par
\par\noindent\hangindent=2.8em\hspace*{1.8em}\phantomsection\label{idx:551}相互（mutual），\hyperref[def:1.3.12]{\pageref*{def:1.3.12}}\par
\par\noindent\hangindent=2.8em\hspace*{1.8em}\phantomsection\label{idx:552}性质（properties），\hyperref[thm:1.3.9]{\pageref*{thm:1.3.9}}\par
\par\noindent\hangindent=2em\hspace*{1em}\phantomsection\label{idx:553}期望值（expected value），\hyperref[thm:4.2.10]{\pageref*{thm:4.2.10}}, \hyperref[thm:4.6.6]{\pageref*{thm:4.6.6}}\par
\par\noindent\hangindent=2em\hspace*{1em}\phantomsection\label{idx:554}随机变量（random variables），\hyperref[def:4.2.5]{\pageref*{def:4.2.5}}, \hyperref[thm:4.3.5]{\pageref*{thm:4.3.5}}, \hyperref[def:4.6.5]{\pageref*{def:4.6.5}}, \hyperref[def:5.1.1]{\pageref*{def:5.1.1}}\par
\par\noindent\hangindent=2.8em\hspace*{1.8em}\phantomsection\label{idx:555}和（sums），\hyperref[thm:4.2.12]{\pageref*{thm:4.2.12}}\par
\par\addvspace{2pt}
\par\noindent\hangindent=1em\hspace*{0em}\phantomsection\label{idx:556}示性函数（Indicator function），\hyperref[def:3.4.5]{\pageref*{def:3.4.5}}\par
\par\addvspace{2pt}
\needspace{3\baselineskip}
\par\noindent\hangindent=1em\hspace*{0em}\phantomsection\label{idx:557}不等式（Inequality）\par
\par\noindent\hangindent=2em\hspace*{1em}\phantomsection\label{idx:558}算术–几何–调和平均数（arithmetic-geometric-harmonic mean），\hyperref[sec:4.9.2]{\pageref*{sec:4.9.2}}\par
\par\noindent\hangindent=2em\hspace*{1em}\phantomsection\label{idx:559}Bonferroni与Boole的相似性（Bonferroni - Boole similarity），\hyperref[idxocc:559:41]{\pageref*{idxocc:559:41}}\par
\par\noindent\hangindent=2em\hspace*{1em}\phantomsection\label{idx:560}Bonferroni（Bonferroni’s），\hyperref[eq:1.2.9]{\pageref*{eq:1.2.9}}, \hyperref[eq:1.2.10]{\pageref*{eq:1.2.10}}, \hyperref[sec:1.8.1]{\pageref*{sec:1.8.1}}, \hyperref[idxocc:560:379]{\pageref*{idxocc:560:379}}\par
\par\noindent\hangindent=2.8em\hspace*{1.8em}\phantomsection\label{idx:561}方差分析（ANOVA），\hyperref[ex:11.2.9]{\pageref*{ex:11.2.9}}\par
\par\noindent\hangindent=2.8em\hspace*{1.8em}\phantomsection\label{idx:562}回归（regression），\hyperref[eq:11.3.42]{\pageref*{eq:11.3.42}}\par
\par\noindent\hangindent=2em\hspace*{1em}\phantomsection\label{idx:563}Boole，\hyperref[thm:1.2.11]{\pageref*{thm:1.2.11}}, \hyperref[sec:1.8.1]{\pageref*{sec:1.8.1}}\par
\par\noindent\hangindent=2em\hspace*{1em}\phantomsection\label{idx:564}Cauchy–Schwarz，\hyperref[thm:4.5.7]{\pageref*{thm:4.5.7}}, \hyperref[thm:4.7.3]{\pageref*{thm:4.7.3}}\par
\par\noindent\hangindent=2em\hspace*{1em}\phantomsection\label{idx:565}切比雪夫（Chebychev’s），\hyperref[thm:3.6.1]{\pageref*{thm:3.6.1}}, \hyperref[thm:3.8.2]{\pageref*{thm:3.8.2}}, \hyperref[thm:5.5.2]{\pageref*{thm:5.5.2}}, \hyperref[sec:5.8.3]{\pageref*{sec:5.8.3}}, \hyperref[idxocc:565:409]{\pageref*{idxocc:565:409}}\par
\par\noindent\hangindent=2em\hspace*{1em}\phantomsection\label{idx:566}协方差（covariance），\hyperref[thm:4.7.9]{\pageref*{thm:4.7.9}}\par
\par\noindent\hangindent=2em\hspace*{1em}\phantomsection\label{idx:567}高斯（Gauss），\hyperref[thm:3.8.4]{\pageref*{thm:3.8.4}}\par
\par\noindent\hangindent=2em\hspace*{1em}\phantomsection\label{idx:568}Hölder（Hölder’s），\hyperref[thm:4.7.2]{\pageref*{thm:4.7.2}}, \hyperref[idxocc:568:471]{\pageref*{idxocc:568:471}}\par
\par\noindent\hangindent=2em\hspace*{1em}\phantomsection\label{idx:569}Jensen（Jensen’s），\hyperref[thm:4.7.7]{\pageref*{thm:4.7.7}}\par
\par\noindent\hangindent=2em\hspace*{1em}\phantomsection\label{idx:570}Liapounov（Liapounov’s），\hyperref[eq:4.7.6]{\pageref*{eq:4.7.6}}\par
\par\noindent\hangindent=2em\hspace*{1em}\phantomsection\label{idx:571}对数（logarithm），\hyperref[exercise:1.28]{\pageref*{exercise:1.28}}\par
\par\noindent\hangindent=2em\hspace*{1em}\phantomsection\label{idx:572}马尔可夫（Markov’s），\hyperref[lem:3.8.3]{\pageref*{lem:3.8.3}}\par
\par\noindent\hangindent=2em\hspace*{1em}\phantomsection\label{idx:573}均值（means），\hyperref[ex:4.7.8]{\pageref*{ex:4.7.8}}\par
\par\noindent\hangindent=2em\hspace*{1em}\phantomsection\label{idx:574}矩母函数（mgf），\hyperref[exercise:3.45]{\pageref*{exercise:3.45}}, \hyperref[exercise:5.38]{\pageref*{exercise:5.38}}\par
\par\noindent\hangindent=2em\hspace*{1em}\phantomsection\label{idx:575}Minkowski（Minkowski’s），\hyperref[thm:4.7.5]{\pageref*{thm:4.7.5}}\par
\par\noindent\hangindent=2em\hspace*{1em}\phantomsection\label{idx:576}正态（normal），\hyperref[ex:3.6.3]{\pageref*{ex:3.6.3}}, \hyperref[exercise:3.47]{\pageref*{exercise:3.47}}\par
\par\noindent\hangindent=2em\hspace*{1em}\phantomsection\label{idx:577}幂平均数（power mean），\hyperref[sec:4.9.2]{\pageref*{sec:4.9.2}}\par
\par\noindent\hangindent=2em\hspace*{1em}\phantomsection\label{idx:578}三角（triangle），\hyperref[eq:4.7.8]{\pageref*{eq:4.7.8}}, \hyperref[exercise:4.64]{\pageref*{exercise:4.64}}\par
\par\noindent\hangindent=2em\hspace*{1em}\phantomsection\label{idx:579}Vysochanskiĭ–Petunin（Vysochanskiĭ-Petunin），\hyperref[thm:3.8.5]{\pageref*{thm:3.8.5}}\par
\par\addvspace{2pt}
\par\noindent\hangindent=1em\hspace*{0em}\phantomsection\label{idx:580}影响函数（Influence function），\hyperref[sec:10.6.4]{\pageref*{sec:10.6.4}}\par
\par\noindent\hangindent=2em\hspace*{1em}\phantomsection\label{idx:581}与渐近方差的关系（relation to asymptotic variance），\hyperref[idxocc:581:450]{\pageref*{idxocc:581:450}}\par
\par\addvspace{2pt}
\par\noindent\hangindent=1em\hspace*{0em}\phantomsection\label{idx:582}信息量（Information number），\hyperref[idxocc:582:305]{\pageref*{idxocc:582:305}}, \hyperref[eq:10.1.7]{\pageref*{eq:10.1.7}}\par
\par\addvspace{2pt}
\par\noindent\hangindent=1em\hspace*{0em}\phantomsection\label{idx:583}工具变量（Instrumental variable），\hyperref[sec:12.6.3]{\pageref*{sec:12.6.3}}\par
\par\addvspace{2pt}
\needspace{3\baselineskip}
\par\noindent\hangindent=1em\hspace*{0em}\phantomsection\label{idx:584}换序（Interchange）\par
\par\noindent\hangindent=2em\hspace*{1em}\phantomsection\label{idx:585}求导与求和（differentiation and summation），\hyperref[thm:2.4.8]{\pageref*{thm:2.4.8}}\par
\par\noindent\hangindent=2em\hspace*{1em}\phantomsection\label{idx:586}积分与求和（integration and summation），\hyperref[thm:2.4.10]{\pageref*{thm:2.4.10}}\par
\par\addvspace{2pt}
\needspace{3\baselineskip}
\par\noindent\hangindent=1em\hspace*{0em}\phantomsection\label{idx:587}不变性（Invariance）\par
\par\noindent\hangindent=2em\hspace*{1em}\phantomsection\label{idx:588}分布（distribution），\hyperref[def:6.4.4]{\pageref*{def:6.4.4}}\par
\par\noindent\hangindent=2em\hspace*{1em}\phantomsection\label{idx:589}形式化（formal），\hyperref[idxocc:589:272]{\pageref*{idxocc:589:272}}\par
\par\noindent\hangindent=2.8em\hspace*{1.8em}\phantomsection\label{idx:590}位置（location），\hyperref[exercise:6.41]{\pageref*{exercise:6.41}}\par
\sourcepageanchor{562}
\par\noindent\hangindent=2em\hspace*{1em}\phantomsection\label{idx:591}正态（normal），\hyperref[ex:6.4.6]{\pageref*{ex:6.4.6}}\par
\par\addvspace{2pt}
\par\noindent\hangindent=1em\hspace*{0em}\phantomsection\label{idx:592}逆二项抽样（Inverse binomial sampling），\hyperref[ex:3.2.6]{\pageref*{ex:3.2.6}}\par
\par\addvspace{2pt}
\par\noindent\hangindent=1em\hspace*{0em}\phantomsection\label{idx:593}雅可比行列式（Jacobian），\hyperref[idxocc:593:160]{\pageref*{idxocc:593:160}}\par
\par\addvspace{2pt}
\par\noindent\hangindent=1em\hspace*{0em}\phantomsection\label{idx:594}Karlin–Rubin定理（Karlin–Rubin Theorem），\hyperref[thm:8.3.17]{\pageref*{thm:8.3.17}}\par
\par\addvspace{2pt}
\par\noindent\hangindent=1em\hspace*{0em}\phantomsection\label{idx:595}核（Kernel），\hyperref[idxocc:595:83]{\pageref*{idxocc:595:83}}\par
\par\addvspace{2pt}
\par\noindent\hangindent=1em\hspace*{0em}\phantomsection\label{idx:596}柯尔莫哥洛夫公理（Kolmogorov’s Axioms），\hyperref[def:1.2.4]{\pageref*{def:1.2.4}}\par
\par\noindent\hangindent=2em\hspace*{1em}\phantomsection\label{idx:597}条件概率（conditional probability），\hyperref[idxocc:597:49]{\pageref*{idxocc:597:49}}\par
\par\noindent\hangindent=2em\hspace*{1em}\phantomsection\label{idx:598}诱导概率（induced probability），\hyperref[eq:1.4.1]{\pageref*{eq:1.4.1}}, \hyperref[idxocc:598:55]{\pageref*{idxocc:598:55}}, \hyperref[idxocc:598:70]{\pageref*{idxocc:598:70}}\par
\par\addvspace{2pt}
\par\noindent\hangindent=1em\hspace*{0em}\phantomsection\label{idx:599}峰度（Kurtosis），\hyperref[exercise:2.28]{\pageref*{exercise:2.28}}, \hyperref[exercise:7.45]{\pageref*{exercise:7.45}}\par
\par\addvspace{2pt}
\par\noindent\hangindent=1em\hspace*{0em}\phantomsection\label{idx:600}Lévy定理（Lévy Theorem），\hyperref[sec:5.8.1]{\pageref*{sec:5.8.1}}\par
\par\addvspace{2pt}
\par\noindent\hangindent=1em\hspace*{0em}\phantomsection\label{idx:601}拉普拉斯变换（Laplace transform），\hyperref[eq:2.3.8]{\pageref*{eq:2.3.8}}, \hyperref[idxocc:601:265]{\pageref*{idxocc:601:265}}\par
\par\addvspace{2pt}
\par\noindent\hangindent=1em\hspace*{0em}\phantomsection\label{idx:602}潜在变量（Latent variables），\hyperref[idxocc:602:499]{\pageref*{idxocc:602:499}}\par
\par\addvspace{2pt}
\needspace{3\baselineskip}
\par\noindent\hangindent=1em\hspace*{0em}\phantomsection\label{idx:603}大数定律（Law of Large Numbers）\par
\par\noindent\hangindent=2em\hspace*{1em}\phantomsection\label{idx:604}强（Strong），\hyperref[thm:5.5.9]{\pageref*{thm:5.5.9}}\par
\par\noindent\hangindent=2em\hspace*{1em}\phantomsection\label{idx:605}弱（Weak），\hyperref[thm:5.5.2]{\pageref*{thm:5.5.2}}\par
\par\addvspace{2pt}
\par\noindent\hangindent=1em\hspace*{0em}\phantomsection\label{idx:606}最小二乘估计（Least squares estimate），\hyperref[sec:11.3.1]{\pageref*{sec:11.3.1}}\par
\needspace{3\baselineskip}
\par\noindent\hangindent=2em\hspace*{1em}\phantomsection\label{idx:607}相合性（consistency）\par
\par\noindent\hangindent=2.8em\hspace*{1.8em}\phantomsection\label{idx:608}变量含误差（EIV），\hyperref[sec:12.6.2]{\pageref*{sec:12.6.2}}\par
\par\noindent\hangindent=2em\hspace*{1em}\phantomsection\label{idx:609}变量含误差（EIV），\hyperref[sec:12.2.2]{\pageref*{sec:12.2.2}}\par
\par\noindent\hangindent=2.8em\hspace*{1.8em}\phantomsection\label{idx:610}截距（intercept），\hyperref[idxocc:610:502]{\pageref*{idxocc:610:502}}\par
\par\noindent\hangindent=2.8em\hspace*{1.8em}\phantomsection\label{idx:611}与最大似然估计量的关系（relation to MLE），\hyperref[exercise:12.4]{\pageref*{exercise:12.4}}\par
\par\noindent\hangindent=2.8em\hspace*{1.8em}\phantomsection\label{idx:612}斜率（slope），\hyperref[eq:12.2.11]{\pageref*{eq:12.2.11}}\par
\par\noindent\hangindent=2em\hspace*{1em}\phantomsection\label{idx:613}正规方程（normal equations），\hyperref[exercise:11.25]{\pageref*{exercise:11.25}}\par
\par\noindent\hangindent=2em\hspace*{1em}\phantomsection\label{idx:614}截距的（of intercept），\hyperref[eq:11.3.9]{\pageref*{eq:11.3.9}}\par
\par\noindent\hangindent=2em\hspace*{1em}\phantomsection\label{idx:615}斜率的（of slope），\hyperref[eq:11.3.10]{\pageref*{eq:11.3.10}}\par
\par\addvspace{2pt}
\par\noindent\hangindent=1em\hspace*{0em}\phantomsection\label{idx:616}勒贝格控制收敛定理（Lebesgue’s Dominated Convergence Theorem），\hyperref[idxocc:616:88]{\pageref*{idxocc:616:88}}\par
\par\addvspace{2pt}
\par\noindent\hangindent=1em\hspace*{0em}\phantomsection\label{idx:617}Lehmann–Scheffé定理（Lehmann–Scheffé Theorem），\hyperref[thm:7.5.1]{\pageref*{thm:7.5.1}}\par
\par\addvspace{2pt}
\par\noindent\hangindent=1em\hspace*{0em}\phantomsection\label{idx:618}Leibnitz法则（Leibnitz’s Rule），\hyperref[thm:2.4.1]{\pageref*{thm:2.4.1}}\par
\par\addvspace{2pt}
\par\noindent\hangindent=1em\hspace*{0em}\phantomsection\label{idx:619}似然函数（Likelihood function），\hyperref[def:6.3.1]{\pageref*{def:6.3.1}}\par
\par\noindent\hangindent=2em\hspace*{1em}\phantomsection\label{idx:620}诱导（induced），\hyperref[eq:7.2.5]{\pageref*{eq:7.2.5}}\par
\par\noindent\hangindent=2em\hspace*{1em}\phantomsection\label{idx:621}对数似然（log likelihood），\hyperref[idxocc:621:288]{\pageref*{idxocc:621:288}}, \hyperref[exercise:7.3]{\pageref*{exercise:7.3}}\par
\par\noindent\hangindent=2em\hspace*{1em}\phantomsection\label{idx:622}负二项（negative binomial），\hyperref[ex:6.3.2]{\pageref*{ex:6.3.2}}\par
\par\noindent\hangindent=2em\hspace*{1em}\phantomsection\label{idx:623}充分性（sufficiency），\hyperref[sec:6.6.3]{\pageref*{sec:6.6.3}}\par
\par\addvspace{2pt}
\par\noindent\hangindent=1em\hspace*{0em}\phantomsection\label{idx:624}似然原则（Likelihood Principle），\hyperref[idxocc:624:268]{\pageref*{idxocc:624:268}}\par
\par\noindent\hangindent=2em\hspace*{1em}\phantomsection\label{idx:625}二项与负二项（binomial and negative binomial），\hyperref[ex:6.3.7]{\pageref*{ex:6.3.7}}\par
\par\noindent\hangindent=2em\hspace*{1em}\phantomsection\label{idx:626}Birnbaum定理（Birnbaum’s Theorem），\hyperref[thm:6.3.6]{\pageref*{thm:6.3.6}}, \hyperref[exercise:6.33]{\pageref*{exercise:6.33}}\par
\par\noindent\hangindent=2em\hspace*{1em}\phantomsection\label{idx:627}形式化（formal），\hyperref[sec:6.3.2]{\pageref*{sec:6.3.2}}, \hyperref[idxocc:627:270]{\pageref*{idxocc:627:270}}\par
\par\noindent\hangindent=2em\hspace*{1em}\phantomsection\label{idx:628}解释（interpretations），\hyperref[ex:6.3.7]{\pageref*{ex:6.3.7}}\par
\par\noindent\hangindent=2em\hspace*{1em}\phantomsection\label{idx:629}样本量（sample size），\hyperref[exercise:6.34]{\pageref*{exercise:6.34}}\par
\par\addvspace{2pt}
\par\noindent\hangindent=1em\hspace*{0em}\phantomsection\label{idx:630}似然比检验（Likelihood ratio test (LRT)），\hyperref[def:8.2.1]{\pageref*{def:8.2.1}}\par
\par\noindent\hangindent=2em\hspace*{1em}\phantomsection\label{idx:631}方差分析（ANOVA），\hyperref[exercise:11.13]{\pageref*{exercise:11.13}}\par
\par\noindent\hangindent=2em\hspace*{1em}\phantomsection\label{idx:632}渐近（asymptotic），\hyperref[thm:10.3.3]{\pageref*{thm:10.3.3}}\par
\par\noindent\hangindent=2em\hspace*{1em}\phantomsection\label{idx:633}贝塔（beta），\hyperref[exercise:8.17]{\pageref*{exercise:8.17}}\par
\par\noindent\hangindent=2em\hspace*{1em}\phantomsection\label{idx:634}指数（exponential），\hyperref[ex:8.2.3]{\pageref*{ex:8.2.3}}, \hyperref[ex:8.2.5]{\pageref*{ex:8.2.5}}, \hyperref[ex:8.3.7]{\pageref*{ex:8.3.7}}, \hyperref[exercise:8.6]{\pageref*{exercise:8.6}}\par
\par\noindent\hangindent=2em\hspace*{1em}\phantomsection\label{idx:635}最大似然估计量（MLE），\hyperref[idxocc:635:335]{\pageref*{idxocc:635:335}}\par
\par\noindent\hangindent=2em\hspace*{1em}\phantomsection\label{idx:636}多项（multinomial），\hyperref[ex:10.3.4]{\pageref*{ex:10.3.4}}\par
\par\noindent\hangindent=2em\hspace*{1em}\phantomsection\label{idx:637}Neyman–Pearson引理（Neyman–Pearson Lemma），\hyperref[exercise:8.24]{\pageref*{exercise:8.24}}\par
\par\noindent\hangindent=2em\hspace*{1em}\phantomsection\label{idx:638}正态（normal），\hyperref[ex:8.2.2]{\pageref*{ex:8.2.2}}, \hyperref[ex:8.2.6]{\pageref*{ex:8.2.6}}, \hyperref[ex:8.3.7]{\pageref*{ex:8.3.7}}, \hyperref[exercise:8.37]{\pageref*{exercise:8.37}}\par
\par\noindent\hangindent=2.8em\hspace*{1.8em}\phantomsection\label{idx:639}滋扰参数（nuisance parameter），\hyperref[ex:8.2.6]{\pageref*{ex:8.2.6}}\par
\par\noindent\hangindent=2.8em\hspace*{1.8em}\phantomsection\label{idx:640}双样本（two-sample），\hyperref[exercise:8.41]{\pageref*{exercise:8.41}}\par
\par\noindent\hangindent=2em\hspace*{1em}\phantomsection\label{idx:641}帕累托（Pareto），\hyperref[exercise:8.5]{\pageref*{exercise:8.5}}\par
\par\noindent\hangindent=2em\hspace*{1em}\phantomsection\label{idx:642}充分统计量（sufficient statistic），\hyperref[thm:8.2.4]{\pageref*{thm:8.2.4}}\par
\par\addvspace{2pt}
\par\noindent\hangindent=1em\hspace*{0em}\phantomsection\label{idx:643}Lindeberg–Feller条件（Lindeberg–Feller Condition），\hyperref[sec:5.8.1]{\pageref*{sec:5.8.1}}\par
\par\addvspace{2pt}
\par\noindent\hangindent=1em\hspace*{0em}\phantomsection\label{idx:644}线性组合（Linear combination），\hyperref[def:11.2.4]{\pageref*{def:11.2.4}}\par
\par\addvspace{2pt}
\needspace{3\baselineskip}
\par\noindent\hangindent=1em\hspace*{0em}\phantomsection\label{idx:645}位置族（Location family）\par
\par\noindent\hangindent=2em\hspace*{1em}\phantomsection\label{idx:646}枢轴量（pivotal quantity），\hyperref[ex:9.2.7]{\pageref*{ex:9.2.7}}\par
\par\addvspace{2pt}
\par\noindent\hangindent=1em\hspace*{0em}\phantomsection\label{idx:647}位置尺度族（Location–scale family），\hyperref[sec:3.5]{\pageref*{sec:3.5}}\par
\par\noindent\hangindent=2em\hspace*{1em}\phantomsection\label{idx:648}位置参数（location parameter），\hyperref[def:3.5.2]{\pageref*{def:3.5.2}}, \hyperref[def:3.5.5]{\pageref*{def:3.5.5}}\par
\par\noindent\hangindent=2em\hspace*{1em}\phantomsection\label{idx:649}枢轴量（pivotal quantity），\hyperref[ex:9.2.7]{\pageref*{ex:9.2.7}}\par
\par\noindent\hangindent=2em\hspace*{1em}\phantomsection\label{idx:650}尺度参数（scale parameter），\hyperref[def:3.5.4]{\pageref*{def:3.5.4}}\par
\par\noindent\hangindent=2em\hspace*{1em}\phantomsection\label{idx:651}标准概率密度函数（standard pdf），\hyperref[idxocc:651:126]{\pageref*{idxocc:651:126}}, \hyperref[def:3.5.5]{\pageref*{def:3.5.5}}\par
\par\noindent\hangindent=2em\hspace*{1em}\phantomsection\label{idx:652}随机序（stochastic order），\hyperref[exercise:8.34]{\pageref*{exercise:8.34}}\par
\par\addvspace{2pt}
\par\noindent\hangindent=1em\hspace*{0em}\phantomsection\label{idx:653}损失函数（Loss function），\hyperref[idxocc:653:313]{\pageref*{idxocc:653:313}}\par
\par\noindent\hangindent=2em\hspace*{1em}\phantomsection\label{idx:654}绝对误差（absolute error），\hyperref[ex:7.3.28]{\pageref*{ex:7.3.28}}\par
\par\noindent\hangindent=2em\hspace*{1em}\phantomsection\label{idx:655}广义零一（generalized zero–one），\hyperref[eq:8.3.11]{\pageref*{eq:8.3.11}}\par
\par\noindent\hangindent=2em\hspace*{1em}\phantomsection\label{idx:656}假设检验（hypothesis testing），\hyperref[eq:8.3.11]{\pageref*{eq:8.3.11}}\par
\par\noindent\hangindent=2em\hspace*{1em}\phantomsection\label{idx:657}区间估计（interval estimation），\hyperref[eq:9.3.4]{\pageref*{eq:9.3.4}}\par
\par\noindent\hangindent=2em\hspace*{1em}\phantomsection\label{idx:658}似然（likelihood），\hyperref[ex:7.3.27]{\pageref*{ex:7.3.27}}\par
\par\noindent\hangindent=2em\hspace*{1em}\phantomsection\label{idx:659}点估计（point estimation），\hyperref[ex:7.3.28]{\pageref*{ex:7.3.28}}\par
\par\noindent\hangindent=2em\hspace*{1em}\phantomsection\label{idx:660}后验期望损失（posterior expected loss），\hyperref[eq:7.3.19]{\pageref*{eq:7.3.19}}\par
\par\noindent\hangindent=2em\hspace*{1em}\phantomsection\label{idx:661}平方误差（squared error），\hyperref[ex:7.3.28]{\pageref*{ex:7.3.28}}\par
\par\noindent\hangindent=2em\hspace*{1em}\phantomsection\label{idx:662}Stein损失（Stein’s loss），\hyperref[ex:7.3.27]{\pageref*{ex:7.3.27}}\par
\par\noindent\hangindent=2em\hspace*{1em}\phantomsection\label{idx:663}零一（zero–one），\hyperref[eq:8.3.11]{\pageref*{eq:8.3.11}}\par
\par\addvspace{2pt}
\par\noindent\hangindent=1em\hspace*{0em}\phantomsection\label{idx:664}映射（Mapping），\hyperref[idxocc:664:69]{\pageref*{idxocc:664:69}}\par
\par\noindent\hangindent=2em\hspace*{1em}\phantomsection\label{idx:665}逆（inverse），\hyperref[idxocc:665:69]{\pageref*{idxocc:665:69}}\par
\par\addvspace{2pt}
\par\noindent\hangindent=1em\hspace*{0em}\phantomsection\label{idx:666}马尔可夫链蒙特卡洛（Markov chain Monte Carlo），\hyperref[sec:5.8.5]{\pageref*{sec:5.8.5}}\par
\par\addvspace{2pt}
\par\noindent\hangindent=1em\hspace*{0em}\phantomsection\label{idx:667}最大似然估计量（Maximum likelihood estimator (MLE)），\hyperref[def:7.2.4]{\pageref*{def:7.2.4}}\par
\par\noindent\hangindent=2em\hspace*{1em}\phantomsection\label{idx:668}渐近效率（asymptotic efficiency），\hyperref[thm:10.1.12]{\pageref*{thm:10.1.12}}\par
\par\noindent\hangindent=2em\hspace*{1em}\phantomsection\label{idx:669}伯努利（Bernoulli），\hyperref[ex:7.2.7]{\pageref*{ex:7.2.7}}\par
\par\noindent\hangindent=2.8em\hspace*{1.8em}\phantomsection\label{idx:670}近似方差（approximate variance），\hyperref[ex:10.1.14]{\pageref*{ex:10.1.14}}\par
\par\noindent\hangindent=2em\hspace*{1em}\phantomsection\label{idx:671}二项（binomial），\hyperref[ex:7.2.9]{\pageref*{ex:7.2.9}}, \hyperref[ex:7.2.13]{\pageref*{ex:7.2.13}}\par
\par\noindent\hangindent=2em\hspace*{1em}\phantomsection\label{idx:672}二元正态（bivariate normal），\hyperref[exercise:7.18]{\pageref*{exercise:7.18}}\par
\par\noindent\hangindent=2em\hspace*{1em}\phantomsection\label{idx:673}微积分（calculus），\hyperref[ex:7.2.12]{\pageref*{ex:7.2.12}}\par
\par\noindent\hangindent=2em\hspace*{1em}\phantomsection\label{idx:674}相合性（consistency），\hyperref[thm:10.1.6]{\pageref*{thm:10.1.6}}\par
\par\noindent\hangindent=2.8em\hspace*{1.8em}\phantomsection\label{idx:675}变量含误差（EIV），\hyperref[idxocc:675:505]{\pageref*{idxocc:675:505}}, \hyperref[idxocc:675:507]{\pageref*{idxocc:675:507}}\par
\par\noindent\hangindent=2em\hspace*{1em}\phantomsection\label{idx:676}不稳定性（instability），\hyperref[ex:7.2.13]{\pageref*{ex:7.2.13}}\par
\par\noindent\hangindent=2em\hspace*{1em}\phantomsection\label{idx:677}不变性（invariance），\hyperref[thm:7.2.10]{\pageref*{thm:7.2.10}}\par
\par\noindent\hangindent=2em\hspace*{1em}\phantomsection\label{idx:678}似然比检验（LRT），\hyperref[idxocc:678:335]{\pageref*{idxocc:678:335}}\par
\par\noindent\hangindent=2em\hspace*{1em}\phantomsection\label{idx:679}正态（normal），\hyperref[ex:7.2.5]{\pageref*{ex:7.2.5}}, \hyperref[ex:7.2.11]{\pageref*{ex:7.2.11}}\par
\par\noindent\hangindent=2em\hspace*{1em}\phantomsection\label{idx:680}与矩估计法的关系（relation to method of moments），\hyperref[sec:7.5.1]{\pageref*{sec:7.5.1}}\par
\par\addvspace{2pt}
\par\noindent\hangindent=1em\hspace*{0em}\phantomsection\label{idx:681}二次函数的最大值（Maximum of quadratic functions），\hyperref[lem:11.2.7]{\pageref*{lem:11.2.7}}, \hyperref[eq:11.3.18]{\pageref*{eq:11.3.18}}, \hyperref[thm:11.3.6]{\pageref*{thm:11.3.6}}, \hyperref[exercise:11.14]{\pageref*{exercise:11.14}}, \hyperref[item:11.40.a]{\pageref*{item:11.40.a}}\par
\par\addvspace{2pt}
\par\noindent\hangindent=1em\hspace*{0em}\phantomsection\label{idx:682}最大概率估计量（Maximum probability estimator），\hyperref[item:11.37.e]{\pageref*{item:11.37.e}}\par
\par\addvspace{2pt}
\par\noindent\hangindent=1em\hspace*{0em}\phantomsection\label{idx:683}均方误差（Mean squared error），\hyperref[def:7.3.1]{\pageref*{def:7.3.1}}, \hyperref[ex:7.3.28]{\pageref*{ex:7.3.28}}\par
\par\addvspace{2pt}
\par\noindent\hangindent=1em\hspace*{0em}\phantomsection\label{idx:684}中位数（Median），\hyperref[exercise:2.17]{\pageref*{exercise:2.17}}\par
\par\noindent\hangindent=2em\hspace*{1em}\phantomsection\label{idx:685}最小化性质（minimizing property），\hyperref[exercise:2.18]{\pageref*{exercise:2.18}}\par
\par\addvspace{2pt}
\par\noindent\hangindent=1em\hspace*{0em}\phantomsection\label{idx:686}德梅雷骰子问题（Meré dice problem），\hyperref[ex:1.3.8]{\pageref*{ex:1.3.8}}, \hyperref[ex:3.2.3]{\pageref*{ex:3.2.3}}\par
\par\addvspace{2pt}
\par\noindent\hangindent=1em\hspace*{0em}\phantomsection\label{idx:687}度量（Metric），\hyperref[item:12.3.b]{\pageref*{item:12.3.b}}\par
\par\addvspace{2pt}
\par\noindent\hangindent=1em\hspace*{0em}\phantomsection\label{idx:688}Metropolis算法（Metropolis Algorithm），\hyperref[idxocc:688:238]{\pageref*{idxocc:688:238}}\par
\sourcepageanchor{563}
\par\addvspace{2pt}
\par\noindent\hangindent=1em\hspace*{0em}\phantomsection\label{idx:689}矩母函数（Moment generating function (mgf)），\hyperref[def:2.3.6]{\pageref*{def:2.3.6}}, \hyperref[sec:2.6.3]{\pageref*{sec:2.6.3}}\par
\par\noindent\hangindent=2em\hspace*{1em}\phantomsection\label{idx:690}收敛（convergence），\hyperref[thm:2.3.12]{\pageref*{thm:2.3.12}}\par
\par\noindent\hangindent=2em\hspace*{1em}\phantomsection\label{idx:691}同分布（identically distributed），\hyperref[thm:2.3.11]{\pageref*{thm:2.3.11}}\par
\par\noindent\hangindent=2em\hspace*{1em}\phantomsection\label{idx:692}独立随机变量（independent random variables），\hyperref[thm:4.2.12]{\pageref*{thm:4.2.12}}\par
\par\noindent\hangindent=2em\hspace*{1em}\phantomsection\label{idx:693}性质（properties），\hyperref[thm:2.3.15]{\pageref*{thm:2.3.15}}\par
\par\noindent\hangindent=2em\hspace*{1em}\phantomsection\label{idx:694}样本均值（sample mean），\hyperref[thm:5.2.7]{\pageref*{thm:5.2.7}}\par
\par\noindent\hangindent=2em\hspace*{1em}\phantomsection\label{idx:695}和（sum），\hyperref[thm:4.6.7]{\pageref*{thm:4.6.7}}\par
\par\noindent\hangindent=2em\hspace*{1em}\phantomsection\label{idx:696}在中心极限定理中的用途（use in CLT），\hyperref[idxocc:696:224]{\pageref*{idxocc:696:224}}\par
\par\addvspace{2pt}
\par\noindent\hangindent=1em\hspace*{0em}\phantomsection\label{idx:697}矩（Moments），\hyperref[def:2.3.1]{\pageref*{def:2.3.1}}\par
\par\noindent\hangindent=2em\hspace*{1em}\phantomsection\label{idx:698}非唯一性（nonuniqueness），\hyperref[ex:2.3.10]{\pageref*{ex:2.3.10}}, \hyperref[exercise:2.35]{\pageref*{exercise:2.35}}\par
\par\noindent\hangindent=2em\hspace*{1em}\phantomsection\label{idx:699}唯一性（uniqueness），\hyperref[sec:2.6.1]{\pageref*{sec:2.6.1}}\par
\par\addvspace{2pt}
\par\noindent\hangindent=1em\hspace*{0em}\phantomsection\label{idx:700}单调似然比（Monotone likelihood ratio (MLR)），\hyperref[def:8.3.16]{\pageref*{def:8.3.16}}\par
\par\noindent\hangindent=2em\hspace*{1em}\phantomsection\label{idx:701}二项（binomial），\hyperref[exercise:8.25]{\pageref*{exercise:8.25}}\par
\par\noindent\hangindent=2em\hspace*{1em}\phantomsection\label{idx:702}柯西（Cauchy），\hyperref[exercise:8.29]{\pageref*{exercise:8.29}}\par
\par\noindent\hangindent=2em\hspace*{1em}\phantomsection\label{idx:703}指数族（exponential family），\hyperref[exercise:8.27]{\pageref*{exercise:8.27}}\par
\par\noindent\hangindent=2em\hspace*{1em}\phantomsection\label{idx:704}逻辑（logistic），\hyperref[exercise:8.28]{\pageref*{exercise:8.28}}\par
\par\noindent\hangindent=2em\hspace*{1em}\phantomsection\label{idx:705}非中心t（noncentral t），\hyperref[exercise:8.35]{\pageref*{exercise:8.35}}\par
\par\noindent\hangindent=2em\hspace*{1em}\phantomsection\label{idx:706}正态（normal），\hyperref[exercise:8.25]{\pageref*{exercise:8.25}}\par
\par\noindent\hangindent=2em\hspace*{1em}\phantomsection\label{idx:707}泊松（Poisson），\hyperref[exercise:8.25]{\pageref*{exercise:8.25}}\par
\par\noindent\hangindent=2em\hspace*{1em}\phantomsection\label{idx:708}随机序（stochastic order），\hyperref[exercise:8.26]{\pageref*{exercise:8.26}}\par
\par\addvspace{2pt}
\needspace{3\baselineskip}
\par\noindent\hangindent=1em\hspace*{0em}\phantomsection\label{idx:709}蒙特卡洛（Monte Carlo）\par
\par\noindent\hangindent=2em\hspace*{1em}\phantomsection\label{idx:710}巧算（swindle），\hyperref[exercise:6.31]{\pageref*{exercise:6.31}}\par
\par\addvspace{2pt}
\par\noindent\hangindent=1em\hspace*{0em}\phantomsection\label{idx:711}蒙提霍尔问题（Monte Hall problem），\hyperref[exercise:1.37]{\pageref*{exercise:1.37}}\par
\par\addvspace{2pt}
\par\noindent\hangindent=1em\hspace*{0em}\phantomsection\label{idx:712}多项式系数（Multinomial coefficient），\hyperref[idxocc:712:179]{\pageref*{idxocc:712:179}}\par
\par\addvspace{2pt}
\par\noindent\hangindent=1em\hspace*{0em}\phantomsection\label{idx:713}多项式定理（Multinomial Theorem），\hyperref[thm:4.6.4]{\pageref*{thm:4.6.4}}\par
\par\addvspace{2pt}
\par\noindent\hangindent=1em\hspace*{0em}\phantomsection\label{idx:714}多重比较（Multiple comparisons），\hyperref[sec:11.5.2]{\pageref*{sec:11.5.2}}\par
\par\noindent\hangindent=2em\hspace*{1em}\phantomsection\label{idx:715}逐次比较错误率（comparisonwise error rate），\hyperref[item:11.17.b]{\pageref*{item:11.17.b}}\par
\par\noindent\hangindent=2em\hspace*{1em}\phantomsection\label{idx:716}整体试验错误率（experimentwise error rate），\hyperref[item:11.17.b]{\pageref*{item:11.17.b}}\par
\par\noindent\hangindent=2em\hspace*{1em}\phantomsection\label{idx:717}受保护的最小显著差（protected LSD），\hyperref[exercise:11.17]{\pageref*{exercise:11.17}}\par
\par\noindent\hangindent=2em\hspace*{1em}\phantomsection\label{idx:718}Tukey的Q方法（Tukey’s Q method），\hyperref[idxocc:718:494]{\pageref*{idxocc:718:494}}\par
\par\addvspace{2pt}
\par\noindent\hangindent=1em\hspace*{0em}\phantomsection\label{idx:719}必要统计量（Necessary statistic），\hyperref[def:6.6.3]{\pageref*{def:6.6.3}}\par
\par\addvspace{2pt}
\par\noindent\hangindent=1em\hspace*{0em}\phantomsection\label{idx:720}邻域（Neighborhood），\hyperref[def:2.3.6]{\pageref*{def:2.3.6}}\par
\par\addvspace{2pt}
\par\noindent\hangindent=1em\hspace*{0em}\phantomsection\label{idx:721}Neyman最短（Neyman shortest），\hyperref[idxocc:721:392]{\pageref*{idxocc:721:392}}\par
\par\addvspace{2pt}
\par\noindent\hangindent=1em\hspace*{0em}\phantomsection\label{idx:722}Neyman–Pearson引理（Neyman–Pearson Lemma），\hyperref[thm:8.3.12]{\pageref*{thm:8.3.12}}\par
\par\noindent\hangindent=2em\hspace*{1em}\phantomsection\label{idx:723}似然比检验（LRT），\hyperref[exercise:8.24]{\pageref*{exercise:8.24}}\par
\par\addvspace{2pt}
\needspace{3\baselineskip}
\par\noindent\hangindent=1em\hspace*{0em}\phantomsection\label{idx:724}尼罗河（Nile）\par
\par\noindent\hangindent=2em\hspace*{1em}\phantomsection\label{idx:725}问题（problem），\hyperref[exercise:6.37]{\pageref*{exercise:6.37}}\par
\par\addvspace{2pt}
\par\noindent\hangindent=1em\hspace*{0em}\phantomsection\label{idx:726}正规方程（Normal equations），\hyperref[exercise:11.25]{\pageref*{exercise:11.25}}\par
\par\addvspace{2pt}
\par\noindent\hangindent=1em\hspace*{0em}\phantomsection\label{idx:727}滋扰参数（Nuisance parameter），\hyperref[idxocc:727:337]{\pageref*{idxocc:727:337}}\par
\par\addvspace{2pt}
\par\noindent\hangindent=1em\hspace*{0em}\phantomsection\label{idx:728}优势（Odds），\hyperref[idxocc:728:414]{\pageref*{idxocc:728:414}}\par
\par\addvspace{2pt}
\par\noindent\hangindent=1em\hspace*{0em}\phantomsection\label{idx:729}一一变换（One-to-one transformation），\hyperref[idxocc:729:71]{\pageref*{idxocc:729:71}}, \hyperref[idxocc:729:160]{\pageref*{idxocc:729:160}}\par
\par\addvspace{2pt}
\par\noindent\hangindent=1em\hspace*{0em}\phantomsection\label{idx:730}满射变换（Onto transformation），\hyperref[idxocc:730:71]{\pageref*{idxocc:730:71}}, \hyperref[idxocc:730:160]{\pageref*{idxocc:730:160}}\par
\par\addvspace{2pt}
\par\noindent\hangindent=1em\hspace*{0em}\phantomsection\label{idx:731}次序统计量（Order statistics），\hyperref[def:5.4.1]{\pageref*{def:5.4.1}}\par
\par\noindent\hangindent=2em\hspace*{1em}\phantomsection\label{idx:732}累积分布函数（cdf），\hyperref[thm:5.4.3]{\pageref*{thm:5.4.3}}\par
\par\noindent\hangindent=2em\hspace*{1em}\phantomsection\label{idx:733}联合累积分布函数（joint cdf），\hyperref[exercise:5.26]{\pageref*{exercise:5.26}}\par
\par\noindent\hangindent=2em\hspace*{1em}\phantomsection\label{idx:734}联合概率密度函数（joint pdf），\hyperref[thm:5.4.6]{\pageref*{thm:5.4.6}}\par
\par\noindent\hangindent=2em\hspace*{1em}\phantomsection\label{idx:735}概率密度函数（pdf），\hyperref[thm:5.4.4]{\pageref*{thm:5.4.4}}\par
\par\noindent\hangindent=2em\hspace*{1em}\phantomsection\label{idx:736}概率质量函数（pmf），\hyperref[thm:5.4.3]{\pageref*{thm:5.4.3}}\par
\par\noindent\hangindent=2em\hspace*{1em}\phantomsection\label{idx:737}间距（spacings），\hyperref[idxocc:737:183]{\pageref*{idxocc:737:183}}\par
\par\noindent\hangindent=2em\hspace*{1em}\phantomsection\label{idx:738}充分（sufficient），\hyperref[exercise:6.8]{\pageref*{exercise:6.8}}\par
\par\addvspace{2pt}
\par\noindent\hangindent=1em\hspace*{0em}\phantomsection\label{idx:739}正交最小二乘（Orthogonal least squares），\hyperref[sec:12.2.2]{\pageref*{sec:12.2.2}}\par
\par\addvspace{2pt}
\needspace{3\baselineskip}
\par\noindent\hangindent=1em\hspace*{0em}\phantomsection\label{idx:740}过参数化模型（Overparameterized model）\par
\par\noindent\hangindent=2em\hspace*{1em}\phantomsection\label{idx:741}方差分析（ANOVA），\hyperref[eq:11.2.2]{\pageref*{eq:11.2.2}}\par
\par\addvspace{2pt}
\par\noindent\hangindent=1em\hspace*{0em}\phantomsection\label{idx:742}p值（P-value），\hyperref[def:8.3.26]{\pageref*{def:8.3.26}}\par
\par\noindent\hangindent=2em\hspace*{1em}\phantomsection\label{idx:743}条件（conditional），\hyperref[eq:8.3.10]{\pageref*{eq:8.3.10}}\par
\par\noindent\hangindent=2em\hspace*{1em}\phantomsection\label{idx:744}与贝叶斯检验的关系（relation to Bayes test），\hyperref[sec:8.5.3]{\pageref*{sec:8.5.3}}\par
\par\addvspace{2pt}
\par\noindent\hangindent=1em\hspace*{0em}\phantomsection\label{idx:745}两两不交集（Pairwise disjoint sets），\hyperref[def:1.1.5]{\pageref*{def:1.1.5}}\par
\par\addvspace{2pt}
\needspace{3\baselineskip}
\par\noindent\hangindent=1em\hspace*{0em}\phantomsection\label{idx:746}悖论（Paradox）\par
\par\noindent\hangindent=2em\hspace*{1em}\phantomsection\label{idx:747}博雷尔（Borel），\hyperref[sec:4.9.3]{\pageref*{sec:4.9.3}}\par
\par\noindent\hangindent=2em\hspace*{1em}\phantomsection\label{idx:748}交换（exchange），\hyperref[sec:4.9.1]{\pageref*{sec:4.9.1}}\par
\par\addvspace{2pt}
\par\noindent\hangindent=1em\hspace*{0em}\phantomsection\label{idx:749}并联系统（Parallel system），\hyperref[exercise:4.55]{\pageref*{exercise:4.55}}\par
\par\addvspace{2pt}
\par\noindent\hangindent=1em\hspace*{0em}\phantomsection\label{idx:750}参数（Parameter），\hyperref[ex:2.1.1]{\pageref*{ex:2.1.1}}\par
\par\addvspace{2pt}
\par\noindent\hangindent=1em\hspace*{0em}\phantomsection\label{idx:751}参数空间（Parameter space），\hyperref[idxocc:751:313]{\pageref*{idxocc:751:313}}\par
\par\addvspace{2pt}
\par\noindent\hangindent=1em\hspace*{0em}\phantomsection\label{idx:752}参数自助法（Parametric bootstrap），\hyperref[ex:5.6.6]{\pageref*{ex:5.6.6}}\par
\par\addvspace{2pt}
\par\noindent\hangindent=1em\hspace*{0em}\phantomsection\label{idx:753}划分（Partition），\hyperref[def:1.1.6]{\pageref*{def:1.1.6}}, \hyperref[thm:1.3.5]{\pageref*{thm:1.3.5}}, \hyperref[idxocc:753:252]{\pageref*{idxocc:753:252}}, \hyperref[idxocc:753:259]{\pageref*{idxocc:753:259}}\par
\par\addvspace{2pt}
\par\noindent\hangindent=1em\hspace*{0em}\phantomsection\label{idx:754}百分位数（Percentile），\hyperref[idxocc:754:216]{\pageref*{idxocc:754:216}}\par
\par\addvspace{2pt}
\par\noindent\hangindent=1em\hspace*{0em}\phantomsection\label{idx:755}Pitman估计量（见\hyperref[idx:787]{点估计量：不变}）（Pitman Estimator (see Point estimator, invariant)），\hyperref[exercise:7.35]{\pageref*{exercise:7.35}}\par
\par\addvspace{2pt}
\par\noindent\hangindent=1em\hspace*{0em}\phantomsection\label{idx:756}枢轴量（Pivotal quantity），\hyperref[def:9.2.6]{\pageref*{def:9.2.6}}\par
\par\noindent\hangindent=2em\hspace*{1em}\phantomsection\label{idx:757}伽马（gamma），\hyperref[ex:9.2.8]{\pageref*{ex:9.2.8}}\par
\par\noindent\hangindent=2em\hspace*{1em}\phantomsection\label{idx:758}一般形式（general form），\hyperref[eq:9.2.11]{\pageref*{eq:9.2.11}}\par
\par\noindent\hangindent=2em\hspace*{1em}\phantomsection\label{idx:759}位置族（location family），\hyperref[ex:9.2.7]{\pageref*{ex:9.2.7}}\par
\par\noindent\hangindent=2em\hspace*{1em}\phantomsection\label{idx:760}位置尺度族（location–scale family），\hyperref[ex:9.2.7]{\pageref*{ex:9.2.7}}\par
\par\noindent\hangindent=2em\hspace*{1em}\phantomsection\label{idx:761}尺度族（scale family），\hyperref[ex:9.2.7]{\pageref*{ex:9.2.7}}\par
\par\addvspace{2pt}
\par\noindent\hangindent=1em\hspace*{0em}\phantomsection\label{idx:762}点估计值（Point estimate），\hyperref[idxocc:762:283]{\pageref*{idxocc:762:283}}\par
\par\addvspace{2pt}
\par\noindent\hangindent=1em\hspace*{0em}\phantomsection\label{idx:763}点估计量（Point estimator），\hyperref[def:7.1.1]{\pageref*{def:7.1.1}}\par
\par\noindent\hangindent=2em\hspace*{1em}\phantomsection\label{idx:764}近似（approximation），\hyperref[eq:5.5.9]{\pageref*{eq:5.5.9}}\par
\needspace{3\baselineskip}
\par\noindent\hangindent=2em\hspace*{1em}\phantomsection\label{idx:765}贝叶斯（Bayes）\par
\par\noindent\hangindent=2.8em\hspace*{1.8em}\phantomsection\label{idx:766}经验（empirical），\hyperref[idxocc:766:331]{\pageref*{idxocc:766:331}}\par
\par\noindent\hangindent=2.8em\hspace*{1.8em}\phantomsection\label{idx:767}分层（hierarchical），\hyperref[idxocc:767:332]{\pageref*{idxocc:767:332}}\par
\par\noindent\hangindent=2.8em\hspace*{1.8em}\phantomsection\label{idx:768}稳健（robust），\hyperref[idxocc:768:331]{\pageref*{idxocc:768:331}}\par
\par\noindent\hangindent=2em\hspace*{1em}\phantomsection\label{idx:769}崩溃值（breakdown value），\hyperref[def:10.2.2]{\pageref*{def:10.2.2}}\par
\par\noindent\hangindent=2em\hspace*{1em}\phantomsection\label{idx:770}相合（consistent），\hyperref[def:10.1.1]{\pageref*{def:10.1.1}}\par
\par\noindent\hangindent=2.8em\hspace*{1.8em}\phantomsection\label{idx:771}最大似然估计量（maximum likelihood estimator），\hyperref[thm:10.1.6]{\pageref*{thm:10.1.6}}\par
\par\noindent\hangindent=2.8em\hspace*{1.8em}\phantomsection\label{idx:772}正态（normal），\hyperref[ex:10.1.2]{\pageref*{ex:10.1.2}}\par
\par\noindent\hangindent=2em\hspace*{1em}\phantomsection\label{idx:773}相合的自助法（consistent bootstrap），\hyperref[item:10.1.bootstrap.b]{\pageref*{item:10.1.bootstrap.b}}\par
\needspace{3\baselineskip}
\par\noindent\hangindent=2em\hspace*{1em}\phantomsection\label{idx:774}相关（correlation）\par
\par\noindent\hangindent=2.8em\hspace*{1.8em}\phantomsection\label{idx:775}极限分布（limit distribution），\hyperref[item:10.17.d]{\pageref*{item:10.17.d}}\par
\par\noindent\hangindent=2em\hspace*{1em}\phantomsection\label{idx:776}决策论（decision theoretic），\hyperref[ex:7.3.28]{\pageref*{ex:7.3.28}}\par
\par\noindent\hangindent=2.8em\hspace*{1.8em}\phantomsection\label{idx:777}二项（binomial），\hyperref[ex:7.3.30]{\pageref*{ex:7.3.30}}\par
\par\noindent\hangindent=2.8em\hspace*{1.8em}\phantomsection\label{idx:778}方差（variance），\hyperref[ex:7.3.26]{\pageref*{ex:7.3.26}}\par
\par\noindent\hangindent=2em\hspace*{1em}\phantomsection\label{idx:779}效率（efficiency），\hyperref[sec:10.1.2]{\pageref*{sec:10.1.2}}, \hyperref[def:10.1.11]{\pageref*{def:10.1.11}}\par
\par\noindent\hangindent=2em\hspace*{1em}\phantomsection\label{idx:780}有效（efficient），\hyperref[idxocc:780:415]{\pageref*{idxocc:780:415}}\par
\needspace{3\baselineskip}
\par\noindent\hangindent=2em\hspace*{1em}\phantomsection\label{idx:781}经验贝叶斯（empirical Bayes）\par
\par\noindent\hangindent=2.8em\hspace*{1.8em}\phantomsection\label{idx:782}正态（normal），\hyperref[exercise:7.25]{\pageref*{exercise:7.25}}\par
\needspace{3\baselineskip}
\par\noindent\hangindent=2em\hspace*{1em}\phantomsection\label{idx:783}等变（equivariant）\par
\par\noindent\hangindent=2.8em\hspace*{1.8em}\phantomsection\label{idx:784}位置（location），\hyperref[exercise:6.41]{\pageref*{exercise:6.41}}\par
\par\noindent\hangindent=2.8em\hspace*{1.8em}\phantomsection\label{idx:785}位置尺度（location–scale），\hyperref[exercise:6.42]{\pageref*{exercise:6.42}}\par
\par\noindent\hangindent=2em\hspace*{1em}\phantomsection\label{idx:786}广义均值（generalized means），\hyperref[exercise:7.16]{\pageref*{exercise:7.16}}\par
\needspace{3\baselineskip}
\par\noindent\hangindent=2em\hspace*{1em}\phantomsection\label{idx:787}不变（invariant）\par
\par\noindent\hangindent=2.8em\hspace*{1.8em}\phantomsection\label{idx:788}形式化（formal），\hyperref[idxocc:788:301]{\pageref*{idxocc:788:301}}\par
\sourcepageanchor{564}
\par\noindent\hangindent=2.8em\hspace*{1.8em}\phantomsection\label{idx:789}均方误差（mean squared error），\hyperref[ex:7.3.6]{\pageref*{ex:7.3.6}}\par
\par\noindent\hangindent=2.8em\hspace*{1.8em}\phantomsection\label{idx:790}Pitman，\hyperref[exercise:7.35]{\pageref*{exercise:7.35}}\par
\par\noindent\hangindent=2em\hspace*{1em}\phantomsection\label{idx:791}刀切法（jackknife），\hyperref[exercise:7.66]{\pageref*{exercise:7.66}}\par
\par\noindent\hangindent=2em\hspace*{1em}\phantomsection\label{idx:792}线性（linear），\hyperref[exercise:7.41]{\pageref*{exercise:7.41}}, \hyperref[eq:11.3.15]{\pageref*{eq:11.3.15}}\par
\needspace{3\baselineskip}
\par\noindent\hangindent=2em\hspace*{1em}\phantomsection\label{idx:793}M估计量（M-estimator）\par
\par\noindent\hangindent=2.8em\hspace*{1.8em}\phantomsection\label{idx:794}渐近有效（asymptotically efficient），\hyperref[exercise:10.29]{\pageref*{exercise:10.29}}\par
\needspace{3\baselineskip}
\par\noindent\hangindent=2em\hspace*{1em}\phantomsection\label{idx:795}最大似然估计量（maximum likelihood estimator）\par
\par\noindent\hangindent=2.8em\hspace*{1.8em}\phantomsection\label{idx:796}渐近效率（asymptotic efficiency），\hyperref[thm:10.1.12]{\pageref*{thm:10.1.12}}\par
\needspace{3\baselineskip}
\par\noindent\hangindent=2em\hspace*{1em}\phantomsection\label{idx:797}均值（mean）\par
\par\noindent\hangindent=2.8em\hspace*{1.8em}\phantomsection\label{idx:798}影响函数（influence function），\hyperref[ex:10.6.2]{\pageref*{ex:10.6.2}}\par
\needspace{3\baselineskip}
\par\noindent\hangindent=2em\hspace*{1em}\phantomsection\label{idx:799}均方误差（mean squared error）\par
\par\noindent\hangindent=2.8em\hspace*{1.8em}\phantomsection\label{idx:800}伯努利（Bernoulli），\hyperref[ex:7.3.5]{\pageref*{ex:7.3.5}}\par
\par\noindent\hangindent=2.8em\hspace*{1.8em}\phantomsection\label{idx:801}方差（variance），\hyperref[eq:7.3.1]{\pageref*{eq:7.3.1}}\par
\needspace{3\baselineskip}
\par\noindent\hangindent=2em\hspace*{1em}\phantomsection\label{idx:802}中位数（median）\par
\par\noindent\hangindent=2.8em\hspace*{1.8em}\phantomsection\label{idx:803}影响函数（influence function），\hyperref[idxocc:803:450]{\pageref*{idxocc:803:450}}\par
\par\noindent\hangindent=2.8em\hspace*{1.8em}\phantomsection\label{idx:804}极限分布（limit distribution），\hyperref[ex:10.2.3]{\pageref*{ex:10.2.3}}\par
\par\noindent\hangindent=2em\hspace*{1em}\phantomsection\label{idx:805}中位数绝对偏差（median absolute deviation），\hyperref[exercise:10.21]{\pageref*{exercise:10.21}}\par
\par\noindent\hangindent=2em\hspace*{1em}\phantomsection\label{idx:806}矩估计法（method of moments），\hyperref[sec:7.2.1]{\pageref*{sec:7.2.1}}\par
\par\noindent\hangindent=2.8em\hspace*{1.8em}\phantomsection\label{idx:807}二项（binomial），\hyperref[ex:7.2.2]{\pageref*{ex:7.2.2}}\par
\par\noindent\hangindent=2.8em\hspace*{1.8em}\phantomsection\label{idx:808}卡方（chi squared），\hyperref[ex:7.2.3]{\pageref*{ex:7.2.3}}\par
\par\noindent\hangindent=2.8em\hspace*{1.8em}\phantomsection\label{idx:809}正态（normal），\hyperref[ex:7.2.1]{\pageref*{ex:7.2.1}}\par
\par\noindent\hangindent=2.8em\hspace*{1.8em}\phantomsection\label{idx:810}与最大似然估计量的关系（relation to MLE），\hyperref[sec:7.5.1]{\pageref*{sec:7.5.1}}\par
\par\noindent\hangindent=2em\hspace*{1em}\phantomsection\label{idx:811}幂平均数（power means），\hyperref[exercise:7.16]{\pageref*{exercise:7.16}}\par
\par\noindent\hangindent=2em\hspace*{1em}\phantomsection\label{idx:812}比值（ratio），\hyperref[ex:5.5.27]{\pageref*{ex:5.5.27}}\par
\par\noindent\hangindent=2em\hspace*{1em}\phantomsection\label{idx:813}稳健（robust），\hyperref[idxocc:813:420]{\pageref*{idxocc:813:420}}\par
\par\noindent\hangindent=2.8em\hspace*{1.8em}\phantomsection\label{idx:814}Huber估计量（Huber estimator），\hyperref[ex:10.2.5]{\pageref*{ex:10.2.5}}\par
\par\noindent\hangindent=2.8em\hspace*{1.8em}\phantomsection\label{idx:815}极限分布（limit distribution），\hyperref[eq:10.2.6]{\pageref*{eq:10.2.6}}\par
\par\noindent\hangindent=2.8em\hspace*{1.8em}\phantomsection\label{idx:816}M估计量（M-estimator），\hyperref[sec:10.2.2]{\pageref*{sec:10.2.2}}\par
\par\noindent\hangindent=2em\hspace*{1em}\phantomsection\label{idx:817}充分性（sufficiency），\hyperref[thm:7.3.17]{\pageref*{thm:7.3.17}}\par
\par\addvspace{2pt}
\par\noindent\hangindent=1em\hspace*{0em}\phantomsection\label{idx:818}泊松基本假设（Poisson postulates），\hyperref[thm:3.8.1]{\pageref*{thm:3.8.1}}\par
\par\addvspace{2pt}
\par\noindent\hangindent=1em\hspace*{0em}\phantomsection\label{idx:819}后验分布（Posterior distribution），\hyperref[idxocc:819:293]{\pageref*{idxocc:819:293}}\par
\par\addvspace{2pt}
\par\noindent\hangindent=1em\hspace*{0em}\phantomsection\label{idx:820}后验期望损失（另见\hyperref[idx:653]{损失函数}）（Posterior expected loss (see also Loss function)），\hyperref[eq:7.3.19]{\pageref*{eq:7.3.19}}\par
\par\addvspace{2pt}
\par\noindent\hangindent=1em\hspace*{0em}\phantomsection\label{idx:821}幂平均数（Power means），\hyperref[sec:4.9.2]{\pageref*{sec:4.9.2}}, \hyperref[exercise:7.16]{\pageref*{exercise:7.16}}\par
\par\addvspace{2pt}
\par\noindent\hangindent=1em\hspace*{0em}\phantomsection\label{idx:822}Pratt定理（Pratt’s Theorem），\hyperref[thm:9.3.9]{\pageref*{thm:9.3.9}}\par
\par\addvspace{2pt}
\par\noindent\hangindent=1em\hspace*{0em}\phantomsection\label{idx:823}预测区间（Prediction interval），\hyperref[def:11.3.5]{\pageref*{def:11.3.5}}\par
\par\addvspace{2pt}
\par\noindent\hangindent=1em\hspace*{0em}\phantomsection\label{idx:824}先验分布（Prior distribution），\hyperref[idxocc:824:293]{\pageref*{idxocc:824:293}}\par
\needspace{3\baselineskip}
\par\noindent\hangindent=2em\hspace*{1em}\phantomsection\label{idx:825}共轭（conjugate）\par
\par\noindent\hangindent=2.8em\hspace*{1.8em}\phantomsection\label{idx:826}伽马（gamma），\hyperref[exercise:7.23]{\pageref*{exercise:7.23}}\par
\par\noindent\hangindent=2.8em\hspace*{1.8em}\phantomsection\label{idx:827}泊松（Poisson），\hyperref[exercise:7.24]{\pageref*{exercise:7.24}}\par
\par\addvspace{2pt}
\par\noindent\hangindent=1em\hspace*{0em}\phantomsection\label{idx:828}概率密度函数（Probability density function (pdf)），\hyperref[def:1.6.3]{\pageref*{def:1.6.3}}\par
\par\noindent\hangindent=2em\hspace*{1em}\phantomsection\label{idx:829}二元变换（bivariate transformation），\hyperref[eq:4.3.2]{\pageref*{eq:4.3.2}}\par
\par\noindent\hangindent=2em\hspace*{1em}\phantomsection\label{idx:830}条件（conditional），\hyperref[def:4.2.3]{\pageref*{def:4.2.3}}, \hyperref[eq:4.6.6]{\pageref*{eq:4.6.6}}\par
\par\noindent\hangindent=2em\hspace*{1em}\phantomsection\label{idx:831}构造（construction），\hyperref[idxocc:831:61]{\pageref*{idxocc:831:61}}\par
\par\noindent\hangindent=2em\hspace*{1em}\phantomsection\label{idx:832}偶（even），\hyperref[exercise:2.25]{\pageref*{exercise:2.25}}\par
\par\noindent\hangindent=2em\hspace*{1em}\phantomsection\label{idx:833}联合（joint），\hyperref[def:4.1.10]{\pageref*{def:4.1.10}}, \hyperref[eq:4.6.2]{\pageref*{eq:4.6.2}}\par
\par\noindent\hangindent=2em\hspace*{1em}\phantomsection\label{idx:834}边缘（marginal），\hyperref[eq:4.1.2]{\pageref*{eq:4.1.2}}, \hyperref[idxocc:834:176]{\pageref*{idxocc:834:176}}\par
\par\noindent\hangindent=2em\hspace*{1em}\phantomsection\label{idx:835}单调变换（monotone transformation），\hyperref[thm:2.1.5]{\pageref*{thm:2.1.5}}\par
\par\noindent\hangindent=2em\hspace*{1em}\phantomsection\label{idx:836}多元变换（multivariate transformation），\hyperref[eq:4.6.7]{\pageref*{eq:4.6.7}}\par
\par\noindent\hangindent=2em\hspace*{1em}\phantomsection\label{idx:837}分段单调变换（piecewise monotone transformation），\hyperref[thm:2.1.8]{\pageref*{thm:2.1.8}}\par
\par\noindent\hangindent=2em\hspace*{1em}\phantomsection\label{idx:838}性质（properties），\hyperref[thm:1.6.5]{\pageref*{thm:1.6.5}}\par
\par\noindent\hangindent=2em\hspace*{1em}\phantomsection\label{idx:839}样本（sample），\hyperref[eq:5.1.1]{\pageref*{eq:5.1.1}}\par
\par\noindent\hangindent=2em\hspace*{1em}\phantomsection\label{idx:840}对称（symmetric），\hyperref[exercise:2.26]{\pageref*{exercise:2.26}}\par
\par\noindent\hangindent=2em\hspace*{1em}\phantomsection\label{idx:841}截断（truncated），\hyperref[exercise:1.52]{\pageref*{exercise:1.52}}\par
\par\noindent\hangindent=2em\hspace*{1em}\phantomsection\label{idx:842}单峰（unimodal），\hyperref[exercise:2.27]{\pageref*{exercise:2.27}}\par
\par\addvspace{2pt}
\par\noindent\hangindent=1em\hspace*{0em}\phantomsection\label{idx:843}概率函数（Probability function），\hyperref[def:1.2.4]{\pageref*{def:1.2.4}}\par
\par\noindent\hangindent=2em\hspace*{1em}\phantomsection\label{idx:844}有限集（finite set），\hyperref[thm:1.2.6]{\pageref*{thm:1.2.6}}\par
\par\noindent\hangindent=2em\hspace*{1em}\phantomsection\label{idx:845}诱导（induced），\hyperref[eq:1.4.1]{\pageref*{eq:1.4.1}}\par
\par\noindent\hangindent=2em\hspace*{1em}\phantomsection\label{idx:846}性质（properties），\hyperref[thm:1.2.8]{\pageref*{thm:1.2.8}}\par
\par\addvspace{2pt}
\par\noindent\hangindent=1em\hspace*{0em}\phantomsection\label{idx:847}概率积分变换（Probability integral transformation），\hyperref[thm:2.1.10]{\pageref*{thm:2.1.10}}, \hyperref[ex:5.6.3]{\pageref*{ex:5.6.3}}\par
\par\noindent\hangindent=2em\hspace*{1em}\phantomsection\label{idx:848}置信区间（confidence interval），\hyperref[thm:9.2.12]{\pageref*{thm:9.2.12}}\par
\par\noindent\hangindent=2em\hspace*{1em}\phantomsection\label{idx:849}离散（discrete），\hyperref[exercise:2.10]{\pageref*{exercise:2.10}}, \hyperref[thm:9.2.14]{\pageref*{thm:9.2.14}}\par
\par\addvspace{2pt}
\par\noindent\hangindent=1em\hspace*{0em}\phantomsection\label{idx:850}概率质量函数（Probability mass function (pmf)），\hyperref[def:1.6.1]{\pageref*{def:1.6.1}}\par
\par\noindent\hangindent=2em\hspace*{1em}\phantomsection\label{idx:851}条件（conditional），\hyperref[def:4.2.1]{\pageref*{def:4.2.1}}, \hyperref[eq:4.6.6]{\pageref*{eq:4.6.6}}\par
\par\noindent\hangindent=2em\hspace*{1em}\phantomsection\label{idx:852}构造（construction），\hyperref[idxocc:852:61]{\pageref*{idxocc:852:61}}\par
\par\noindent\hangindent=2em\hspace*{1em}\phantomsection\label{idx:853}联合（joint），\hyperref[eq:4.6.1]{\pageref*{eq:4.6.1}}\par
\par\noindent\hangindent=2.8em\hspace*{1.8em}\phantomsection\label{idx:854}二元（bivariate），\hyperref[def:4.1.3]{\pageref*{def:4.1.3}}\par
\par\noindent\hangindent=2em\hspace*{1em}\phantomsection\label{idx:855}边缘（marginal），\hyperref[thm:4.1.6]{\pageref*{thm:4.1.6}}, \hyperref[idxocc:855:176]{\pageref*{idxocc:855:176}}\par
\par\noindent\hangindent=2em\hspace*{1em}\phantomsection\label{idx:856}性质（properties），\hyperref[thm:1.6.5]{\pageref*{thm:1.6.5}}\par
\par\noindent\hangindent=2em\hspace*{1em}\phantomsection\label{idx:857}样本（sample），\hyperref[eq:5.1.1]{\pageref*{eq:5.1.1}}\par
\par\addvspace{2pt}
\par\noindent\hangindent=1em\hspace*{0em}\phantomsection\label{idx:858}错误覆盖概率（Probability of false coverage），\hyperref[eq:9.3.2]{\pageref*{eq:9.3.2}}\par
\par\addvspace{2pt}
\par\noindent\hangindent=1em\hspace*{0em}\phantomsection\label{idx:859}尼罗河问题（Problem of the Nile），\hyperref[exercise:7.54]{\pageref*{exercise:7.54}}\par
\par\addvspace{2pt}
\par\noindent\hangindent=1em\hspace*{0em}\phantomsection\label{idx:860}四分位数（Quartile），\hyperref[idxocc:860:216]{\pageref*{idxocc:860:216}}\par
\par\addvspace{2pt}
\par\noindent\hangindent=1em\hspace*{0em}\phantomsection\label{idx:861}随机数生成（Random number generation），\hyperref[idxocc:861:76]{\pageref*{idxocc:861:76}}, \hyperref[exercise:5.50]{\pageref*{exercise:5.50}}\par
\par\addvspace{2pt}
\needspace{3\baselineskip}
\par\noindent\hangindent=1em\hspace*{0em}\phantomsection\label{idx:862}随机样本（Random sample）\par
\par\noindent\hangindent=2em\hspace*{1em}\phantomsection\label{idx:863}生成（generating），\hyperref[sec:5.6]{\pageref*{sec:5.6}}\par
\par\addvspace{2pt}
\par\noindent\hangindent=1em\hspace*{0em}\phantomsection\label{idx:864}随机变量（Random variable），\hyperref[def:1.4.1]{\pageref*{def:1.4.1}}\par
\par\noindent\hangindent=2em\hspace*{1em}\phantomsection\label{idx:865}绝对连续（absolutely continuous），\hyperref[idxocc:865:62]{\pageref*{idxocc:865:62}}\par
\par\noindent\hangindent=2em\hspace*{1em}\phantomsection\label{idx:866}删失（censored），\hyperref[exercise:4.26]{\pageref*{exercise:4.26}}\par
\par\noindent\hangindent=2em\hspace*{1em}\phantomsection\label{idx:867}连续（continuous），\hyperref[def:1.5.7]{\pageref*{def:1.5.7}}\par
\par\noindent\hangindent=2em\hspace*{1em}\phantomsection\label{idx:868}离散（discrete），\hyperref[def:1.5.7]{\pageref*{def:1.5.7}}\par
\par\noindent\hangindent=2em\hspace*{1em}\phantomsection\label{idx:869}同分布（identically distributed），\hyperref[def:1.5.8]{\pageref*{def:1.5.8}}\par
\par\noindent\hangindent=2.8em\hspace*{1.8em}\phantomsection\label{idx:870}矩（moments），\hyperref[thm:2.3.11]{\pageref*{thm:2.3.11}}\par
\par\addvspace{2pt}
\par\noindent\hangindent=1em\hspace*{0em}\phantomsection\label{idx:871}随机向量（Random vector），\hyperref[def:4.1.1]{\pageref*{def:4.1.1}}\par
\par\noindent\hangindent=2em\hspace*{1em}\phantomsection\label{idx:872}连续（continuous），\hyperref[def:4.1.10]{\pageref*{def:4.1.10}}\par
\par\noindent\hangindent=2em\hspace*{1em}\phantomsection\label{idx:873}离散（discrete），\hyperref[idxocc:873:146]{\pageref*{idxocc:873:146}}\par
\par\addvspace{2pt}
\needspace{3\baselineskip}
\par\noindent\hangindent=1em\hspace*{0em}\phantomsection\label{idx:874}随机完全区组设计（Randomized complete block design (RCB)）\par
\par\noindent\hangindent=2em\hspace*{1em}\phantomsection\label{idx:875}与变量含误差的关系（relation to EIV），\hyperref[idxocc:875:506]{\pageref*{idxocc:875:506}}\par
\par\addvspace{2pt}
\par\noindent\hangindent=1em\hspace*{0em}\phantomsection\label{idx:876}随机停止的和（Randomly stopped sum），\hyperref[exercise:4.33]{\pageref*{exercise:4.33}}\par
\par\addvspace{2pt}
\par\noindent\hangindent=1em\hspace*{0em}\phantomsection\label{idx:877}Rao–Blackwell定理（Rao–Blackwell Theorem），\hyperref[thm:7.3.17]{\pageref*{thm:7.3.17}}\par
\par\addvspace{2pt}
\par\noindent\hangindent=1em\hspace*{0em}\phantomsection\label{idx:878}回归（Regression），\hyperref[idxocc:878:452]{\pageref*{idxocc:878:452}}, \hyperref[sec:11.3]{\pageref*{sec:11.3}}\par
\par\noindent\hangindent=2em\hspace*{1em}\phantomsection\label{idx:879}方差分析表（ANOVA table），\hyperref[tab:11.3.2]{\pageref*{tab:11.3.2}}\par
\par\noindent\hangindent=2em\hspace*{1em}\phantomsection\label{idx:880}二元正态（bivariate normal），\hyperref[sec:11.3.3.2]{\pageref*{sec:11.3.3.2}}\par
\par\noindent\hangindent=2em\hspace*{1em}\phantomsection\label{idx:881}截距的最优线性无偏估计量（BLUE of intercept），\hyperref[eq:11.3.21]{\pageref*{eq:11.3.21}}\par
\par\noindent\hangindent=2em\hspace*{1em}\phantomsection\label{idx:882}斜率的最优线性无偏估计量（BLUE of slope），\hyperref[eq:11.3.19]{\pageref*{eq:11.3.19}}\par
\par\noindent\hangindent=2em\hspace*{1em}\phantomsection\label{idx:883}决定系数（coefficient of determination），\hyperref[idxocc:883:481]{\pageref*{idxocc:883:481}}\par
\par\noindent\hangindent=2em\hspace*{1em}\phantomsection\label{idx:884}条件期望（conditional expectation），\hyperref[exercise:4.13]{\pageref*{exercise:4.13}}, \hyperref[eq:11.3.3]{\pageref*{eq:11.3.3}}\par
\par\noindent\hangindent=2em\hspace*{1em}\phantomsection\label{idx:885}条件正态（conditional normal），\hyperref[sec:11.3.3.1]{\pageref*{sec:11.3.3.1}}\par
\par\noindent\hangindent=2em\hspace*{1em}\phantomsection\label{idx:886}置信带（confidence band），\hyperref[idxocc:886:485]{\pageref*{idxocc:886:485}}\par
\par\noindent\hangindent=2em\hspace*{1em}\phantomsection\label{idx:887}置信区间（confidence interval），\hyperref[eq:11.3.37]{\pageref*{eq:11.3.37}}\par
\par\noindent\hangindent=2em\hspace*{1em}\phantomsection\label{idx:888}数据拟合（data fitting），\hyperref[idxocc:888:468]{\pageref*{idxocc:888:468}}, \hyperref[idxocc:888:471]{\pageref*{idxocc:888:471}}\par
\par\noindent\hangindent=2em\hspace*{1em}\phantomsection\label{idx:889}因变量（dependent variable），\hyperref[idxocc:889:466]{\pageref*{idxocc:889:466}}\par
\sourcepageanchor{565}
\par\noindent\hangindent=2em\hspace*{1em}\phantomsection\label{idx:890}设计（design），\hyperref[idxocc:890:473]{\pageref*{idxocc:890:473}}\par
\par\noindent\hangindent=2em\hspace*{1em}\phantomsection\label{idx:891}外推（extrapolation），\hyperref[idxocc:891:486]{\pageref*{idxocc:891:486}}\par
\par\noindent\hangindent=2em\hspace*{1em}\phantomsection\label{idx:892}自变量（independent variable），\hyperref[idxocc:892:466]{\pageref*{idxocc:892:466}}\par
\par\noindent\hangindent=2em\hspace*{1em}\phantomsection\label{idx:893}推断（inference），\hyperref[idxocc:893:468]{\pageref*{idxocc:893:468}}\par
\par\noindent\hangindent=2em\hspace*{1em}\phantomsection\label{idx:894}最小绝对偏差（least absolute deviation），\hyperref[idxocc:894:514]{\pageref*{idxocc:894:514}}, \hyperref[exercise:12.25]{\pageref*{exercise:12.25}}\par
\par\noindent\hangindent=2em\hspace*{1em}\phantomsection\label{idx:895}线性（linear），\hyperref[idxocc:895:467]{\pageref*{idxocc:895:467}}, \hyperref[sec:11.3.3]{\pageref*{sec:11.3.3}}\par
\par\noindent\hangindent=2em\hspace*{1em}\phantomsection\label{idx:896}关于x线性（linear in x），\hyperref[idxocc:896:467]{\pageref*{idxocc:896:467}}\par
\par\noindent\hangindent=2em\hspace*{1em}\phantomsection\label{idx:897}最大概率估计量（maximum probability estimator），\hyperref[item:11.37.e]{\pageref*{item:11.37.e}}\par
\par\noindent\hangindent=2em\hspace*{1em}\phantomsection\label{idx:898}最大似然估计量（MLE），\hyperref[sec:11.3.4]{\pageref*{sec:11.3.4}}\par
\par\noindent\hangindent=2.8em\hspace*{1.8em}\phantomsection\label{idx:899}有偏（biased），\hyperref[idxocc:899:477]{\pageref*{idxocc:899:477}}\par
\par\noindent\hangindent=2.8em\hspace*{1.8em}\phantomsection\label{idx:900}抽样分布（sampling distribution），\hyperref[thm:11.3.3]{\pageref*{thm:11.3.3}}\par
\par\noindent\hangindent=2em\hspace*{1em}\phantomsection\label{idx:901}总体（population），\hyperref[idxocc:901:452]{\pageref*{idxocc:901:452}}, \hyperref[eq:11.3.2]{\pageref*{eq:11.3.2}}\par
\par\noindent\hangindent=2em\hspace*{1em}\phantomsection\label{idx:902}预测区间（prediction interval），\hyperref[eq:11.3.41]{\pageref*{eq:11.3.41}}\par
\par\noindent\hangindent=2em\hspace*{1em}\phantomsection\label{idx:903}预测变量（predictor variable），\hyperref[idxocc:903:467]{\pageref*{idxocc:903:467}}\par
\par\noindent\hangindent=2em\hspace*{1em}\phantomsection\label{idx:904}$r^2$，\hyperref[idxocc:904:481]{\pageref*{idxocc:904:481}}\par
\par\noindent\hangindent=2em\hspace*{1em}\phantomsection\label{idx:905}残差（residuals），\hyperref[eq:11.3.27]{\pageref*{eq:11.3.27}}\par
\par\noindent\hangindent=2em\hspace*{1em}\phantomsection\label{idx:906}响应变量（response variable），\hyperref[idxocc:906:467]{\pageref*{idxocc:906:467}}\par
\par\noindent\hangindent=2em\hspace*{1em}\phantomsection\label{idx:907}平方和（sum of squares），\hyperref[eq:11.3.6]{\pageref*{eq:11.3.6}}, \hyperref[eq:11.3.36]{\pageref*{eq:11.3.36}}\par
\par\noindent\hangindent=2em\hspace*{1em}\phantomsection\label{idx:908}t统计量（t statistic），\hyperref[eq:11.3.33]{\pageref*{eq:11.3.33}}\par
\par\noindent\hangindent=2em\hspace*{1em}\phantomsection\label{idx:909}t检验（t test），\hyperref[eq:11.3.35]{\pageref*{eq:11.3.35}}, \hyperref[eq:11.3.38]{\pageref*{eq:11.3.38}}\par
\par\noindent\hangindent=2em\hspace*{1em}\phantomsection\label{idx:910}向均值回归（toward the mean），\hyperref[idxocc:910:467]{\pageref*{idxocc:910:467}}\par
\par\noindent\hangindent=2em\hspace*{1em}\phantomsection\label{idx:911}零截距（zero intercept），\hyperref[exercise:7.19]{\pageref*{exercise:7.19}}\par
\par\addvspace{2pt}
\par\noindent\hangindent=1em\hspace*{0em}\phantomsection\label{idx:912}正则条件（Regularity conditions），\hyperref[sec:10.6.2]{\pageref*{sec:10.6.2}}\par
\par\addvspace{2pt}
\par\noindent\hangindent=1em\hspace*{0em}\phantomsection\label{idx:913}重抽样（Resampling），\hyperref[ex:10.1.19]{\pageref*{ex:10.1.19}}, \hyperref[eq:10.1.11]{\pageref*{eq:10.1.11}}\par
\par\addvspace{2pt}
\par\noindent\hangindent=1em\hspace*{0em}\phantomsection\label{idx:914}风险函数（Risk function），\hyperref[eq:7.3.16]{\pageref*{eq:7.3.16}}\par
\par\addvspace{2pt}
\needspace{3\baselineskip}
\par\noindent\hangindent=1em\hspace*{0em}\phantomsection\label{idx:915}稳健（Robust）\par
\par\noindent\hangindent=2em\hspace*{1em}\phantomsection\label{idx:916}Tukey模型（Tukey model），\hyperref[exercise:10.23]{\pageref*{exercise:10.23}}\par
\par\addvspace{2pt}
\par\noindent\hangindent=1em\hspace*{0em}\phantomsection\label{idx:917}样本（Sample），\hyperref[def:5.1.1]{\pageref*{def:5.1.1}}\par
\par\noindent\hangindent=2em\hspace*{1em}\phantomsection\label{idx:918}指数族（exponential family），\hyperref[thm:5.2.11]{\pageref*{thm:5.2.11}}\par
\par\noindent\hangindent=2em\hspace*{1em}\phantomsection\label{idx:919}有限总体（finite population），\hyperref[idxocc:919:201]{\pageref*{idxocc:919:201}}\par
\par\noindent\hangindent=2em\hspace*{1em}\phantomsection\label{idx:920}无限总体（infinite population），\hyperref[idxocc:920:201]{\pageref*{idxocc:920:201}}\par
\par\noindent\hangindent=2em\hspace*{1em}\phantomsection\label{idx:921}均值（mean），\hyperref[def:5.2.2]{\pageref*{def:5.2.2}}\par
\par\noindent\hangindent=2.8em\hspace*{1.8em}\phantomsection\label{idx:922}柯西（Cauchy），\hyperref[ex:5.2.10]{\pageref*{ex:5.2.10}}\par
\par\noindent\hangindent=2.8em\hspace*{1.8em}\phantomsection\label{idx:923}期望值（expected value），\hyperref[thm:5.2.6]{\pageref*{thm:5.2.6}}\par
\par\noindent\hangindent=2.8em\hspace*{1.8em}\phantomsection\label{idx:924}位置尺度（location–scale），\hyperref[idxocc:924:208]{\pageref*{idxocc:924:208}}\par
\par\noindent\hangindent=2.8em\hspace*{1.8em}\phantomsection\label{idx:925}矩母函数（mgf），\hyperref[thm:5.2.7]{\pageref*{thm:5.2.7}}\par
\par\noindent\hangindent=2.8em\hspace*{1.8em}\phantomsection\label{idx:926}正态（normal），\hyperref[ex:5.2.8]{\pageref*{ex:5.2.8}}, \hyperref[thm:5.3.1]{\pageref*{thm:5.3.1}}\par
\par\noindent\hangindent=2.8em\hspace*{1.8em}\phantomsection\label{idx:927}概率密度函数（pdf），\hyperref[exercise:5.5]{\pageref*{exercise:5.5}}\par
\par\noindent\hangindent=2.8em\hspace*{1.8em}\phantomsection\label{idx:928}递推关系（recursion relation），\hyperref[exercise:5.15]{\pageref*{exercise:5.15}}\par
\par\noindent\hangindent=2.8em\hspace*{1.8em}\phantomsection\label{idx:929}截尾（trimmed），\hyperref[idxocc:929:442]{\pageref*{idxocc:929:442}}\par
\par\noindent\hangindent=2em\hspace*{1em}\phantomsection\label{idx:930}中位数（median），\hyperref[eq:5.4.1]{\pageref*{eq:5.4.1}}\par
\par\noindent\hangindent=2em\hspace*{1em}\phantomsection\label{idx:931}中程数（midrange），\hyperref[ex:5.4.7]{\pageref*{ex:5.4.7}}\par
\par\noindent\hangindent=2em\hspace*{1em}\phantomsection\label{idx:932}概率密度函数与概率质量函数（pdf, pmf），\hyperref[eq:5.1.1]{\pageref*{eq:5.1.1}}\par
\par\noindent\hangindent=2em\hspace*{1em}\phantomsection\label{idx:933}极差（range），\hyperref[idxocc:933:215]{\pageref*{idxocc:933:215}}\par
\par\noindent\hangindent=2em\hspace*{1em}\phantomsection\label{idx:934}简单随机（simple random），\hyperref[idxocc:934:202]{\pageref*{idxocc:934:202}}\par
\par\noindent\hangindent=2em\hspace*{1em}\phantomsection\label{idx:935}标准差（standard deviation），\hyperref[def:5.2.3]{\pageref*{def:5.2.3}}\par
\par\noindent\hangindent=2em\hspace*{1em}\phantomsection\label{idx:936}方差（variance），\hyperref[def:5.2.3]{\pageref*{def:5.2.3}}\par
\par\noindent\hangindent=2.8em\hspace*{1.8em}\phantomsection\label{idx:937}偏差（bias），\hyperref[sec:5.8.2]{\pageref*{sec:5.8.2}}\par
\par\noindent\hangindent=2.8em\hspace*{1.8em}\phantomsection\label{idx:938}期望值（expected value），\hyperref[thm:5.2.6]{\pageref*{thm:5.2.6}}\par
\par\noindent\hangindent=2.8em\hspace*{1.8em}\phantomsection\label{idx:939}恒等式（identity），\hyperref[exercise:5.8]{\pageref*{exercise:5.8}}\par
\par\noindent\hangindent=2.8em\hspace*{1.8em}\phantomsection\label{idx:940}正态（normal），\hyperref[thm:5.3.1]{\pageref*{thm:5.3.1}}\par
\par\noindent\hangindent=2.8em\hspace*{1.8em}\phantomsection\label{idx:941}递推关系（recursion relation），\hyperref[exercise:5.15]{\pageref*{exercise:5.15}}\par
\par\noindent\hangindent=2em\hspace*{1em}\phantomsection\label{idx:942}有放回（with replacement），\hyperref[idxocc:942:202]{\pageref*{idxocc:942:202}}\par
\par\noindent\hangindent=2em\hspace*{1em}\phantomsection\label{idx:943}无放回（without replacement），\hyperref[idxocc:943:202]{\pageref*{idxocc:943:202}}\par
\par\addvspace{2pt}
\needspace{3\baselineskip}
\par\noindent\hangindent=1em\hspace*{0em}\phantomsection\label{idx:944}样本极差（Sample range）\par
\par\noindent\hangindent=2em\hspace*{1em}\phantomsection\label{idx:945}均匀（uniform），\hyperref[ex:5.4.7]{\pageref*{ex:5.4.7}}\par
\par\addvspace{2pt}
\needspace{3\baselineskip}
\par\noindent\hangindent=1em\hspace*{0em}\phantomsection\label{idx:946}样本量（Sample size）\par
\par\noindent\hangindent=2em\hspace*{1em}\phantomsection\label{idx:947}假设检验：正态（hypothesis test – normal），\hyperref[ex:8.3.4]{\pageref*{ex:8.3.4}}\par
\par\addvspace{2pt}
\par\noindent\hangindent=1em\hspace*{0em}\phantomsection\label{idx:948}样本空间（Sample space），\hyperref[def:1.1.1]{\pageref*{def:1.1.1}}\par
\par\addvspace{2pt}
\needspace{3\baselineskip}
\par\noindent\hangindent=1em\hspace*{0em}\phantomsection\label{idx:949}抽样（Sampling）\par
\par\noindent\hangindent=2em\hspace*{1em}\phantomsection\label{idx:950}有放回（with replacement），\hyperref[ex:10.1.19]{\pageref*{ex:10.1.19}}\par
\par\addvspace{2pt}
\par\noindent\hangindent=1em\hspace*{0em}\phantomsection\label{idx:951}抽样分布（Sampling distribution），\hyperref[def:5.2.1]{\pageref*{def:5.2.1}}\par
\par\addvspace{2pt}
\par\noindent\hangindent=1em\hspace*{0em}\phantomsection\label{idx:952}Satterthwaite近似（Satterthwaite approximation），\hyperref[ex:7.2.3]{\pageref*{ex:7.2.3}}, \hyperref[exercise:8.42]{\pageref*{exercise:8.42}}\par
\par\addvspace{2pt}
\needspace{3\baselineskip}
\par\noindent\hangindent=1em\hspace*{0em}\phantomsection\label{idx:953}尺度族（Scale family）\par
\par\noindent\hangindent=2em\hspace*{1em}\phantomsection\label{idx:954}枢轴量（pivotal quantity），\hyperref[ex:9.2.7]{\pageref*{ex:9.2.7}}\par
\par\addvspace{2pt}
\par\noindent\hangindent=1em\hspace*{0em}\phantomsection\label{idx:955}Scheffé方法（Scheffé’s method），\hyperref[thm:11.2.10]{\pageref*{thm:11.2.10}}\par
\par\noindent\hangindent=2em\hspace*{1em}\phantomsection\label{idx:956}方差分析（ANOVA），\hyperref[exercise:11.14]{\pageref*{exercise:11.14}}\par
\par\addvspace{2pt}
\needspace{3\baselineskip}
\par\noindent\hangindent=1em\hspace*{0em}\phantomsection\label{idx:957}集合估计量（Set estimator）\par
\par\noindent\hangindent=2em\hspace*{1em}\phantomsection\label{idx:958}与置信区间的关系（relation to confidence interval），\hyperref[idxocc:958:369]{\pageref*{idxocc:958:369}}\par
\par\addvspace{2pt}
\needspace{3\baselineskip}
\par\noindent\hangindent=1em\hspace*{0em}\phantomsection\label{idx:959}集合运算（Set operations）\par
\par\noindent\hangindent=2em\hspace*{1em}\phantomsection\label{idx:960}结合律（associativity），\hyperref[thm:1.1.4]{\pageref*{thm:1.1.4}}\par
\par\noindent\hangindent=2em\hspace*{1em}\phantomsection\label{idx:961}交换律（commutativity），\hyperref[thm:1.1.4]{\pageref*{thm:1.1.4}}\par
\par\noindent\hangindent=2em\hspace*{1em}\phantomsection\label{idx:962}求补（complementation），\hyperref[idxocc:962:33]{\pageref*{idxocc:962:33}}\par
\par\noindent\hangindent=2em\hspace*{1em}\phantomsection\label{idx:963}德摩根律（DeMorgan’s Laws），\hyperref[thm:1.1.4]{\pageref*{thm:1.1.4}}\par
\par\noindent\hangindent=2em\hspace*{1em}\phantomsection\label{idx:964}分配律（Distributive Law），\hyperref[thm:1.1.4]{\pageref*{thm:1.1.4}}\par
\par\noindent\hangindent=2em\hspace*{1em}\phantomsection\label{idx:965}恒等式（identities），\hyperref[exercise:1.2]{\pageref*{exercise:1.2}}\par
\par\noindent\hangindent=2em\hspace*{1em}\phantomsection\label{idx:966}交（intersection），\hyperref[idxocc:966:33]{\pageref*{idxocc:966:33}}\par
\par\noindent\hangindent=2.8em\hspace*{1.8em}\phantomsection\label{idx:967}可数（countable），\hyperref[idxocc:967:34]{\pageref*{idxocc:967:34}}\par
\par\noindent\hangindent=2.8em\hspace*{1.8em}\phantomsection\label{idx:968}不可数（uncountable），\hyperref[idxocc:968:35]{\pageref*{idxocc:968:35}}\par
\par\noindent\hangindent=2em\hspace*{1em}\phantomsection\label{idx:969}并（union），\hyperref[idxocc:969:33]{\pageref*{idxocc:969:33}}\par
\par\noindent\hangindent=2.8em\hspace*{1.8em}\phantomsection\label{idx:970}可数（countable），\hyperref[idxocc:970:34]{\pageref*{idxocc:970:34}}\par
\par\noindent\hangindent=2.8em\hspace*{1.8em}\phantomsection\label{idx:971}不可数（uncountable），\hyperref[idxocc:971:35]{\pageref*{idxocc:971:35}}\par
\par\addvspace{2pt}
\par\noindent\hangindent=1em\hspace*{0em}\phantomsection\label{idx:972}\ensuremath{\sigma}代数（Sigma algebra），\hyperref[def:1.2.1]{\pageref*{def:1.2.1}}\par
\par\noindent\hangindent=2em\hspace*{1em}\phantomsection\label{idx:973}可数（countable），\hyperref[ex:1.2.2]{\pageref*{ex:1.2.2}}\par
\par\noindent\hangindent=2em\hspace*{1em}\phantomsection\label{idx:974}生成的（generated），\hyperref[idxocc:974:59]{\pageref*{idxocc:974:59}}\par
\par\noindent\hangindent=2em\hspace*{1em}\phantomsection\label{idx:975}性质（properties），\hyperref[exercise:1.11]{\pageref*{exercise:1.11}}\par
\par\noindent\hangindent=2em\hspace*{1em}\phantomsection\label{idx:976}平凡（trivial），\hyperref[idxocc:976:36]{\pageref*{idxocc:976:36}}\par
\par\noindent\hangindent=2em\hspace*{1em}\phantomsection\label{idx:977}不可数（uncountable），\hyperref[ex:1.2.3]{\pageref*{ex:1.2.3}}\par
\par\addvspace{2pt}
\par\noindent\hangindent=1em\hspace*{0em}\phantomsection\label{idx:978}偏度（Skewness），\hyperref[exercise:2.28]{\pageref*{exercise:2.28}}\par
\par\addvspace{2pt}
\par\noindent\hangindent=1em\hspace*{0em}\phantomsection\label{idx:979}Slutsky定理（Slutsky’s Theorem），\hyperref[thm:5.5.17]{\pageref*{thm:5.5.17}}, \hyperref[idxocc:979:429]{\pageref*{idxocc:979:429}}\par
\par\addvspace{2pt}
\par\noindent\hangindent=1em\hspace*{0em}\phantomsection\label{idx:980}标准差（Standard deviation），\hyperref[def:2.3.2]{\pageref*{def:2.3.2}}\par
\par\addvspace{2pt}
\par\noindent\hangindent=1em\hspace*{0em}\phantomsection\label{idx:981}统计量（Statistic），\hyperref[def:5.2.1]{\pageref*{def:5.2.1}}\par
\par\addvspace{2pt}
\needspace{3\baselineskip}
\par\noindent\hangindent=1em\hspace*{0em}\phantomsection\label{idx:982}Stein估计（Stein estimation）\par
\par\noindent\hangindent=2em\hspace*{1em}\phantomsection\label{idx:983}方差分析（ANOVA），\hyperref[sec:11.5.6]{\pageref*{sec:11.5.6}}\par
\par\addvspace{2pt}
\par\noindent\hangindent=1em\hspace*{0em}\phantomsection\label{idx:984}Stein引理（Stein’s Lemma），\hyperref[lem:3.6.5]{\pageref*{lem:3.6.5}}\par
\par\addvspace{2pt}
\par\noindent\hangindent=1em\hspace*{0em}\phantomsection\label{idx:985}Stirling公式（Stirling’s formula），\hyperref[exercise:1.28]{\pageref*{exercise:1.28}}, \hyperref[exercise:5.35]{\pageref*{exercise:5.35}}\par
\par\addvspace{2pt}
\needspace{3\baselineskip}
\par\noindent\hangindent=1em\hspace*{0em}\phantomsection\label{idx:986}随机序（Stochastic order）\par
\par\noindent\hangindent=2em\hspace*{1em}\phantomsection\label{idx:987}递减（decreasing），\hyperref[exercise:3.43]{\pageref*{exercise:3.43}}\par
\par\noindent\hangindent=2em\hspace*{1em}\phantomsection\label{idx:988}较大（greater），\hyperref[exercise:1.49]{\pageref*{exercise:1.49}}, \hyperref[exercise:2.10]{\pageref*{exercise:2.10}}\par
\par\noindent\hangindent=2em\hspace*{1em}\phantomsection\label{idx:989}递增（increasing），\hyperref[exercise:3.41]{\pageref*{exercise:3.41}}, \hyperref[idxocc:989:380]{\pageref*{idxocc:989:380}}\par
\par\noindent\hangindent=2.8em\hspace*{1.8em}\phantomsection\label{idx:990}卡方（chi squared），\hyperref[exercise:5.19]{\pageref*{exercise:5.19}}\par
\par\noindent\hangindent=2.8em\hspace*{1.8em}\phantomsection\label{idx:991}F，\hyperref[exercise:5.19]{\pageref*{exercise:5.19}}\par
\par\noindent\hangindent=2em\hspace*{1em}\phantomsection\label{idx:992}位置尺度族（location–scale family），\hyperref[exercise:8.34]{\pageref*{exercise:8.34}}\par
\sourcepageanchor{566}
\par\noindent\hangindent=2em\hspace*{1em}\phantomsection\label{idx:993}单调似然比（MLR），\hyperref[exercise:8.26]{\pageref*{exercise:8.26}}\par
\par\noindent\hangindent=2em\hspace*{1em}\phantomsection\label{idx:994}功效（power），\hyperref[sec:8.5.1]{\pageref*{sec:8.5.1}}\par
\par\addvspace{2pt}
\par\noindent\hangindent=1em\hspace*{0em}\phantomsection\label{idx:995}强大数定律（Strong Law of Large Numbers），\hyperref[sec:5.8.4]{\pageref*{sec:5.8.4}}\par
\par\addvspace{2pt}
\par\noindent\hangindent=1em\hspace*{0em}\phantomsection\label{idx:996}结构推断（Structural inference），\hyperref[idxocc:996:376]{\pageref*{idxocc:996:376}}\par
\par\addvspace{2pt}
\par\noindent\hangindent=1em\hspace*{0em}\phantomsection\label{idx:997}充分性原则（Sufficiency Principle），\hyperref[idxocc:997:252]{\pageref*{idxocc:997:252}}\par
\par\noindent\hangindent=2em\hspace*{1em}\phantomsection\label{idx:998}形式化（formal），\hyperref[idxocc:998:269]{\pageref*{idxocc:998:269}}\par
\par\noindent\hangindent=2em\hspace*{1em}\phantomsection\label{idx:999}解释（interpretations），\hyperref[ex:6.3.7]{\pageref*{ex:6.3.7}}\par
\par\addvspace{2pt}
\par\noindent\hangindent=1em\hspace*{0em}\phantomsection\label{idx:1000}充分统计量（Sufficient statistic），\hyperref[def:6.2.1]{\pageref*{def:6.2.1}}\par
\par\noindent\hangindent=2em\hspace*{1em}\phantomsection\label{idx:1001}方差分析（ANOVA），\hyperref[item:11.6.a]{\pageref*{item:11.6.a}}\par
\par\noindent\hangindent=2em\hspace*{1em}\phantomsection\label{idx:1002}伯努利（Bernoulli），\hyperref[ex:6.2.3]{\pageref*{ex:6.2.3}}\par
\par\noindent\hangindent=2em\hspace*{1em}\phantomsection\label{idx:1003}离散均匀（discrete uniform），\hyperref[ex:6.2.8]{\pageref*{ex:6.2.8}}\par
\par\noindent\hangindent=2em\hspace*{1em}\phantomsection\label{idx:1004}指数族（exponential family），\hyperref[thm:6.2.10]{\pageref*{thm:6.2.10}}\par
\par\noindent\hangindent=2em\hspace*{1em}\phantomsection\label{idx:1005}因子分解定理（Factorization Theorem），\hyperref[thm:6.2.6]{\pageref*{thm:6.2.6}}\par
\par\noindent\hangindent=2em\hspace*{1em}\phantomsection\label{idx:1006}最小（minimal），\hyperref[def:6.2.11]{\pageref*{def:6.2.11}}, \hyperref[exercise:6.26]{\pageref*{exercise:6.26}}\par
\par\noindent\hangindent=2.8em\hspace*{1.8em}\phantomsection\label{idx:1007}刻画（characterization），\hyperref[thm:6.2.13]{\pageref*{thm:6.2.13}}\par
\par\noindent\hangindent=2.8em\hspace*{1.8em}\phantomsection\label{idx:1008}必要（necessary），\hyperref[thm:6.6.4]{\pageref*{thm:6.6.4}}\par
\par\noindent\hangindent=2.8em\hspace*{1.8em}\phantomsection\label{idx:1009}正态（normal），\hyperref[ex:6.2.14]{\pageref*{ex:6.2.14}}\par
\par\noindent\hangindent=2.8em\hspace*{1.8em}\phantomsection\label{idx:1010}均匀（uniform），\hyperref[ex:6.2.15]{\pageref*{ex:6.2.15}}, \hyperref[exercise:6.10]{\pageref*{exercise:6.10}}\par
\par\noindent\hangindent=2em\hspace*{1em}\phantomsection\label{idx:1011}正态（normal），\hyperref[ex:6.2.4]{\pageref*{ex:6.2.4}}, \hyperref[ex:6.2.7]{\pageref*{ex:6.2.7}}, \hyperref[ex:6.2.9]{\pageref*{ex:6.2.9}}, \hyperref[ex:6.2.12]{\pageref*{ex:6.2.12}}\par
\par\addvspace{2pt}
\needspace{3\baselineskip}
\par\noindent\hangindent=1em\hspace*{0em}\phantomsection\label{idx:1012}平方和（Sum of squares）\par
\par\noindent\hangindent=2em\hspace*{1em}\phantomsection\label{idx:1013}残差（residual），\hyperref[idxocc:1013:469]{\pageref*{idxocc:1013:469}}\par
\par\addvspace{2pt}
\par\noindent\hangindent=1em\hspace*{0em}\phantomsection\label{idx:1014}超效率（Superefficiency），\hyperref[sec:10.6.1]{\pageref*{sec:10.6.1}}\par
\par\addvspace{2pt}
\par\noindent\hangindent=1em\hspace*{0em}\phantomsection\label{idx:1015}支持集（Support），\hyperref[idxocc:1015:71]{\pageref*{idxocc:1015:71}}\par
\par\addvspace{2pt}
\needspace{3\baselineskip}
\par\noindent\hangindent=1em\hspace*{0em}\phantomsection\label{idx:1016}t检验（t test）\par
\par\noindent\hangindent=2em\hspace*{1em}\phantomsection\label{idx:1017}方差分析（ANOVA），\hyperref[eq:11.2.7]{\pageref*{eq:11.2.7}}\par
\par\noindent\hangindent=2em\hspace*{1em}\phantomsection\label{idx:1018}近似（approximate），\hyperref[exercise:8.42]{\pageref*{exercise:8.42}}\par
\par\noindent\hangindent=2em\hspace*{1em}\phantomsection\label{idx:1019}似然比检验（LRT），\hyperref[ex:8.2.6]{\pageref*{ex:8.2.6}}\par
\par\noindent\hangindent=2.8em\hspace*{1.8em}\phantomsection\label{idx:1020}配对（paired），\hyperref[exercise:8.40]{\pageref*{exercise:8.40}}\par
\par\noindent\hangindent=2.8em\hspace*{1.8em}\phantomsection\label{idx:1021}双侧（two-sided），\hyperref[ex:8.2.8]{\pageref*{ex:8.2.8}}\par
\par\noindent\hangindent=2em\hspace*{1em}\phantomsection\label{idx:1022}配对（paired），\hyperref[exercise:8.39]{\pageref*{exercise:8.39}}\par
\par\noindent\hangindent=2em\hspace*{1em}\phantomsection\label{idx:1023}回归（regression），\hyperref[eq:11.3.38]{\pageref*{eq:11.3.38}}\par
\par\noindent\hangindent=2em\hspace*{1em}\phantomsection\label{idx:1024}双样本（two-sample），\hyperref[exercise:8.41]{\pageref*{exercise:8.41}}\par
\par\noindent\hangindent=2em\hspace*{1em}\phantomsection\label{idx:1025}双侧（two-sided），\hyperref[exercise:8.41]{\pageref*{exercise:8.41}}\par
\needspace{3\baselineskip}
\par\noindent\hangindent=2em\hspace*{1em}\phantomsection\label{idx:1026}一致最强（UMP）\par
\par\noindent\hangindent=2.8em\hspace*{1.8em}\phantomsection\label{idx:1027}无偏（unbiased），\hyperref[exercise:8.37]{\pageref*{exercise:8.37}}\par
\par\noindent\hangindent=2em\hspace*{1em}\phantomsection\label{idx:1028}并交（union–intersection），\hyperref[ex:8.3.8]{\pageref*{ex:8.3.8}}\par
\par\addvspace{2pt}
\par\noindent\hangindent=1em\hspace*{0em}\phantomsection\label{idx:1029}泰勒级数（Taylor series），\hyperref[def:5.5.20]{\pageref*{def:5.5.20}}, \hyperref[idxocc:1029:423]{\pageref*{idxocc:1029:423}}\par
\par\addvspace{2pt}
\par\noindent\hangindent=1em\hspace*{0em}\phantomsection\label{idx:1030}泰勒定理（Taylor’s Theorem），\hyperref[thm:5.5.21]{\pageref*{thm:5.5.21}}\par
\par\noindent\hangindent=2em\hspace*{1em}\phantomsection\label{idx:1031}在中心极限定理中的用途（use in CLT），\hyperref[idxocc:1031:224]{\pageref*{idxocc:1031:224}}\par
\par\addvspace{2pt}
\par\noindent\hangindent=1em\hspace*{0em}\phantomsection\label{idx:1032}检验函数（Test function），\hyperref[idxocc:1032:346]{\pageref*{idxocc:1032:346}}\par
\par\addvspace{2pt}
\par\noindent\hangindent=1em\hspace*{0em}\phantomsection\label{idx:1033}检验统计量（Test statistic），\hyperref[idxocc:1033:334]{\pageref*{idxocc:1033:334}}\par
\par\addvspace{2pt}
\par\noindent\hangindent=1em\hspace*{0em}\phantomsection\label{idx:1034}阈值参数（Threshold parameter），\hyperref[ex:3.5.3]{\pageref*{ex:3.5.3}}\par
\par\addvspace{2pt}
\par\noindent\hangindent=1em\hspace*{0em}\phantomsection\label{idx:1035}全最小二乘（Total least squares），\hyperref[idxocc:1035:501]{\pageref*{idxocc:1035:501}}\par
\par\addvspace{2pt}
\needspace{3\baselineskip}
\par\noindent\hangindent=1em\hspace*{0em}\phantomsection\label{idx:1036}变换（Transformation）\par
\par\noindent\hangindent=2em\hspace*{1em}\phantomsection\label{idx:1037}Box–Cox（Box-Cox），\hyperref[exercise:7.16]{\pageref*{exercise:7.16}}\par
\par\noindent\hangindent=2em\hspace*{1em}\phantomsection\label{idx:1038}Box–Muller（Box-Muller），\hyperref[ex:5.6.4]{\pageref*{ex:5.6.4}}\par
\par\noindent\hangindent=2em\hspace*{1em}\phantomsection\label{idx:1039}Fisher的z变换（Fisher’s z），\hyperref[item:10.17.e]{\pageref*{item:10.17.e}}\par
\par\noindent\hangindent=2em\hspace*{1em}\phantomsection\label{idx:1040}单调（monotone），\hyperref[idxocc:1040:71]{\pageref*{idxocc:1040:71}}\par
\par\noindent\hangindent=2em\hspace*{1em}\phantomsection\label{idx:1041}概率积分（probability integral），\hyperref[ex:5.6.3]{\pageref*{ex:5.6.3}}\par
\par\addvspace{2pt}
\par\noindent\hangindent=1em\hspace*{0em}\phantomsection\label{idx:1042}截尾均值（Trimmed mean），\hyperref[idxocc:1042:442]{\pageref*{idxocc:1042:442}}\par
\par\addvspace{2pt}
\par\noindent\hangindent=1em\hspace*{0em}\phantomsection\label{idx:1043}调节参数（Tuning parameter），\hyperref[eq:10.2.2]{\pageref*{eq:10.2.2}}\par
\par\addvspace{2pt}
\par\noindent\hangindent=1em\hspace*{0em}\phantomsection\label{idx:1044}无偏估计量（Unbiased estimator），\hyperref[idxocc:1044:206]{\pageref*{idxocc:1044:206}}, \hyperref[def:7.3.2]{\pageref*{def:7.3.2}}\par
\par\noindent\hangindent=2em\hspace*{1em}\phantomsection\label{idx:1045}最优（best），\hyperref[def:7.3.7]{\pageref*{def:7.3.7}}\par
\par\noindent\hangindent=2.8em\hspace*{1.8em}\phantomsection\label{idx:1046}二项（binomial），\hyperref[ex:7.3.24]{\pageref*{ex:7.3.24}}\par
\par\noindent\hangindent=2.8em\hspace*{1.8em}\phantomsection\label{idx:1047}刻画（characterization），\hyperref[thm:7.3.20]{\pageref*{thm:7.3.20}}\par
\par\noindent\hangindent=2.8em\hspace*{1.8em}\phantomsection\label{idx:1048}完备性（completeness），\hyperref[thm:7.3.23]{\pageref*{thm:7.3.23}}\par
\par\noindent\hangindent=2.8em\hspace*{1.8em}\phantomsection\label{idx:1049}泊松（Poisson），\hyperref[ex:7.3.12]{\pageref*{ex:7.3.12}}\par
\par\noindent\hangindent=2.8em\hspace*{1.8em}\phantomsection\label{idx:1050}充分性（sufficiency），\hyperref[thm:7.3.17]{\pageref*{thm:7.3.17}}\par
\par\noindent\hangindent=2.8em\hspace*{1.8em}\phantomsection\label{idx:1051}均匀（uniform），\hyperref[ex:7.3.22]{\pageref*{ex:7.3.22}}\par
\par\noindent\hangindent=2.8em\hspace*{1.8em}\phantomsection\label{idx:1052}唯一（unique），\hyperref[thm:7.3.19]{\pageref*{thm:7.3.19}}\par
\par\noindent\hangindent=2em\hspace*{1em}\phantomsection\label{idx:1053}线性（linear），\hyperref[exercise:7.41]{\pageref*{exercise:7.41}}\par
\par\noindent\hangindent=2em\hspace*{1em}\phantomsection\label{idx:1054}位置（location），\hyperref[exercise:6.41]{\pageref*{exercise:6.41}}\par
\par\noindent\hangindent=2em\hspace*{1em}\phantomsection\label{idx:1055}正态（normal），\hyperref[ex:7.3.3]{\pageref*{ex:7.3.3}}\par
\par\noindent\hangindent=2.8em\hspace*{1.8em}\phantomsection\label{idx:1056}标准差（standard deviation），\hyperref[exercise:7.50]{\pageref*{exercise:7.50}}\par
\par\noindent\hangindent=2em\hspace*{1em}\phantomsection\label{idx:1057}泊松（Poisson），\hyperref[ex:7.3.8]{\pageref*{ex:7.3.8}}\par
\par\noindent\hangindent=2em\hspace*{1em}\phantomsection\label{idx:1058}与贝叶斯估计量的关系（relation to Bayes estimator），\hyperref[sec:7.5.2]{\pageref*{sec:7.5.2}}\par
\par\noindent\hangindent=2em\hspace*{1em}\phantomsection\label{idx:1059}均匀（uniform），\hyperref[ex:7.3.13]{\pageref*{ex:7.3.13}}, \hyperref[ex:7.3.21]{\pageref*{ex:7.3.21}}\par
\par\noindent\hangindent=2em\hspace*{1em}\phantomsection\label{idx:1060}方差界（variance bound），\hyperref[ex:7.3.14]{\pageref*{ex:7.3.14}}\par
\par\noindent\hangindent=2em\hspace*{1em}\phantomsection\label{idx:1061}零（zero），\hyperref[ex:7.3.21]{\pageref*{ex:7.3.21}}\par
\par\addvspace{2pt}
\par\noindent\hangindent=1em\hspace*{0em}\phantomsection\label{idx:1062}方差（Variance），\hyperref[def:2.3.2]{\pageref*{def:2.3.2}}\par
\par\noindent\hangindent=2em\hspace*{1em}\phantomsection\label{idx:1063}渐近（asymptotic），\hyperref[idxocc:1063:410]{\pageref*{idxocc:1063:410}}\par
\par\noindent\hangindent=2em\hspace*{1em}\phantomsection\label{idx:1064}条件（conditional），\hyperref[idxocc:1064:155]{\pageref*{idxocc:1064:155}}, \hyperref[thm:4.4.7]{\pageref*{thm:4.4.7}}\par
\par\noindent\hangindent=2em\hspace*{1em}\phantomsection\label{idx:1065}极限（limiting），\hyperref[def:10.1.7]{\pageref*{def:10.1.7}}\par
\par\noindent\hangindent=2em\hspace*{1em}\phantomsection\label{idx:1066}性质（properties），\hyperref[thm:2.3.4]{\pageref*{thm:2.3.4}}\par
\par\noindent\hangindent=2em\hspace*{1em}\phantomsection\label{idx:1067}均值的二次函数（quadratic function of mean），\hyperref[idxocc:1067:110]{\pageref*{idxocc:1067:110}}\par
\par\noindent\hangindent=2em\hspace*{1em}\phantomsection\label{idx:1068}和（sum），\hyperref[thm:4.5.6]{\pageref*{thm:4.5.6}}, \hyperref[exercise:4.44]{\pageref*{exercise:4.44}}, \hyperref[lem:5.2.5]{\pageref*{lem:5.2.5}}\par
\par\noindent\hangindent=2em\hspace*{1em}\phantomsection\label{idx:1069}泰勒近似（Taylor approximation），\hyperref[eq:5.5.9]{\pageref*{eq:5.5.9}}\par
\par\addvspace{2pt}
\par\noindent\hangindent=1em\hspace*{0em}\phantomsection\label{idx:1070}韦恩图（Venn diagram），\hyperref[idxocc:1070:34]{\pageref*{idxocc:1070:34}}\par
\par\addvspace{2pt}
\par\noindent\hangindent=1em\hspace*{0em}\phantomsection\label{idx:1071}看守问题（Warden problem），\hyperref[ex:1.3.4]{\pageref*{ex:1.3.4}}, \hyperref[exercise:1.37]{\pageref*{exercise:1.37}}\par
\end{multicols}
\endgroup
